% Localisation et nature de l’information génétique déterminant les caractères héréditaires — SVTEEHB 3ème — Cours signé OnBuch+
% Programme officiel MINESEC. Compiler avec : tectonic N02-localisation-nature-information-genetique-determinant-c.tex
\newcommand{\DOCMATIERE}{SVTEEHB}
\newcommand{\DOCNIVEAU}{3ème}
\newcommand{\DOCMODULE}{SVTEEHB – classe de 3ème}
\newcommand{\DOCLECON}{N02}
\newcommand{\DOCTITRE}{Localisation et nature de l’information génétique déterminant les caractères héréditaires}
% OnBuch+ — préambule « cours signés » ENRICHI (compilé avec tectonic / XeLaTeX)
% Système visuel « Pop » d'OnBuch+ : fond crème (jamais blanc), bordures encre
% épaisses, ombres « dures » décalées, titres Archivo Black en capitales.
% Version MATHÉMATIQUES : graphes (pgfplots/tikz), boîtes pédagogiques
% (définition, propriété, méthode, attention, le savais-tu, exercice résolu,
% exercice, TP, bilan), page de garde et signature OnBuch+ ancrée en bas.
%
% Macros attendues dans main.tex AVANT \input{preamble} :
%   \DOCMATIERE  \DOCNIVEAU  \DOCMODULE  \DOCLECON (numéro)  \DOCTITRE
\documentclass[11pt]{article}

\usepackage{fontspec}
\usepackage[french]{babel}
\usepackage[a4paper,margin=2cm,top=3.4cm,bottom=2.8cm,headheight=1.4cm,footskip=1.3cm]{geometry}
\usepackage{amsmath,amssymb,amsfonts}
\usepackage{mathtools} % \xmapsto, \coloneqq... souvent utilisés par les modèles
\usepackage{cancel} % simplifications barrées (bilans de forces)
% \abs{z} (module, valeur absolue) et \norm{u} (norme), utilisés par les modèles
\providecommand{\abs}{}\let\abs\relax\DeclarePairedDelimiter{\abs}{\lvert}{\rvert}
\providecommand{\norm}{}\let\norm\relax\DeclarePairedDelimiter{\norm}{\lVert}{\rVert}
\usepackage[table]{xcolor} % [table] : fournit aussi \rowcolors (alternance de lignes)
\usepackage{pagecolor}
\usepackage{tikz}
\usetikzlibrary{arrows.meta,positioning,calc,shapes.geometric,shapes.misc,shapes.arrows,decorations.pathmorphing,decorations.pathreplacing,decorations.markings,patterns,shadows,mindmap,trees,fit,backgrounds,matrix,intersections,angles,quotes}
\usepackage{pgfplots}
\usepackage[version=4]{mhchem} % équations nucléaires \ce{^{235}_{92}U}
\usepackage[european,siunitx]{circuitikz} % schémas électriques
\pgfplotsset{compat=1.18}
\usepgfplotslibrary{fillbetween,groupplots} % aires entre courbes (calcul intégral)
\usepackage{enumitem}
\usepackage{fancyhdr}
% [most] charge toutes les bibliothèques tcolorbox (listings, minted, raster,
% documentation...) et gonfle inutilement la mémoire TeX sur un document long
% (~300 boîtes) : on ne charge que ce qui est réellement utilisé.
\usepackage[skins,breakable]{tcolorbox}
\usepackage{graphicx}
\usepackage{colortbl}
\usepackage{booktabs}
\usepackage{tabularx}
\usepackage{multirow}
\usepackage{diagbox} % case d'en-tête diagonale des tableaux à double entrée
% Colonnes extensibles centrée / gauche / droite (C, L, R), souvent utilisées par les modèles
\newcolumntype{C}{>{\centering\arraybackslash}X}
\newcolumntype{L}{>{\raggedright\arraybackslash}X}
\newcolumntype{R}{>{\raggedleft\arraybackslash}X}
\usepackage{array}
\usepackage{multicol}
\usepackage{siunitx}
% parse-numbers=false : accepte aussi \SI{2,5 \times 10^{-3}}{...}
\sisetup{locale=FR,output-decimal-marker={,},group-minimum-digits=4,parse-numbers=false}
\DeclareSIUnit\ohm{\ensuremath{\symup{\Omega}}} % U+2126 absent de la police
\DeclareSIUnit\division{div} % réglages d'oscilloscope (ms/div, V/div)
\DeclareSIUnit\tr{tr} % tours (vitesse de rotation, ex. \SI{1500}{\tr\per\minute})
\DeclareSIUnit\molar{M} % concentration molaire (ex. \SI{3}{\molar}, \SI{1}{\micro\molar})
\DeclareSIUnit\an{an} % année (le modèle confond parfois avec \year, un compteur TeX) — ex. \SI{5}{\kilo\meter\per\an}
\DeclareSIUnit\FCFA{FCFA} % franc CFA (ex. \SI{500}{\FCFA\per\kilo\gram})
\DeclareSIUnit\degreeBrix{\textdegree Brix} % teneur en sucre d'un sirop/jus (réfractométrie)
\DeclareSIUnit\month{mois} % le modèle utilise parfois \month, un compteur TeX, comme unité de durée
\DeclareSIUnit\microliter{\micro\liter} % le modèle écrit parfois \microliter en un seul token
\DeclareSIUnit\million{millions} % dénombrement (ex. \SI{15}{\million\per\mL} spermatozoïdes)
\usepackage{qrcode}
% \degree : commande fréquemment utilisée par les modèles pour un angle ou une
% température en dehors de \SI{}{} (non fournie par les paquets chargés).
\providecommand{\degree}{^{\circ}}
% \qed (fin de démonstration), utilisé par les modèles sans amsthm
\providecommand{\qed}{\hfill$\square$}
% \barre : double barre des valeurs interdites dans un tableau de variations (tkz-tab)
\providecommand{\barre}{\ensuremath{\|}}
% environnement proof (amsthm non chargé) utilisé par les modèles
\ifdefined\proof\else\newenvironment{proof}[1][Démonstration]{\par\noindent\textit{#1.}\ }{\qed\par}\fi
\usepackage{lastpage}
\usepackage{hyperref}
\hypersetup{hidelinks}

% ============================================================
% Palette « Pop » OnBuch+ — mêmes hex que la classe P (plus_ui.dart)
% ============================================================
\definecolor{poporange}{HTML}{F59321}
\definecolor{popdark}{HTML}{E07A0C}
\definecolor{popcream}{HTML}{FFF6E8}
\definecolor{popink}{HTML}{1C1714}
\definecolor{poppurple}{HTML}{7A5AE0}
\definecolor{popgreen}{HTML}{2E9E6A}
\definecolor{poppink}{HTML}{E0457B}
\definecolor{popblue}{HTML}{2F7DE1}
\definecolor{popgold}{HTML}{C98A12}
\definecolor{popmuted}{HTML}{8A7F76}
\definecolor{popline}{HTML}{EADFD2}
\definecolor{popgray}{HTML}{8A7F76}
\colorlet{popgrey}{popgray}
% Alias tolérants (le modèle nomme parfois les couleurs différemment)
\colorlet{popwhite}{white}
\colorlet{popblack}{popink}
\colorlet{popred}{poppink}
\colorlet{popyellow}{popgold}
% Teintes claires (fonds de tableaux/encadrés) — le modèle nomme parfois
% les couleurs avec un suffixe L (ex. popblueL) sans les avoir définies.
\colorlet{poporangeL}{poporange!20}
\colorlet{popdarkL}{popdark!20}
\colorlet{poppurpleL}{poppurple!20}
\colorlet{popgreenL}{popgreen!20}
\colorlet{poppinkL}{poppink!20}
\colorlet{popblueL}{popblue!20}
\colorlet{popgoldL}{popgold!20}
\colorlet{popredL}{popred!20}
\colorlet{popgrayL}{popgray!20}
% Teintes claires pour fonds de tableaux / zones
\colorlet{popblueL}{popblue!10!white}
\colorlet{poporangeL}{poporange!14!white}
\colorlet{popgreenL}{popgreen!12!white}
\colorlet{poppurpleL}{poppurple!12!white}
\colorlet{poppinkL}{poppink!12!white}
\colorlet{popgoldL}{popgold!14!white}

\pagecolor{popcream}

% ============================================================
% Typographie
% ============================================================
\defaultfontfeatures{Ligatures=TeX}
\setmainfont{PlusJakartaSans-Medium.ttf}[Path=../../../fonts/,BoldFont=PlusJakartaSans-Bold.ttf]
% Maths en sans-serif (Fira Math) avec chiffres et lettres droites de Plus
% Jakarta, pour que les formules se fondent dans le texte.
\usepackage[math-style=ISO,bold-style=ISO]{unicode-math}
\setmathfont{FiraMath-Regular.otf}[Path=../../../fonts/]
\setmathfont{PlusJakartaSans-Medium.ttf}[Path=../../../fonts/,range={up/num,up/{Latin,latin},\mathup}]
\usepackage{icomma} % virgule décimale sans espace en mode math
% Caractères mathématiques Unicode absents des polices de texte (Plus Jakarta,
% Archivo Black) : rendus par leur équivalent mathématique, en texte comme en
% formule — sinon ils s'affichent en carré vide (ex. « Arithmétique dans ℤ »).
\usepackage{newunicodechar}
\newunicodechar{ℝ}{\ensuremath{\mathbb{R}}}
\newunicodechar{ℤ}{\ensuremath{\mathbb{Z}}}
\newunicodechar{ℂ}{\ensuremath{\mathbb{C}}}
\newunicodechar{ℕ}{\ensuremath{\mathbb{N}}}
\newunicodechar{ℚ}{\ensuremath{\mathbb{Q}}}
\newunicodechar{𝔻}{\ensuremath{\mathbb{D}}}
\newunicodechar{α}{\ensuremath{\alpha}}
\newunicodechar{β}{\ensuremath{\beta}}
\newunicodechar{γ}{\ensuremath{\gamma}}
\newunicodechar{δ}{\ensuremath{\delta}}
\newunicodechar{ε}{\ensuremath{\varepsilon}}
\newunicodechar{θ}{\ensuremath{\theta}}
\newunicodechar{λ}{\ensuremath{\lambda}}
\newunicodechar{μ}{\ensuremath{\mu}}
\newunicodechar{σ}{\ensuremath{\sigma}}
\newunicodechar{τ}{\ensuremath{\tau}}
\newunicodechar{φ}{\ensuremath{\varphi}}
\newunicodechar{ψ}{\ensuremath{\psi}}
\newunicodechar{ω}{\ensuremath{\omega}}
\newunicodechar{Σ}{\ensuremath{\Sigma}}
% majuscules produites par \MakeUppercase dans les titres (α-aminés -> Α)
\newunicodechar{Α}{\ensuremath{\alpha}}
\newunicodechar{Β}{\ensuremath{\beta}}
\newunicodechar{Γ}{\ensuremath{\Gamma}}
\newunicodechar{Δ}{\ensuremath{\Delta}}
\newunicodechar{Ε}{\ensuremath{\varepsilon}}
\newunicodechar{Θ}{\ensuremath{\Theta}}
\newunicodechar{Λ}{\ensuremath{\Lambda}}
\newunicodechar{Μ}{\ensuremath{\mu}}
\newunicodechar{Τ}{\ensuremath{\tau}}
\newunicodechar{Φ}{\ensuremath{\Phi}}
\newunicodechar{Ψ}{\ensuremath{\Psi}}
\newunicodechar{Ω}{\ensuremath{\Omega}}
\newunicodechar{↦}{\ensuremath{\mapsto}}
\newunicodechar{∈}{\ensuremath{\in}}
\newunicodechar{∉}{\ensuremath{\notin}}
\newunicodechar{∧}{\ensuremath{\wedge}}
\newunicodechar{⇔}{\ensuremath{\Leftrightarrow}}
\newunicodechar{⟺}{\ensuremath{\Longleftrightarrow}}
\newunicodechar{⇒}{\ensuremath{\Rightarrow}}
\newunicodechar{⟹}{\ensuremath{\Longrightarrow}}
\newunicodechar{∘}{\ensuremath{\circ}}
\newunicodechar{∪}{\ensuremath{\cup}}
\newunicodechar{∩}{\ensuremath{\cap}}
\newunicodechar{≡}{\ensuremath{\equiv}}
\newunicodechar{⊂}{\ensuremath{\subset}}
\newunicodechar{⊥}{\ensuremath{\perp}}
\newunicodechar{∃}{\ensuremath{\exists}}
\newunicodechar{∀}{\ensuremath{\forall}}
\newunicodechar{∛}{\ensuremath{\sqrt[3]{\,}}}
\newunicodechar{∜}{\ensuremath{\sqrt[4]{\,}}}
\newunicodechar{⃗}{\ensuremath{{}^{\to}}}
\newunicodechar{ⁿ}{\ifmmode^{n}\else\textsuperscript{\ensuremath{n}}\fi}
\newunicodechar{ˣ}{\ifmmode^{x}\else\textsuperscript{\ensuremath{x}}\fi}
\newunicodechar{ᵇ}{\ifmmode^{b}\else\textsuperscript{\ensuremath{b}}\fi}
\newunicodechar{ʳ}{\ifmmode^{r}\else\textsuperscript{\ensuremath{r}}\fi}
\newunicodechar{ᵘ}{\ifmmode^{u}\else\textsuperscript{\ensuremath{u}}\fi}
\newunicodechar{ʸ}{\ifmmode^{y}\else\textsuperscript{\ensuremath{y}}\fi}
\newunicodechar{ᵃ}{\ifmmode^{a}\else\textsuperscript{\ensuremath{a}}\fi}
\newunicodechar{ᵉ}{\ifmmode^{e}\else\textsuperscript{\ensuremath{e}}\fi}
\newunicodechar{ᵏ}{\ifmmode^{k}\else\textsuperscript{\ensuremath{k}}\fi}
\newunicodechar{ᵖ}{\ifmmode^{p}\else\textsuperscript{\ensuremath{p}}\fi}
\newunicodechar{ᴺ}{\ifmmode^{N}\else\textsuperscript{\ensuremath{N}}\fi}
\newunicodechar{ᶜ}{\ifmmode^{c}\else\textsuperscript{\ensuremath{c}}\fi}
\newunicodechar{ʰ}{\ifmmode^{h}\else\textsuperscript{\ensuremath{h}}\fi}
\newunicodechar{ᵠ}{\ifmmode^{q}\else\textsuperscript{\ensuremath{q}}\fi}
\newunicodechar{ₙ}{\ifmmode_{n}\else\textsubscript{\ensuremath{n}}\fi}
\newunicodechar{ₐ}{\ifmmode_{a}\else\textsubscript{\ensuremath{a}}\fi}
\newunicodechar{ᵢ}{\ifmmode_{i}\else\textsubscript{\ensuremath{i}}\fi}
\newunicodechar{ᵣ}{\ifmmode_{r}\else\textsubscript{\ensuremath{r}}\fi}
\newunicodechar{ₖ}{\ifmmode_{k}\else\textsubscript{\ensuremath{k}}\fi}
\newunicodechar{ᵦ}{\ifmmode_{\beta}\else\textsubscript{\ensuremath{\beta}}\fi}
\newunicodechar{ₓ}{\ifmmode_{x}\else\textsubscript{\ensuremath{x}}\fi}
\newfontfamily\popdisplay{ArchivoBlack-Regular.ttf}[Path=../../../fonts/]
\newfontfamily\poplogo{SpaceGrotesk-Bold.ttf}[Path=../../../fonts/]
\newfontfamily\popsemi{PlusJakartaSans-SemiBold.ttf}[Path=../../../fonts/]
\newfontfamily\popxbold{PlusJakartaSans-ExtraBold.ttf}[Path=../../../fonts/]
\newfontfamily\popxboldls{PlusJakartaSans-ExtraBold.ttf}[Path=../../../fonts/,LetterSpace=3.0]

\setlist[enumerate]{leftmargin=1.7em,itemsep=5pt,topsep=4pt}
\setlist[itemize]{leftmargin=1.7em,itemsep=4pt,topsep=4pt,label={\color{poporange}\textbullet}}
\setlength{\parindent}{0pt}
\setlength{\parskip}{5pt}
\renewcommand{\arraystretch}{1.25}

% pgfplots : style Pop par défaut
\pgfplotsset{
  every axis/.append style={axis line style={popink,line width=1pt},tick style={popink},
    grid style={popline,line width=0.5pt},label style={font=\small},tick label style={font=\footnotesize}},
  popcurve/.style={poporange,line width=1.8pt,no markers,smooth},
}
% Même style disponible dans un \draw TikZ simple (hors pgfplots)
\tikzset{popcurve/.style={poporange,line width=1.8pt}}
\tikzset{popgrid/.style={popline,very thin}}
\colorlet{popcurve}{poporange} % le nom du style est parfois utilisé comme couleur

% ============================================================
% Pill / squiggle
% ============================================================
\newcommand{\poppill}[3]{%
  \tikz[baseline=-0.55ex]\node[draw=popink,line width=1.2pt,rounded corners=2.6pt,%
    fill=#2,inner xsep=6.5pt,inner ysep=3pt]{%
    \popxboldls\fontsize{7}{7}\selectfont\color{#3}\MakeUppercase{#1}};%
}
\newcommand{\popsquiggle}[1]{%
  \tikz[baseline=-0.4ex]{
    \draw[#1,line width=2.6pt,line cap=round]
      plot [smooth] coordinates {(0,0.09) (0.34,0.01) (0.68,0.21) (1.02,0.01) (1.36,0.10)};
  }%
}
% Mot-clé surligné (marqueur orange sous le texte)
\newcommand{\cle}[1]{{\popxbold\color{popdark}#1}}

% ============================================================
% En-tête / pied
% ============================================================
\pagestyle{fancy}
\fancyhf{}
\renewcommand{\headrulewidth}{0pt}
\renewcommand{\footrulewidth}{0pt}
\lhead{\raisebox{-0.15ex}{\poplogo\fontsize{13}{13}\selectfont\color{popink}OnBuch\color{poporange}+}}
\rhead{\poppill{\DOCMATIERE\ \textbullet\ \DOCNIVEAU\ \textbullet\ Leçon \DOCLECON}{white}{popink}}
\lfoot{\popxbold\fontsize{7.6}{7.6}\selectfont\color{popmuted}\MakeUppercase{Cours signé OnBuch+}\ \textbullet\ {\popsemi\DOCTITRE}}
\rfoot{\tikz[baseline=-0.6ex]\node[rounded corners=8.5pt,fill=popink,minimum height=17pt,minimum width=17pt,inner xsep=5pt,inner sep=0]{\popxbold\fontsize{8}{8}\selectfont\color{popcream}\thepage\,/\,\pageref*{LastPage}};}

% ============================================================
% Titres
% ============================================================
\newcommand{\chaptitre}[2]{%
  \begin{tcolorbox}[enhanced,colback=poporange,colframe=popink,boxrule=2.6pt,arc=11pt,
      drop shadow=popink,left=15pt,right=15pt,top=12pt,bottom=11pt,before skip=3pt,after skip=18pt]
    \poppill{#2}{popink}{popcream}\par\vspace{9pt}
    {\raggedright\hyphenpenalty=10000\exhyphenpenalty=10000\popdisplay\fontsize{22}{24}\selectfont\color{popcream}\MakeUppercase{#1}\par}
  \end{tcolorbox}
}
% Section numérotée : pastille orange avec le numéro + titre Archivo
\newcounter{popsec}
\newcommand{\coursec}[1]{%
  \stepcounter{popsec}\setcounter{popsub}{0}%
  \par\vspace{16pt}\noindent
  \begin{minipage}{\linewidth}
  \tikz[baseline=-0.7ex]\node[draw=popink,line width=1.6pt,fill=poporange,rounded corners=4pt,minimum size=22pt,inner sep=0]{\popdisplay\fontsize{12}{12}\selectfont\color{popcream}\thepopsec};%
  \hspace{8pt}{\popdisplay\fontsize{15}{17}\selectfont\color{popink}\MakeUppercase{#1}}\par
  \vspace{4pt}\noindent{\color{poporange}\rule{36pt}{3.6pt}}
  \end{minipage}\par\nopagebreak\vspace{8pt}
}
\newcounter{popsub}
\newcommand{\courssub}[1]{%
  \stepcounter{popsub}%
  \par\vspace{9pt}\noindent{\popxbold\fontsize{11.5}{13}\selectfont\color{popink}{\color{poporange}\thepopsec.\thepopsub}\hspace{6pt}#1}\par\nopagebreak\vspace{4pt}
}

% ============================================================
% Boîtes pédagogiques — même recette : fond blanc, bordure épaisse colorée,
% ombre dure encre, badge plein. Titre optionnel [#1].
% ============================================================
\tcbset{popbase/.style={breakable,enhanced,colback=white,boxrule=2.3pt,arc=9pt,
  shadow={3pt}{-3pt}{0pt}{popink},left=14pt,right=14pt,top=11pt,bottom=11pt,before skip=11pt,after skip=15pt,
  fontupper=\popsemi\fontsize{10.6}{14.5}\selectfont\color{popink},
  fontlower=\popsemi\fontsize{10.6}{14.5}\selectfont\color{popink}}}
% Coche dessinée (le glyphe ✓ n'existe pas dans les polices du document)
\newcommand{\popcheck}[1]{\tikz[baseline=-0.5ex]\draw[#1,line width=1.6pt,line cap=round,line join=round](-0.13,0.0)--(-0.04,-0.09)--(0.13,0.1);}
\providecommand{\coche}{\popcheck{popgreen}}
\providecommand{\resultat}[1]{\par\smallskip\noindent\fcolorbox{popgreen}{popgreenL}{\parbox{\dimexpr\linewidth-2\fboxsep-2\fboxrule\relax}{\textbf{Résultat.} #1}}\par\smallskip}
\newcommand{\popboxhead}[3]{\poppill{#1}{#2}{white}\ifx\relax#3\relax\else\hspace{6pt}{\popxbold\fontsize{10.6}{12}\selectfont\color{popink}#3}\fi\par\vspace{7pt}}

\newtcolorbox{aretenir}[1][]{popbase,colframe=popgreen,before upper={\popboxhead{À retenir}{popgreen}{#1}}}
\newtcolorbox{exemplebox}[1][]{popbase,colframe=poporange,before upper={\popboxhead{Exemple}{poporange}{#1}}}
\newtcolorbox{definition}[1][]{popbase,colframe=popblue,before upper={\popboxhead{Définition}{popblue}{#1}}}
\newtcolorbox{propriete}[1][]{popbase,colframe=poppurple,before upper={\popboxhead{Propriété}{poppurple}{#1}}}
\newtcolorbox{methode}[1][]{popbase,colframe=popgold,before upper={\popboxhead{Méthode}{popgold}{#1}}}
\newtcolorbox{attention}[1][]{popbase,colframe=poppink,before upper={\popboxhead{Attention}{poppink}{#1}}}
\newtcolorbox{experience}[1][]{popbase,colframe=popblue,colback=popblueL,before upper={\popboxhead{Expérience}{popblue}{#1}}}
% « Le savais-tu ? » : carte violette pleine (contexte, culture, Cameroun)
\newtcolorbox{savaistu}[1][]{popbase,colback=poppurple,colframe=popink,
  fontupper=\popsemi\fontsize{10.4}{14.2}\selectfont\color{white},
  before upper={\poppill{Le savais-tu\,?}{white}{poppurple}\ifx\relax#1\relax\else\hspace{6pt}{\popxbold\color{white}#1}\fi\par\vspace{7pt}}}
% Exercice résolu : énoncé en haut, \tcblower puis la solution sur fond crème
\newcounter{popexo}\newcounter{popexr}
\newtcolorbox{exoresolu}[1][]{popbase,colframe=poporange,colbacklower=popcream,
  segmentation style={popink,line width=1.2pt,dashed},
  before upper={\stepcounter{popexr}\popboxhead{Exercice résolu \thepopexr}{poporange}{#1}},
  before lower={\poppill{Solution}{popink}{popcream}\par\vspace{6pt}}}
% Exercice (non résolu en place) — la correction est dans la partie Corrigés
\newtcolorbox{exercice}[1][]{popbase,colframe=popink,
  before upper={\stepcounter{popexo}\popboxhead{Exercice \thepopexo}{popink}{#1}}}
% Corrigé
\newtcolorbox{corrige}[1][]{popbase,colframe=popgreen,colback=popgreenL,
  before upper={\popboxhead{Corrigé}{popgreen}{#1}}}
% Objectifs / prérequis / bilan
\newtcolorbox{objectifs}{popbase,colframe=popink,colback=poporangeL,
  before upper={\popboxhead{Objectifs de la leçon}{popink}{}}}
\newtcolorbox{prerequis}{popbase,colframe=popink,colback=popblueL,
  before upper={\popboxhead{Prérequis}{popblue}{}}}
\newtcolorbox{bilan}{popbase,colframe=popink,colback=white,boxrule=2.8pt,
  before upper={\popboxhead{Fiche bilan}{poporange}{L'essentiel à connaître pour l'examen}}}
% Figure encadrée avec légende
\newtcolorbox{popfigure}[1][]{enhanced,breakable=false,colback=white,colframe=popline,boxrule=1.4pt,arc=8pt,
  left=10pt,right=10pt,top=10pt,bottom=8pt,before skip=10pt,after skip=12pt,halign=center,
  lower separated=false,halign lower=center,
  fontlower=\popsemi\fontsize{9}{11.5}\selectfont\color{popmuted},
  #1}
\newcommand{\legende}[1]{\par\vspace{6pt}{\popsemi\fontsize{9}{11.5}\selectfont\color{popmuted}#1}}

% ============================================================
% Signature OnBuch+
% ============================================================
\newcommand{\poplogoinline}[1]{{\poplogo\fontsize{#1}{#1}\selectfont\color{popink}OnBuch\color{poporange}+}}
\newcommand{\signatureonbuch}{%
  \begin{tcolorbox}[enhanced,colback=popink,colframe=popink,boxrule=2.2pt,arc=10pt,
      drop shadow=poporange,left=14pt,right=14pt,top=10pt,bottom=10pt,before skip=0pt,after skip=0pt]
    \tikz[baseline=-0.7ex]\node[circle,fill=popgreen,draw=popcream,line width=1.3pt,minimum size=22pt,inner sep=0]{\popcheck{white}};%
    \hspace{9pt}\begin{minipage}[c]{0.62\linewidth}
      {\popxbold\fontsize{12}{13}\selectfont\color{popcream}Signé\ \textbullet\ {\poplogo OnBuch\color{poporange}+}}\par\vspace{2pt}
      {\raggedright\popsemi\fontsize{8}{10.5}\selectfont\color{popline}\MakeUppercase{Équipe pédagogique OnBuch+}\\\MakeUppercase{Conforme au programme officiel MINESEC}\par}
    \end{minipage}\hfill
    {\poplogo\fontsize{22}{22}\selectfont\color{popcream}OnBuch\color{poporange}+}
  \end{tcolorbox}%
}

% ============================================================
% Page de garde
% #1 titre · #2 matière · #3 niveau · #4 module · #5 n° leçon · #6 durée ·
% #7 description / objectifs (texte ou itemize)
% ============================================================
\newcommand{\pagegarde}[7]{%
  \thispagestyle{empty}
  \newgeometry{a4paper,margin=1.7cm,top=1.6cm,bottom=1.4cm}%
  \begin{tikzpicture}[remember picture,overlay]
    \fill[poporange!18] ([xshift=-3.2cm,yshift=-3.6cm]current page.north east) circle (4.6cm);
    \fill[poppurple!14] ([xshift=2.4cm,yshift=7.6cm]current page.south west) circle (3.8cm);
  \end{tikzpicture}%
  {\poplogo\fontsize{30}{30}\selectfont\color{popink}OnBuch\color{poporange}+}\hfill
  \raisebox{0.6ex}{\poppill{Cours signé \textbullet\ Premium}{popink}{popcream}}\par
  \vspace{1pt}\popsquiggle{poppurple}\par
  \vspace{8pt}
  \begin{tcolorbox}[enhanced,colback=poporange,colframe=popink,boxrule=2.8pt,arc=13pt,
      drop shadow=popink,left=18pt,right=18pt,top=12pt,bottom=13pt,after skip=10pt]
    \poppill{#2 \textbullet\ #3}{popink}{popcream}\hspace{5pt}\poppill{#4}{white}{popink}\par\vspace{8pt}
    {\popdisplay\fontsize{14}{14}\selectfont\color{popink}LEÇON #5}\par\vspace{4pt}
    {\raggedright\hyphenpenalty=10000\exhyphenpenalty=10000\popdisplay\fontsize{25}{27}\selectfont\color{popcream}\MakeUppercase{#1}\par}
    \vspace{8pt}
    {\popxbold\fontsize{9.5}{9.5}\selectfont\color{popink}\MakeUppercase{Durée conseillée : #6}}
  \end{tcolorbox}
  \begin{tcolorbox}[enhanced,colback=white,colframe=popink,boxrule=2.2pt,arc=10pt,
      drop shadow=popink,left=16pt,right=16pt,top=10pt,bottom=10pt,after skip=8pt]
    \poppill{Dans cette leçon}{poppurple}{white}\par\vspace{5pt}
    \popsemi\fontsize{9.3}{12.6}\selectfont\color{popink}#7
  \end{tcolorbox}
  \noindent\begin{minipage}[t]{0.487\linewidth}
    \begin{tcolorbox}[enhanced,colback=popgreen,colframe=popink,boxrule=2.2pt,arc=10pt,
        drop shadow=popink,left=11pt,right=11pt,top=10pt,bottom=10pt,height=3.1cm,valign=center]
      \tcbox[colback=white,colframe=popink,boxrule=1.3pt,arc=5pt,boxsep=0pt,nobeforeafter,
        left=3pt,right=3pt,top=3pt,bottom=3pt,valign=center]{\qrcode[height=1.9cm]{https://play.google.com/store/apps/details?id=com.luvvix.onbuch&hl=fr}}%
      \hspace{8pt}\begin{minipage}[c]{0.5\linewidth}
        {\popdisplay\fontsize{11}{12.5}\selectfont\color{white}\MakeUppercase{Télécharge OnBuch}}\par\vspace{4pt}
        {\raggedright\popsemi\fontsize{8.4}{11}\selectfont\color{white}Gratuit \textbullet\ Google Play\\Tuteur IA, annales, concours, résultats\par}
      \end{minipage}
    \end{tcolorbox}
  \end{minipage}\hfill
  \begin{minipage}[t]{0.487\linewidth}
    \begin{tcolorbox}[enhanced,colback=poppurple,colframe=popink,boxrule=2.2pt,arc=10pt,
        drop shadow=popink,left=11pt,right=11pt,top=10pt,bottom=10pt,height=3.1cm,valign=center]
      \tikz[baseline=-0.7ex]\node[circle,fill=white,draw=popink,line width=1.3pt,minimum size=1.6cm,inner sep=0]{\popdisplay\fontsize{24}{24}\selectfont\color{poppurple}+};%
      \hspace{8pt}\begin{minipage}[c]{0.6\linewidth}
        {\popdisplay\fontsize{11}{12.5}\selectfont\color{white}\MakeUppercase{Passe à OnBuch+}}\par\vspace{4pt}
        {\raggedright\popsemi\fontsize{8.4}{11}\selectfont\color{white}Tous les cours signés, Tuteur IA illimité, fiches PDF — onglet OnBuch+ de l'app\par}
      \end{minipage}
    \end{tcolorbox}
  \end{minipage}
  \vspace{8pt}
  \popcontenu
  \vfill
  \signatureonbuch
  \restoregeometry
  \newpage
}

% Rangée « Ce cours contient » (page de garde)
\newcommand{\poptile}[3]{% #1 couleur #2 gros texte #3 légende
  \begin{tcolorbox}[enhanced,colback=#1,colframe=popink,boxrule=2pt,arc=9pt,shadow={2.5pt}{-2.5pt}{0pt}{popink},
      left=6pt,right=6pt,top=9pt,bottom=9pt,halign=center,valign=center,height=1.75cm,width=\linewidth,nobeforeafter]
    {\popdisplay\fontsize{12.5}{13}\selectfont\color{white}\MakeUppercase{#2}}\par\vspace{3pt}
    {\centering\popsemi\fontsize{7.8}{9.5}\selectfont\color{white}#3\par}
  \end{tcolorbox}}
\newcommand{\popcontenu}{%
  \par\noindent{\popxboldls\fontsize{7.5}{7.5}\selectfont\color{popmuted}CE COURS SIGNÉ CONTIENT}\par\vspace{7pt}
  \noindent\begin{minipage}[t]{0.235\linewidth}\poptile{popblue}{Cours}{complet, illustré, pas à pas}\end{minipage}\hfill
  \begin{minipage}[t]{0.235\linewidth}\poptile{poporange}{Méthodes}{et exercices résolus}\end{minipage}\hfill
  \begin{minipage}[t]{0.235\linewidth}\poptile{poppink}{Exercices}{type examen + corrigés}\end{minipage}\hfill
  \begin{minipage}[t]{0.235\linewidth}\poptile{popgreen}{Fiche bilan}{l'essentiel à retenir}\end{minipage}\par}

% Sommaire
\newtcolorbox{sommaire}{enhanced,colback=white,colframe=popink,boxrule=2.2pt,arc=10pt,
  drop shadow=popink,left=18pt,right=18pt,top=13pt,bottom=13pt,before skip=6pt,after skip=20pt,
  before upper={\poppill{Sommaire}{poppurple}{white}\par\vspace{9pt}\popsemi\fontsize{10.6}{14.5}\selectfont\color{popink}}}

% Signature de fin de cours — poussée tout en bas de la dernière page
\newcommand{\findecours}{%
  \par\vfill\nopagebreak
  \begin{center}\popsquiggle{poporange}\end{center}\vspace{4pt}
  \signatureonbuch
}
\AtBeginDocument{\def\llbracket{\mathopen{[\mkern-3.5mu[}}\def\rrbracket{\mathclose{]\mkern-3.5mu]}}} % intervalles d'entiers (glyphe ⟦ absent de la police)
% Glyphes absents de Fira Math : redéfinis à partir de glyphes disponibles
\AtBeginDocument{%
  \def\setminus{\mathbin{\backslash}}\let\smallsetminus\setminus
  \def\vdots{\mathord{\vbox{\baselineskip3.2pt\lineskiplimit0pt\kern4pt\hbox{.}\hbox{.}\hbox{.}}}}%
  \def\ddots{\mathinner{\mkern1mu\raise6.5pt\hbox{.}\mkern2mu\raise3.6pt\hbox{.}\mkern2mu\raise0.7pt\hbox{.}\mkern1mu}}%
  \def\triangle{\mathord{\scalebox{0.9}{$\symup{\Delta}$}}}%
  \def\nabla{\mathord{\raisebox{\depth}{\scalebox{1}[-1]{$\symup{\Delta}$}}}}%
  \def\checkmark{\popcheck{popdark}}%
  \def\star{\mathbin{\ast}}\def\bigstar{\mathord{\ast}}%
  \def\diamond{\mathbin{\scalebox{0.6}{$\Diamond$}}}%
  \def\bigcup{\mathop{\mathchoice{\raisebox{-0.3em}{\scalebox{1.7}{$\cup$}}}{\scalebox{1.25}{$\cup$}}{\cup}{\cup}}\displaylimits}%
  \def\bigcap{\mathop{\mathchoice{\raisebox{-0.3em}{\scalebox{1.7}{$\cap$}}}{\scalebox{1.25}{$\cap$}}{\cap}{\cap}}\displaylimits}%
  \def\lesseqqgtr{\mathrel{\lessgtr}}\def\bigoplus{\mathop{\oplus}\displaylimits}%
}
\providecommand{\ssi}{si et seulement si} % abréviation fréquente
\AtBeginDocument{\providecommand{\wideparen}{\overparen}} % arc de cercle

%% FIGFIX-BEGIN v5 — correctifs globaux des figures (docs/onbuch-generation/tools/figfix)
\usepackage{adjustbox}\usepackage{etoolbox}\tcbuselibrary{raster}
% toute figure TikZ/pgfplots plus large que la ligne est réduite à la largeur disponible
\BeforeBeginEnvironment{tikzpicture}{\begin{adjustbox}{max width=\linewidth}}
\AfterEndEnvironment{tikzpicture}{\end{adjustbox}}
% légendes pgfplots lisibles par-dessus les courbes
\pgfplotsset{every axis/.append style={legend style={fill=white,fill opacity=0.92,text opacity=1,draw=black!15,inner sep=3pt}}}
% étiquettes d'arêtes (syntaxe edge["..."]) avec fond blanc
\tikzset{every edge quotes/.append style={fill=white,inner sep=1.2pt}}
%% FIGFIX-END
\begin{document}

\pagegarde{Localisation et nature de l’information génétique déterminant les caractères héréditaires}{SVTEEHB}{3ème}{SVTEEHB – classe de 3ème}{N02}{4 séances (fiche de progression harmonisée)}{Cette leçon établit, à partir d'expériences historiques de transfert de noyau, que l'information génétique responsable des caractères héréditaires est localisée dans le noyau des cellules. Elle précise ensuite la nature de cette information : les chromosomes, supports de l'hérédité, constitués d'ADN.
\vspace{4pt}
\begin{multicols}{2}\small
\begin{itemize}
\item À la fin de cette leçon, je sais décrire et interpréter les expériences de transfert de noyau (Hammerling, Gurdon) pour localiser l'information génétique.
\item À la fin de cette leçon, je sais conclure, à partir de résultats expérimentaux, que l'information génétique est localisée dans le noyau.
\item À la fin de cette leçon, je sais définir le chromosome et décrire sa structure (chromatides, centromère) et ses états (monochromatidien, bichromatidien).
\item À la fin de cette leçon, je sais définir l'ADN et décrire sa structure en double hélice (nucléotides, bases A-T-G-C).
\item À la fin de cette leçon, je sais réaliser le protocole de mise en évidence du matériel génétique d'une cellule (TP N°1).
\item À la fin de cette leçon, je sais concevoir des maquettes de paires de chromosomes homologues monochromatidiens et bichromatidiens (TP N°2).
\end{itemize}\end{multicols}}

\begin{sommaire}
\begin{multicols}{2}
\begin{enumerate}
\item Problème et hypothèses : où se trouve l'information génétique ?
\item Localisation de l'information génétique : les expériences de transfert de noyau
\item Les chromosomes, supports de l'information génétique
\item L'ADN, molécule de l'information génétique
\item TP N°1 : Mise en évidence du matériel génétique d'une cellule
\item TP N°2 : Maquettes de paires de chromosomes homologues et applications
\item Activité d'intégration
\item Méthodes et exercices résolus
\item Exercices
\item Corrigés des exercices
\item Fiche bilan
\end{enumerate}
\end{multicols}
\end{sommaire}

\begin{objectifs}
\begin{itemize}
\item À la fin de cette leçon, je sais décrire et interpréter les expériences de transfert de noyau (Hammerling, Gurdon) pour localiser l'information génétique.
\item À la fin de cette leçon, je sais conclure, à partir de résultats expérimentaux, que l'information génétique est localisée dans le noyau.
\item À la fin de cette leçon, je sais définir le chromosome et décrire sa structure (chromatides, centromère) et ses états (monochromatidien, bichromatidien).
\item À la fin de cette leçon, je sais définir l'ADN et décrire sa structure en double hélice (nucléotides, bases A-T-G-C).
\item À la fin de cette leçon, je sais réaliser le protocole de mise en évidence du matériel génétique d'une cellule (TP N°1).
\item À la fin de cette leçon, je sais concevoir des maquettes de paires de chromosomes homologues monochromatidiens et bichromatidiens (TP N°2).
\item À la fin de cette leçon, je sais appliquer ces notions à des situations concrètes (ressemblances familiales, tests de paternité, caryotype).
\item À la fin de cette leçon, je sais rédiger et justifier une démarche scientifique, une réponse ou une conclusion.
\end{itemize}
\end{objectifs}

\begin{prerequis}
\begin{itemize}
\item La cellule : structure générale (membrane, cytoplasme, noyau) vue en 4ème
\item Notion de caractères héréditaires et de transmission des caractères (leçon N°1 de 3ème)
\item Reproduction sexuée : gamètes, fécondation (4ème)
\item Observation au microscope et réalisation d'un protocole expérimental simple
\item Notion de molécule et d'éléments chimiques (C, H, O, N, P)
\end{itemize}
\end{prerequis}

\newpage

\coursec{Problème et hypothèses : où se trouve l'information génétique ?}

Dans nos villages, on dit souvent : \emph{« cet enfant a les yeux de son père et le teint de sa mère »}. À la maternité de l'hôpital de district, une famille s'interroge : le nouveau-né présente une particularité qu'on ne retrouve chez personne dans la famille proche. Chacun y va de son commentaire, mais personne ne peut répondre à la vraie question : \textbf{où se trouve donc, dans nos cellules, l'« information » qui détermine nos ressemblances et nos différences ? Et de quoi cette information est-elle faite ?}

C'est à ces questions que cette leçon va répondre, en suivant le cheminement des savants eux-mêmes : des expériences célèbres de \cle{transfert de noyau} et l'observation des \cle{chromosomes} vont nous permettre de localiser cette information, puis d'en découvrir la nature.

\courssub{Les caractères héréditaires : un constat de la vie quotidienne}

Observe autour de toi : dans une même famille, les enfants ressemblent à leurs parents, mais jamais totalement, et ils diffèrent aussi entre frères et sœurs. Certains traits semblent « passer » d'une génération à l'autre : la couleur des yeux, la forme du nez, le teint de la peau, la forme du lobe de l'oreille (collé ou décollé), le groupe sanguin.

\begin{definition}[Caractère héréditaire]
Un \cle{caractère héréditaire} est un caractère (une particularité observable d'un individu) qui se \textbf{transmet des parents à la descendance}, c'est-à-dire d'une génération à la suivante.
\end{definition}

\begin{exemplebox}[Des caractères héréditaires chez l'Homme]
Chez l'Homme, on peut citer de nombreux caractères héréditaires :
\begin{itemize}
\item le \textbf{groupe sanguin} : un individu est de groupe A, B, AB ou O, et ce caractère se transmet des parents aux enfants selon des règles précises ;
\item la \textbf{couleur des yeux} (yeux bruns, noirs, clairs...) ;
\item la \textbf{forme du lobe de l'oreille} : lobe décollé ou lobe collé ;
\item la \textbf{couleur de la peau}, la forme des cheveux (crépus, ondulés, lisses).
\end{itemize}
Ces caractères ne s'acquièrent pas au cours de la vie : un enfant naît déjà avec son groupe sanguin, déterminé avant sa naissance.
\end{exemplebox}

\begin{attention}[Héréditaire ou acquis ?]
Ne confonds pas caractère \textbf{héréditaire} et caractère \textbf{acquis}. Une cicatrice, une amputation, des muscles développés par le sport sont des caractères \emph{acquis} au cours de la vie : ils \textbf{ne se transmettent pas} aux enfants. Seuls les caractères inscrits dans l'information reçue des parents sont héréditaires.
\end{attention}

Puisque ces caractères passent des parents aux enfants, c'est qu'il existe, quelque part, une \textbf{information} qui les détermine et qui est transmise lors de la reproduction. Cette information s'appelle l'\cle{information génétique}. Mais où est-elle rangée ?

\courssub{La cellule et ses constituants : où chercher ?}

Tu as appris dans les classes précédentes que tout être vivant est constitué de \textbf{cellules}. Le corps humain en contient des dizaines de milliers de milliards ! Or, chacun de nous s'est formé à partir d'\textbf{une seule cellule} : la cellule-œuf (ou \cle{zygote}), issue de l'union d'un spermatozoïde du père et d'un ovule de la mère. C'est donc dans cette cellule, et dans ses constituants, que l'information génétique doit se trouver.

Rappelons la structure d'une cellule animale typique. Elle comprend trois constituants principaux :
\begin{itemize}
\item la \cle{membrane cellulaire} : fine enveloppe qui délimite la cellule et contrôle les échanges avec le milieu extérieur ;
\item le \cle{cytoplasme} : milieu gélatineux où se déroulent les activités chimiques de la cellule (respiration cellulaire, fabrication de substances...) ;
\item le \cle{noyau} : petite sphère située généralement au centre de la cellule, longtemps restée mystérieuse.
\end{itemize}

\begin{popfigure}
\begin{center}
\begin{tikzpicture}[scale=1.0]
% Membrane et cytoplasme
\fill[popblueL] (0,0) ellipse (4.2 and 2.9);
\draw[popblue, very thick] (0,0) ellipse (4.2 and 2.9);
% Noyau
\fill[poporangeL] (-0.3,0.2) circle (1.15);
\draw[poporange, very thick] (-0.3,0.2) circle (1.15);
\fill[poporange!60] (-0.3,0.2) circle (0.3);
% Point d'interrogation sur le noyau
\node[popdark, font=\Large\bfseries] at (-0.3,0.85) {?};
% Petites structures dans le cytoplasme (mitochondries)
\fill[popgreen!50] (2.2,1.2) ellipse (0.35 and 0.18);
\fill[popgreen!50] (-2.6,-1.0) ellipse (0.35 and 0.18);
\fill[popgreen!50] (1.8,-1.4) ellipse (0.35 and 0.18);
% Étiquettes
\draw[->, popdark, thick] (4.2,1.4) -- (5.3,2.0) node[right, font=\small] {\textbf{1. Membrane cellulaire}};
\draw[->, popdark, thick] (2.6,-1.9) -- (4.0,-2.6) node[right, font=\small] {\textbf{2. Cytoplasme}};
\draw[->, popdark, thick] (-0.3,-0.95) -- (-1.5,-2.6) node[below, font=\small] {\textbf{3. Noyau} : l'information est-elle ici ?};
\end{tikzpicture}
\end{center}
\legende{Schéma d'une cellule animale montrant ses trois constituants principaux : (1) la membrane cellulaire, (2) le cytoplasme avec quelques mitochondries, (3) le noyau avec son nucléoïde. Dans lequel de ces constituants se trouve l'information qui détermine les caractères héréditaires ?}
\end{popfigure}

\courssub{Formulation du problème}

Nous voici donc face à une question précise, que l'on peut formuler comme un véritable \textbf{problème scientifique} :

\begin{aretenir}[Le problème posé]
\textbf{Dans quelle partie de la cellule est localisée l'information qui détermine les caractères héréditaires ?} Est-ce la membrane ? Le cytoplasme ? Le noyau ?
\end{aretenir}

Remarque bien la démarche : on part d'un \textbf{constat} (les caractères se transmettent des parents aux enfants), on identifie les \textbf{candidats possibles} (les trois constituants de la cellule), et on pose une \textbf{question} à laquelle seule l'expérience pourra répondre.

\courssub{Émission d'hypothèses vérifiables}

En science, pour répondre à un problème, on commence par émettre des \cle{hypothèses} : ce sont des \textbf{réponses possibles, provisoires}, que l'on va ensuite soumettre à l'épreuve de l'expérience. Face à notre problème, trois hypothèses sont envisageables :

\begin{itemize}
\item \textbf{Hypothèse 1 (H1)} : l'information génétique est localisée dans le \textbf{noyau} de la cellule ;
\item \textbf{Hypothèse 2 (H2)} : l'information génétique est localisée dans le \textbf{cytoplasme} ;
\item \textbf{Hypothèse 3 (H3)} : l'information génétique est localisée dans la \textbf{membrane cellulaire}.
\end{itemize}

\begin{attention}[Une hypothèse n'est pas une conclusion !]
Une hypothèse est une réponse \textbf{provisoire}, une supposition. Elle n'a aucune valeur tant qu'elle n'a pas été \textbf{testée par l'expérience}. Dire « je pense que c'est le noyau » ne prouve rien : il faut une expérience dont les résultats permettent de \emph{valider} (confirmer) ou d'\emph{invalider} (rejeter) l'hypothèse. Au BEPC, ne rédige jamais une conclusion avant d'avoir interprété les résultats expérimentaux !
\end{attention}

\courssub{La démarche expérimentale : comment trancher ?}

\begin{methode}[La démarche d'investigation scientifique en 5 étapes]
Pour résoudre un problème scientifique, on suit toujours la même démarche rigoureuse :
\begin{enumerate}
\item \textbf{Poser le problème} : formuler clairement la question à partir d'observations ;
\item \textbf{Émettre des hypothèses} : proposer des réponses possibles et vérifiables ;
\item \textbf{Concevoir une expérience} : imaginer un protocole capable de tester les hypothèses, et prévoir les résultats attendus pour chacune ;
\item \textbf{Réaliser l'expérience et recueillir les résultats} : observer, mesurer, noter avec précision ;
\item \textbf{Interpréter et conclure} : comparer les résultats obtenus aux résultats attendus, valider ou invalider les hypothèses, puis répondre au problème.
\end{enumerate}
\end{methode}

Appliquons cette démarche à notre problème. Comment savoir si c'est le noyau ou le cytoplasme qui porte l'information ? L'idée géniale des chercheurs consiste à \textbf{séparer puis échanger} le noyau et le cytoplasme entre deux cellules d'individus différents : c'est la technique du \cle{transfert de noyau} (ou greffe de noyau).

Le principe est le suivant : on prend deux individus qui se distinguent par un caractère héréditaire bien visible (par exemple la couleur). On retire le noyau d'une cellule du premier individu et on le place dans une cellule du second individu dont on a enlevé le noyau. On obtient ainsi une cellule « recombinée » : \textbf{noyau de l'un, cytoplasme de l'autre}. On observe ensuite quel caractère se développe :
\begin{itemize}
\item si le caractère obtenu est celui du \textbf{donneur de noyau}, alors l'information est dans le noyau ;
\item si le caractère obtenu est celui du \textbf{donneur de cytoplasme}, alors l'information est dans le cytoplasme.
\end{itemize}

\courssub{Rédaction d'un protocole et des résultats attendus}

Un bon protocole expérimental doit prévoir, \emph{avant} de faire l'expérience, les résultats attendus \textbf{pour chaque hypothèse}. Ainsi, quel que soit le résultat observé, on saura l'interpréter.

\begin{experience}[Protocole de transfert de noyau et résultats attendus]
\textbf{Matériel} : deux individus A et B d'une même espèce, différant par un caractère héréditaire visible ; microscope ; micropipettes pour manipuler les noyaux.

\textbf{Protocole} :
\begin{enumerate}
\item Prélever une cellule chez l'individu A et en extraire le noyau ;
\item Prélever une cellule chez l'individu B et retirer son noyau (cellule \emph{énucléée}) ;
\item Introduire le noyau de A dans la cellule énucléée de B ;
\item Suivre le développement de cette cellule recombinée et observer le caractère obtenu.
\end{enumerate}

\textbf{Résultats attendus selon les hypothèses} :
\begin{center}
\begin{tabularx}{\linewidth}{|l|X|X|}
\hline
\rowcolor{popblueL} \textbf{Hypothèse} & \textbf{Résultat attendu} & \textbf{Interprétation} \\
\hline
H1 : information dans le noyau & Le caractère obtenu est celui de A (donneur de noyau) & L'information est dans le noyau \\
\hline
H2 : information dans le cytoplasme & Le caractère obtenu est celui de B (donneur de cytoplasme) & L'information est dans le cytoplasme \\
\hline
H3 : information dans la membrane & Le caractère obtenu est celui de B (donneur de membrane et cytoplasme) & L'information est dans la membrane \\
\hline
\end{tabularx}
\end{center}
\end{experience}

\begin{popfigure}
\begin{center}
\begin{tikzpicture}[
  hyp/.style={rectangle, rounded corners, draw=popblue, fill=popblueL, text width=3.1cm, align=center, font=\small},
  exp/.style={rectangle, rounded corners, draw=poporange, fill=poporangeL, text width=3.4cm, align=center, font=\small},
  res/.style={rectangle, rounded corners, draw=popgreen, fill=popgreenL, text width=3.4cm, align=center, font=\small},
  arr/.style={->, thick, popdark}
]
\node[hyp] (h1) at (0,3.2) {\textbf{H1} : l'information est dans le \textbf{noyau}};
\node[hyp] (h2) at (0,0) {\textbf{H2} : l'information est dans le \textbf{cytoplasme}};
\node[hyp] (h3) at (0,-3.2) {\textbf{H3} : l'information est dans la \textbf{membrane}};
\node[exp] (e1) at (5.6,3.2) {Transfert du noyau de A dans la cellule énucléée de B};
\node[exp] (e2) at (5.6,0) {Même expérience de transfert de noyau};
\node[exp] (e3) at (5.6,-3.2) {Même expérience de transfert de noyau};
\node[res] (r1) at (11.4,3.2) {Caractère de \textbf{A} (donneur du noyau)};
\node[res] (r2) at (11.4,0) {Caractère de \textbf{B} (donneur du cytoplasme)};
\node[res] (r3) at (11.4,-3.2) {Caractère de \textbf{B} (donneur de la membrane)};
\draw[arr] (h1) -- (e1);
\draw[arr] (h2) -- (e2);
\draw[arr] (h3) -- (e3);
\draw[arr] (e1) -- (r1);
\draw[arr] (e2) -- (r2);
\draw[arr] (e3) -- (r3);
\node[font=\small\itshape, popmuted] at (0,4.6) {Hypothèses};
\node[font=\small\itshape, popmuted] at (5.6,4.6) {Expérience prévue};
\node[font=\small\itshape, popmuted] at (11.4,4.6) {Résultat attendu};
\end{tikzpicture}
\end{center}
\legende{Arbre de décision : chaque hypothèse (H1, H2, H3) est testée par la même expérience de transfert de noyau, mais chacune prédit un résultat attendu différent. La comparaison du résultat réel aux résultats attendus permettra de conclure.}
\end{popfigure}

\begin{exoresolu}[Rédiger et justifier une démarche]
\textbf{Énoncé.} Un élève de 3ème affirme : « L'information génétique se trouve dans le cytoplasme, car c'est là que la cellule fabrique toutes ses substances. » On dispose de deux mouches d'une même espèce : une mouche A à yeux rouges et une mouche B à yeux blancs (caractère héréditaire). Propose une expérience permettant de tester l'affirmation de cet élève, et précise les résultats attendus selon qu'elle est vraie ou fausse.
\tcblower
\textbf{Solution détaillée.}
\begin{enumerate}
\item \textbf{Problème} : l'information déterminant la couleur des yeux est-elle dans le cytoplasme ou dans le noyau ?
\item \textbf{Hypothèses} : H1, l'information est dans le noyau ; H2, l'information est dans le cytoplasme (affirmation de l'élève).
\item \textbf{Expérience} : on prélève le noyau d'une cellule de la mouche A (yeux rouges) et on l'introduit dans une cellule de la mouche B (yeux blancs) préalablement énucléée. On suit le développement de cette cellule recombinée.
\item \textbf{Résultats attendus} :
\begin{itemize}
\item si H2 est vraie (l'élève a raison) : le caractère obtenu sera « yeux blancs », celui du donneur de cytoplasme B ;
\item si H1 est vraie : le caractère obtenu sera « yeux rouges », celui du donneur de noyau A.
\end{itemize}
\item \textbf{Conclusion attendue} : le résultat réel de l'expérience permettra de valider l'une des deux hypothèses et d'invalider l'autre. (Nous verrons dans la suite de la leçon que c'est H1 qui se vérifie.)
\end{enumerate}
\end{exoresolu}

\begin{savaistu}[Et au Cameroun ?]
Les groupes sanguins A, B, AB et O sont des caractères héréditaires bien connus des laboratoires d'analyses médicales du Cameroun. Avant une transfusion sanguine à l'hôpital, on détermine toujours le groupe sanguin du malade : ce caractère, reçu de ses parents, ne change jamais au cours de la vie. C'est une preuve quotidienne que nos cellules transportent une information héréditaire stable !
\end{savaistu}

\begin{aretenir}[L'essentiel de la section]
\begin{itemize}
\item Les \textbf{caractères héréditaires} (groupe sanguin, couleur des yeux, forme du lobe de l'oreille...) se transmettent des parents aux enfants : il existe donc une \textbf{information génétique} transmise lors de la reproduction.
\item La cellule comprend trois constituants : la \textbf{membrane}, le \textbf{cytoplasme} et le \textbf{noyau}.
\item \textbf{Problème} : dans quelle partie de la cellule se trouve l'information génétique ?
\item Trois hypothèses vérifiables : le noyau, le cytoplasme ou la membrane.
\item Pour trancher, on réalise des expériences de \textbf{transfert de noyau} : on échange le noyau entre deux cellules et on observe quel caractère se manifeste.
\item Une hypothèse n'est jamais une conclusion : seuls les \textbf{résultats expérimentaux} permettent de la valider ou de l'invalider.
\end{itemize}
\end{aretenir}

Nous avons posé le problème et construit notre plan d'attaque. Il est temps maintenant de passer au laboratoire : dans la prochaine section, nous étudierons les véritables expériences de transfert de noyau, réalisées sur des algues et des grenouilles, qui ont permis de localiser définitivement l'information génétique.

\coursec{Localisation de l'information génétique : les expériences de transfert de noyau}

Dans la section précédente, nous avons posé le problème : \og~où se trouve l'information génétique qui détermine nos caractères héréditaires ?~\fg{} Nous avons formulé deux hypothèses principales : cette information se trouve soit dans le \cle{noyau}, soit dans le \cle{cytoplasme} de la cellule. Pour trancher, la science ne se contente pas de spéculer : elle conçoit des \cle{expériences de transfert de noyau}. Le raisonnement est simple : si l'on déplace le noyau d'une cellule dans une autre cellule et que les caractères \emph{suivent le noyau}, alors l'information est dans le noyau ; si les caractères restent attachés au cytoplasme receveur, alors l'information est cytoplasmique. C'est exactement la démarche qu'ont menée Joachim Hämmerling dans les années 1930, puis John Gurdon dans les années 1960. Étudions leurs travaux fondateurs.

\courssub{L'expérience de Hämmerling sur l'algue Acetabularia : le noyau commande la forme}

\begin{experience}[Expérience de Hämmerling (années 1930) : greffes sur l'algue \emph{Acetabularia}]
\textbf{Matériel biologique :} l'algue unicellulaire géante \emph{Acetabularia} (surnommée \og~parapluie de mer~\fg{}), qui peut mesurer plusieurs centimètres et est visible à l'œil nu. Son corps (thalle) comporte trois parties : un \cle{rhizoïde} (pied fixateur qui contient le \textbf{seul noyau} de la cellule), un \cle{pédicelle} (tige allongée, sans noyau) et un \cle{chapeau} (partie terminale dont la forme est caractéristique de l'espèce).

\textbf{Espèces utilisées :}
\begin{itemize}
    \item \emph{Acetabularia acetabulum} : chapeau en forme de \textbf{disque lisse} (comme une assiette) ;
    \item \emph{Acetabularia crenulata} : chapeau \textbf{dentelé, lobé} (comme une fleur découpée).
\end{itemize}

\textbf{Protocole :}
\begin{enumerate}
    \item On sectionne le chapeau et le pédicelle des deux espèces, en conservant les pieds (qui contiennent les noyaux).
    \item On réalise des \textbf{greffes croisées} : on fixe le pied (rhizoïde + noyau) d'\emph{A. crenulata} sur le pédicelle (cytoplasme) d'\emph{A. acetabulum}, et réciproquement.
    \item On laisse les algues greffées régénérer un nouveau chapeau.
\end{enumerate}

\textbf{Observations et résultats :}
\begin{itemize}
    \item Le pédicelle d'\emph{A. acetabulum} (cytoplasme \og~lisse~\fg{}) portant le pied d'\emph{A. crenulata} (noyau \og~dentelé~\fg{}) régénère un chapeau \textbf{dentelé}, caractéristique d'\emph{A. crenulata}.
    \item Réciproquement, le pied d'\emph{A. acetabulum} (noyau \og~lisse~\fg{}) greffé sur le pédicelle d'\emph{A. crenulata} produit un chapeau \textbf{lisse}.
    \item \textbf{Témoin :} un fragment d'algue \textbf{privé de noyau} ne régénère pas de chapeau normal et finit par dégénérer.
\end{itemize}

\textbf{Interprétation :} le caractère \og~forme du chapeau~\fg{} suit toujours l'origine du \textbf{noyau} (contenu dans le pied), et jamais l'origine du cytoplasme (contenu dans le pédicelle). L'information qui détermine la forme du chapeau est donc localisée dans le noyau. De plus, la présence du noyau est \textbf{indispensable} à la régénération.
\end{experience}

\begin{popfigure}
\begin{tikzpicture}[scale=0.8, every node/.style={font=\small}]
    \colorlet{capA}{popblue!70}
    \colorlet{capB}{poporange!80}
    \colorlet{piedA}{popblue!40}
    \colorlet{piedB}{poporange!40}
    \colorlet{tige}{popgreen!60!black}
    \colorlet{noyau}{poppurple}

    % Étape 1
    \node[align=center, font=\bfseries] at (0, 5.6) {Étape 1 : deux espèces aux chapeaux différents};
    % A. acetabulum
    \draw[capA, line width=1.5pt, fill=capA!30] (-5.6, 4.6) ellipse (1 and 0.3);
    \draw[tige, line width=2.5pt] (-5.6, 4.3) -- (-5.6, 1.2);
    \draw[piedA, line width=2pt, fill=piedA!50] (-5.6, 1.2) -- (-5.9, 0.4) -- (-5.3, 0.4) -- cycle;
    \fill[noyau] (-5.6, 0.7) circle (0.13);
    \node[align=center] at (-5.6, -0.4) {\emph{A. acetabulum}\\chapeau lisse};
    % A. crenulata
    \draw[capB, line width=1.5pt, fill=capB!30] (5.6, 4.6) ellipse (1 and 0.3);
    \foreach \x in {-0.8,-0.4,0,0.4,0.8} {\draw[capB, line width=1pt] (5.6,4.6) -- (5.6+\x,5.1);}
    \draw[tige, line width=2.5pt] (5.6, 4.3) -- (5.6, 1.2);
    \draw[piedB, line width=2pt, fill=piedB!50] (5.6, 1.2) -- (5.3, 0.4) -- (5.9, 0.4) -- cycle;
    \fill[noyau] (5.6, 0.7) circle (0.13);
    \node[align=center] at (5.6, -0.4) {\emph{A. crenulata}\\chapeau dentelé};

    % Étape 2
    \node[align=center, font=\bfseries] at (0, -1.6) {Étape 2 : greffes croisées (pied d'une espèce sur pédicelle de l'autre)};
    % Greffe 1
    \draw[tige, line width=2.5pt] (-5.6, -2.4) -- (-5.6, -4.4);
    \draw[piedB, line width=2pt, fill=piedB!50] (-5.6, -4.4) -- (-5.9, -5.2) -- (-5.3, -5.2) -- cycle;
    \fill[noyau] (-5.6, -4.9) circle (0.13);
    \node[align=center] at (-5.6, -6) {pied \emph{crenulata}\\sur pédicelle \emph{acetabulum}};
    % Greffe 2
    \draw[tige, line width=2.5pt] (5.6, -2.4) -- (5.6, -4.4);
    \draw[piedA, line width=2pt, fill=piedA!50] (5.6, -4.4) -- (5.3, -5.2) -- (5.9, -5.2) -- cycle;
    \fill[noyau] (5.6, -4.9) circle (0.13);
    \node[align=center] at (5.6, -6) {pied \emph{acetabulum}\\sur pédicelle \emph{crenulata}};

    % Étape 3
    \node[align=center, font=\bfseries] at (0, -7.4) {Étape 3 : régénération du chapeau};
    % Résultat 1
    \draw[capB, line width=1.5pt, fill=capB!30] (-5.6, -8.2) ellipse (1 and 0.3);
    \foreach \x in {-0.8,-0.4,0,0.4,0.8} {\draw[capB, line width=1pt] (-5.6,-8.2) -- (-5.6+\x,-7.7);}
    \draw[tige, line width=2.5pt] (-5.6, -8.5) -- (-5.6, -10.2);
    \draw[piedB, line width=2pt, fill=piedB!50] (-5.6, -10.2) -- (-5.9, -11) -- (-5.3, -11) -- cycle;
    \fill[noyau] (-5.6, -10.7) circle (0.13);
    \node[align=center] at (-5.6, -11.8) {chapeau \textbf{dentelé}\\= caractère du donneur\\du noyau};
    % Résultat 2
    \draw[capA, line width=1.5pt, fill=capA!30] (5.6, -8.2) ellipse (1 and 0.3);
    \draw[tige, line width=2.5pt] (5.6, -8.5) -- (5.6, -10.2);
    \draw[piedA, line width=2pt, fill=piedA!50] (5.6, -10.2) -- (5.3, -11) -- (5.9, -11) -- cycle;
    \fill[noyau] (5.6, -10.7) circle (0.13);
    \node[align=center] at (5.6, -11.8) {chapeau \textbf{lisse}\\= caractère du donneur\\du noyau};

    % Légende
    \fill[noyau] (-1.2, 0.7) circle (0.13);
    \node[right] at (-0.9, 0.7) {noyau (dans le pied)};
\end{tikzpicture}
\legende{Expérience de Hämmerling sur \emph{Acetabularia} : le chapeau régénéré adopte toujours la forme de l'espèce donneuse du noyau (le pied), quelle que soit l'origine du cytoplasme (le pédicelle). L'information qui détermine la forme est donc dans le noyau.}
\end{popfigure}

\begin{propriete}[Localisation de l'information génétique (conclusion de l'expérience de Hämmerling)]
L'information génétique qui détermine les caractères héréditaires (ici, la forme du chapeau) est localisée dans le \textbf{noyau} de la cellule. Le cytoplasme seul ne peut pas assurer le développement normal ni transmettre les caractères propres de l'espèce.
\end{propriete}

\begin{attention}[Piège fréquent : ne pas conclure que le cytoplasme est \og inutile \fg]
Le cytoplasme est indispensable à la vie de la cellule : c'est là que se déroulent la plupart des réactions chimiques qui fournissent l'énergie et fabriquent les matériaux de la cellule. Mais l'\textbf{information} (le \og programme~\fg{}, la \og recette~\fg{} des caractères) se trouve dans le noyau. Dans l'expérience, le cytoplasme receveur \emph{exécute} le programme apporté par le noyau donneur. Dire \og l'information est dans le noyau~\fg{} ne signifie donc pas \og le cytoplasme ne sert à rien~\fg{}.
\end{attention}

\begin{methode}[Interpréter une expérience de transfert de noyau]
Pour conclure sur la localisation de l'information génétique à partir d'une greffe ou d'un transfert de noyau :
\begin{enumerate}
    \item Identifier le \textbf{donneur du noyau} et le \textbf{donneur du cytoplasme} (receveur).
    \item Observer les \textbf{caractères} de l'individu obtenu (forme, couleur, etc.).
    \item \textbf{Comparer} ces caractères à ceux du donneur du noyau et à ceux du donneur du cytoplasme.
    \item \textbf{Conclure} : si les caractères correspondent au donneur du noyau, l'information génétique est dans le noyau ; s'ils correspondent au donneur du cytoplasme, elle est dans le cytoplasme.
\end{enumerate}
\end{methode}

\courssub{L'expérience de Gurdon sur la grenouille : le noyau d'une cellule spécialisée contient toute l'information}

L'expérience de Hämmerling montre que le noyau dirige le développement chez une algue. Mais chez les animaux, les cellules se \cle{spécialisent} au cours du développement : cellules de la peau, de l'intestin, des muscles\ldots{} Une question se pose alors : une cellule spécialisée a-t-elle \textbf{perdu} une partie de son information génétique en se spécialisant ? John Gurdon (prix Nobel de médecine 2012) a répondu par une expérience de \cle{transfert de noyau} chez la grenouille africaine \emph{Xenopus laevis}.

\begin{experience}[Expérience de Gurdon (1962) : transfert de noyau chez la grenouille \emph{Xenopus}]
\textbf{Objectif :} vérifier si le noyau d'une cellule spécialisée (cellule intestinale d'un têtard) contient encore \textbf{toute} l'information nécessaire pour fabriquer un individu complet.

\textbf{Protocole :}
\begin{enumerate}
    \item \textbf{Préparation du receveur :} on prélève un \textbf{ovule} (cellule reproductrice femelle) de grenouille et on le détruit de son noyau : on dit qu'on l'\textbf{énuclée} (par irradiation aux ultraviolets ou à l'aide d'une micropipette). L'ovule énucléé conserve son cytoplasme, riche en réserves nutritives.
    \item \textbf{Préparation du donneur :} on prélève une \textbf{cellule de l'intestin} d'un \textbf{têtard} (cellule déjà spécialisée) et on en isole le noyau.
    \item \textbf{Transfert :} on introduit le noyau de la cellule intestinale dans l'ovule énucléé.
    \item \textbf{Développement :} on laisse l'ovule \og reconstitué~\fg{} se développer dans des conditions convenables.
\end{enumerate}

\textbf{Résultats :}
\begin{itemize}
    \item Certains ovules reconstitués se divisent puis se développent normalement et donnent des \textbf{têtards} complets et viables, qui deviennent des grenouilles adultes.
    \item Ces têtards possèdent les \textbf{caractères du têtard donneur du noyau} (et non ceux de la grenouille donneuse de l'ovule) : ce sont des \textbf{clones} du donneur du noyau.
    \item Le taux de réussite est faible, mais le résultat est décisif.
\end{itemize}

\textbf{Interprétation :} la spécialisation d'une cellule ne s'accompagne \textbf{pas d'une perte d'information génétique} : le noyau de la cellule intestinale contient encore \textbf{toute l'information} nécessaire à la construction d'un organisme entier. Le cytoplasme de l'ovule fournit les conditions (réserves, matériel cellulaire) permettant à ce programme nucléaire de s'exprimer.
\end{experience}

\begin{popfigure}
\begin{tikzpicture}[scale=0.9, every node/.style={font=\small}]
    \colorlet{ovule}{poppink!50}
    \colorlet{noyauT}{popblue!80}
    \colorlet{tadpole}{popgreen!60!black}
    \colorlet{cellInt}{poporange!70}

    % 1. Donneur
    \node[align=center, font=\bfseries] at (-5, 3.6) {1. Donneur du noyau};
    \draw[tadpole, line width=1.5pt, fill=tadpole!20] (-5, 2.6) ellipse (1.1 and 0.55);
    \draw[tadpole, line width=1.5pt, fill=tadpole!20] (-6.3, 2.6) circle (0.45);
    \draw[tadpole, line width=1pt] (-4, 2.6) -- (-3.2, 2.9) -- (-3.6, 2.3);
    \node[align=center] at (-5, 1.6) {têtard};
    \draw[cellInt, line width=1.5pt, fill=cellInt!20] (-5.7, 0.9) rectangle (-4.3, 0.1);
    \fill[noyauT] (-5, 0.5) circle (0.13);
    \node[align=center] at (-5, -0.5) {cellule intestinale\\(spécialisée)};

    % Flèche prélèvement
    \draw[->, very thick, popdark] (-4.3, 0.5) -- (-2.2, 0.5);
    \node[above, align=center] at (-3.2, 0.6) {prélèvement\\du noyau};

    % 2. Receveur
    \node[align=center, font=\bfseries] at (0.8, 3.6) {2. Receveur : ovule énucléé};
    \draw[ovule, line width=1.5pt, fill=ovule!25] (0.8, 1.8) circle (1.1);
    \node[align=center] at (0.8, 1.8) {ovule\\sans noyau\\(cytoplasme\\riche en réserves)};
    \draw[<-, thick, popdark] (1.5, 2.7) -- (2.6, 3.1);
    \node[right, align=left] at (2.7, 3.1) {énucléation};

    % Flèche transfert
    \draw[->, very thick, poppurple] (-2.2, 0.5) -- (-0.6, 0.5) -- (-0.1, 1.1);
    \node[below, align=center, text=poppurple] at (-1.4, 0.3) {\textbf{transfert du noyau}};

    % 3. Résultat
    \node[align=center, font=\bfseries] at (5.6, 3.6) {3. Résultat};
    \draw[ovule, line width=1.5pt, fill=ovule!25] (5.6, 1.8) circle (1.1);
    \fill[noyauT] (5.6, 1.8) circle (0.15);
    \node[align=center] at (5.6, 0.3) {ovule reconstitué\\(noyau du têtard donneur)};
    \draw[->, very thick, popdark] (5.6, -0.4) -- (5.6, -1.2);
    \node[align=center] at (5.6, -1.7) {développement};
    \draw[tadpole, line width=1.5pt, fill=tadpole!20] (5.6, -3) ellipse (1.1 and 0.55);
    \draw[tadpole, line width=1.5pt, fill=tadpole!20] (4.3, -3) circle (0.45);
    \draw[tadpole, line width=1pt] (6.6, -3) -- (7.4, -2.7) -- (7, -3.3);
    \node[align=center] at (5.6, -4.1) {têtard \textbf{clone} : mêmes caractères\\que le donneur du noyau};

    % Légende
    \fill[noyauT] (-6.8, -1.7) circle (0.13);
    \node[right] at (-6.5, -1.7) {noyau du donneur (cellule intestinale)};
\end{tikzpicture}
\legende{Expérience de Gurdon : le transfert du noyau d'une cellule intestinale de têtard dans un ovule énucléé permet le développement d'un têtard complet, clone du donneur du noyau. Le noyau d'une cellule spécialisée contient donc toute l'information génétique de l'espèce.}
\end{popfigure}

\begin{definition}[Énucléer ; clone]
\cle{Énucléer} une cellule, c'est \textbf{retirer ou détruire son noyau} (à l'aide d'une micropipette ou par irradiation aux ultraviolets). La cellule énucléée conserve son cytoplasme mais perd son information génétique : elle ne peut plus se diviser ni se développer.\par
Un \cle{clone} est un individu (ou une cellule) qui possède \textbf{exactement la même information génétique} qu'un autre individu : il en est la copie génétique.
\end{definition}

\begin{savaistu}[De la grenouille à la brebis Dolly : l'aventure du clonage]
L'expérience de Gurdon (1962) a démontré le principe du clonage : le noyau d'une cellule spécialisée, placé dans le cytoplasme d'un ovule énucléé, peut commander le développement d'un individu complet. En \textbf{1996}, l'équipe d'Ian Wilmut en Écosse réussit le premier clonage d'un \textbf{mammifère} à partir d'une cellule adulte : la brebis \textbf{Dolly}, née le 5 juillet 1996. Le noyau provenait d'une cellule de la \textbf{glande mammaire} d'une brebis adulte ; il a été introduit dans un ovule énucléé d'une autre brebis, puis l'embryon a été implanté dans l'utérus d'une mère porteuse. Dolly était le clone de la brebis donneuse du noyau. Ces recherches ouvrent des perspectives médicales importantes, mais le clonage reproductif humain est interdit dans le monde entier pour des raisons éthiques.
\end{savaistu}

\courssub{Synthèse comparative : les résultats tranchent entre les deux hypothèses}

Pour ancrer le raisonnement scientifique, comparons les prévisions des deux hypothèses de départ avec les résultats réellement observés.

\begin{center}
\begin{tabularx}{\linewidth}{|l|X|X|X|}
\hline
\rowcolor{popblueL}
\textbf{Expérience} & \textbf{Hypothèse 1 : information dans le noyau} & \textbf{Hypothèse 2 : information dans le cytoplasme} & \textbf{Résultat observé (décision)} \\
\hline
\textbf{Hämmerling} (greffes sur \emph{Acetabularia}) & Le nouveau chapeau a la forme de l'espèce donneuse du \textbf{pied (noyau)}. & Le nouveau chapeau a la forme de l'espèce donneuse du \textbf{pédicelle (cytoplasme)}. & Le chapeau a la forme de l'espèce \textbf{donneuse du noyau}. $\rightarrow$ Hypothèse 1 validée. \\
\hline
\textbf{Gurdon} (transfert de noyau chez la grenouille) & L'individu obtenu a les caractères du \textbf{donneur du noyau} (cellule intestinale). & L'individu obtenu a les caractères de la \textbf{donneuse de l'ovule} (cytoplasme). & L'individu a les caractères du \textbf{donneur du noyau} : c'est son clone. $\rightarrow$ Hypothèse 1 validée. \\
\hline
\end{tabularx}
\end{center}

\begin{aretenir}[Savoir établi : localisation de l'information génétique]
Au terme de ces deux expériences historiques, indépendantes et convergentes (l'une sur un végétal unicellulaire, l'autre sur un animal), nous construisons le savoir fondamental :
\begin{itemize}
    \item L'\textbf{information génétique} qui détermine les caractères héréditaires est localisée dans le \textbf{noyau} des cellules.
    \item Le noyau de chaque cellule contient le \textbf{programme complet} de développement de l'organisme, même lorsque la cellule est spécialisée.
    \item Le cytoplasme permet l'\textbf{expression} de cette information (il fournit l'énergie et le matériel cellulaire), mais il ne \textbf{contient} pas l'information des caractères héréditaires.
\end{itemize}
\end{aretenir}

\courssub{Application : la fécondation, rencontre de deux informations génétiques}

\begin{exemplebox}[Pourquoi ressemblons-nous à nos deux parents ?]
La fécondation est l'union de deux cellules reproductrices (gamètes) : le \textbf{spermatozoïde}, qui apporte essentiellement son \textbf{noyau} (il contient très peu de cytoplasme), et l'\textbf{ovule}, qui apporte son \textbf{noyau} et tout son cytoplasme riche en réserves.
\begin{itemize}
    \item La cellule-œuf (zygote) reçoit donc deux lots d'information génétique : l'un apporté par le noyau du père, l'autre par le noyau de la mère.
    \item Puisque l'information génétique se trouve dans le noyau, l'enfant reçoit \textbf{la moitié de son information génétique de son père et la moitié de sa mère}.
    \item C'est la base de l'\textbf{hérédité} : chaque enfant est un mélange unique des caractères de ses deux parents, ce qui explique à la fois les ressemblances et les différences observées au sein d'une même famille.
\end{itemize}
\end{exemplebox}

\begin{exoresolu}[Analyse d'une expérience de transfert de noyau]
\textbf{Énoncé :} On reprend le dispositif de Gurdon chez la grenouille. On prélève le noyau d'une cellule de la \textbf{peau} d'une grenouille à peau \textbf{verte} et on l'introduit dans un ovule \textbf{énucléé} prélevé chez une grenouille à peau \textbf{brune}. L'ovule reconstitué se développe et donne un têtard.
\begin{enumerate}
    \item Qui est le donneur du noyau ? Qui est le donneur du cytoplasme ?
    \item Prévoir la couleur de peau du têtard obtenu. Justifier la réponse.
    \item Ce têtard est-il un clone ? Si oui, de quelle grenouille ?
\end{enumerate}
\tcblower
\textbf{Solution détaillée :}
\begin{enumerate}
    \item Le \textbf{donneur du noyau} est la grenouille à peau \textbf{verte} (on a prélevé le noyau d'une de ses cellules de peau). Le \textbf{donneur du cytoplasme} est la grenouille à peau \textbf{brune} (son ovule a fourni le cytoplasme, mais son noyau a été retiré).
    \item Le têtard aura la peau \textbf{verte}. \textbf{Justification :} les expériences de transfert de noyau (Hämmerling, Gurdon) montrent que les caractères de l'individu obtenu sont ceux du \textbf{donneur du noyau}, car c'est dans le noyau que se trouve l'information génétique. Le cytoplasme de l'ovule (grenouille brune) ne contient pas l'information du caractère \og couleur de la peau~\fg{}.
    \item \textbf{Oui}, ce têtard est un \textbf{clone de la grenouille à peau verte}, donneuse du noyau : il possède exactement la même information génétique qu'elle.
\end{enumerate}
\textbf{Point méthode :} dans toute expérience de transfert de noyau, identifier d'abord \textbf{qui fournit le noyau} : c'est lui qui détermine les caractères héréditaires de l'individu obtenu.
\end{exoresolu}

\begin{exercice}[Application : transfert de noyau chez la souris]
On réalise une expérience de transfert de noyau chez la souris. On introduit le noyau d'une cellule de la peau d'une souris à pelage \textbf{noir} dans un ovule \textbf{énucléé} prélevé chez une souris à pelage \textbf{blanc}. L'embryon obtenu est ensuite placé dans l'utérus d'une troisième souris (mère porteuse) à pelage blanc. Une souris naît.
\begin{enumerate}
    \item Quelle est la couleur du pelage de la souris nouveau-née ? Justifier.
    \item Cette souris est-elle un clone ? Si oui, de quelle souris ?
    \item La mère porteuse a-t-elle transmis de l'information génétique au petit ? Justifier.
\end{enumerate}
\end{exercice}

\begin{corrige}[Exercice : transfert de noyau chez la souris]
\begin{enumerate}
    \item La souris nouveau-née a le pelage \textbf{noir}. En effet, l'information génétique qui détermine la couleur du pelage se trouve dans le noyau, et le noyau provient de la souris noire. Le cytoplasme de l'ovule (souris blanche) ne contient pas l'information de ce caractère.
    \item \textbf{Oui}, c'est un \textbf{clone de la souris noire} donneuse du noyau : elle possède exactement la même information génétique qu'elle.
    \item \textbf{Non.} La mère porteuse a seulement assuré le développement de l'embryon dans son utérus (protection, nutrition). Elle n'a transmis \textbf{aucun noyau}, donc aucune information génétique, au petit.
\end{enumerate}
\end{corrige}

\coursec{Les chromosomes, supports de l'information génétique}

\courssub{De la localisation nucléaire à l'identification des porteurs de l'hérédité}

Dans la section précédente, les expériences de transfert de noyau (notamment celles de \textbf{Gurdon} sur la grenouille \emph{Xenopus} et, plus près de nous, le clonage de la brebis Dolly) ont démontré sans ambiguïté que le \cle{noyau} de la cellule est le siège de l'information génétique. Mais le noyau n'est pas une « boîte noire » homogène : lorsqu'une cellule se prépare à se diviser, son contenu nucléaire se condense et devient visible au microscope optique sous la forme de petits bâtonnets colorés. C'est l'observation de ces structures, constantes pour une espèce donnée, qui a permis d'identifier les \cle{chromosomes} comme les supports physiques de l'hérédité. Nous allons maintenant décrire leur structure, leur organisation en paires et l'outil majeur qui permet de les étudier : le caryotype.

\textbf{Observation au microscope : les chromosomes, corps colorés du noyau en division}\par

\paragraph{Contexte historique et étymologique.} Dès la fin du XIX\textsuperscript{e} siècle, les biologistes (Waldeyer, Boveri, Sutton) observent que, lors de la division cellulaire (mitose), le noyau disparaît et laisse place à des filaments qui s'épaississent, se raccourcissent et se teintent fortement avec les colorants basiques (hématoxyline, bleu de méthylène, Giemsa). Le terme \textbf{chromosome} vient du grec \emph{chroma} (couleur) et \emph{soma} (corps) : ce sont littéralement des « corps colorés ».

\begin{experience}[Observation des chromosomes au microscope optique (racine d'oignon ou pointe de racine de \emph{Vicia faba})]
\textbf{Matériel :} Pointe de racine d'oignon (ou fève) en croissance, acide chlorhydrique \SI{1}{N}, colorant (orceine acétique ou bleu de toluidine), lame, lamelle, microscope optique (objectif $\times 40$ ou $\times 100$ à immersion).
\textbf{Protocole :}
\begin{enumerate}
    \item Faire germer des bulbes d'oignon dans l'eau pendant \SIrange{2}{3}{jours} (racines de \SIrange{1}{2}{cm}).
    \item Prélèver l'extrémité des racines (méristème apical, zone de division active) et les placer dans l'HCl \SI{1}{N} à \SI{60}{\celsius} pendant \SI{1}{min} (ramollissement des parois et arrêt de la division en métaphase).
    \item Rincer à l'eau distillée, puis colorer \SIrange{10}{15}{min} dans l'orceine acétique tiédie.
    \item Écraser la pointe de racine sur une lame dans une goutte de colorant, poser la lamelle, appuyer fermement (écrasement pour étaler les cellules en monocouche).
    \item Observer au microscope.
\end{enumerate}
\textbf{Observations :} On distingue des cellules aux contours nets. Dans la plupart, le noyau est au repos (interphase) : on voit un noyau clair avec un nucléole. Dans d'autres (cellules en \textbf{métaphase}), le noyau a disparu et on observe, au centre de la cellule, un ensemble de \textbf{bâtonnets courts, épais, fortement colorés}, disposés sur un plan équatorial : ce sont les \textbf{chromosomes}.
\textbf{Interprétation :} La condensation de la chromatine (ADN + protéines) rend les chromosomes visibles et dénombrables uniquement pendant la division cellulaire (mitose ou méiose). À l'interphase, ils sont décondensés (chromatine) et invisibles individuellement.
\textbf{Conclusion :} Les chromosomes sont les structures individualisées du matériel génétique lors de la division cellulaire.
\end{experience}

\begin{definition}[Chromosome]
Un \textbf{chromosome} est une structure filamenteuse, située dans le noyau des cellules eucaryotes, devenue visible et individualisée au microscope optique lors de la division cellulaire (mitose ou méiose) par condensation de la chromatine. Il est constitué d'une molécule d'ADN très enroulée autour de protéines (histones) et porte une partie de l'information génétique de l'individu.
\end{definition}

\textbf{Structure d'un chromosome : monochromatidien et bichromatidien}\par

L'aspect d'un chromosome change au cours du cycle cellulaire. Il existe deux états morphologiques fondamentaux qu'il faut savoir reconnaître et schématiser.

\begin{popfigure}
\begin{tikzpicture}[scale=1.2, every node/.style={font=\small}]
    % Chromosome monochromatidien (G1 / après mitose)
    \begin{scope}[xshift=-5cm]
        \draw[popdark, thick, fill=popblueL!30, rounded corners=2pt] (-0.6,-1.5) -- (-0.6,1.5) -- (0.6,1.5) -- (0.6,-1.5) -- cycle;
        \draw[popdark, thick] (-0.6,0) -- (0.6,0) node[midway, above, xshift=0.8cm] {\textbf{Centromère}};
        \node[popdark, align=center] at (0, -2.2) {\textbf{Chromosome monochromatidien}\\ (1 chromatide)};
        \node[popdark, align=center] at (0, 2.2) {État G1 / après mitose};
        % Bras
        \node[popdark] at (-1.2, 0.8) {Bras court (p)};
        \node[popdark] at (-1.2, -0.8) {Bras long (q)};
        \draw[->, popdark] (-1.0, 0.6) -- (-0.65, 0.8);
        \draw[->, popdark] (-1.0, -0.6) -- (-0.65, -0.8);
    \end{scope}

    % Flèche cycle
    \draw[poporange, very thick, ->] (0.5, 0) -- node[above] {Réplication ADN (Phase S)} node[below] {Interphase} (2.5, 0);

    % Chromosome bichromatidien (G2 / Métaphase)
    \begin{scope}[xshift=5.5cm]
        % Chromatide sœur 1 (gauche)
        \draw[popdark, thick, fill=popgreenL!30, rounded corners=2pt] (-0.6,-1.5) -- (-0.6,1.5) -- (0,1.5) -- (0,-1.5) -- cycle;
        % Chromatide sœur 2 (droite)
        \draw[popdark, thick, fill=poporangeL!30, rounded corners=2pt] (0,-1.5) -- (0,1.5) -- (0.6,1.5) -- (0.6,-1.5) -- cycle;
        % Centromère commun
        \draw[poppurple, very thick] (0, -1.5) -- (0, 1.5) node[midway, right, xshift=0.5cm] {\textbf{Centromère (kinétochore)}};
        \node[popdark, align=center] at (0, -2.2) {\textbf{Chromosome bichromatidien}\\ (2 chromatides sœurs)};
        \node[popdark, align=center] at (0, 2.2) {État G2 / Métaphase};
        % Annotation chromatides
        \node[popgreen!80!black] at (-1.5, 0) {\textbf{Chromatide sœur 1}};
        \node[poporange!80!black] at (1.5, 0) {\textbf{Chromatide sœur 2}};
        \draw[->, popgreen!80!black] (-1.1, 0.2) -- (-0.5, 0.5);
        \draw[->, poporange!80!black] (1.1, 0.2) -- (0.5, 0.5);
    \end{scope}
\end{tikzpicture}
\legende{Schéma comparatif d'un chromosome monochromatidien (une chromatide, après division) et bichromatidien (deux chromatides sœurs issues de la réplication de l'ADN, avant division). Le centromère unit les deux chromatides et délimite le bras court (p) et le bras long (q).}
\end{popfigure}

\begin{propriete}[États structuraux du chromosome]
\begin{itemize}
    \item \textbf{Chromosome monochromatidien :} Composé d'une \textbf{seule chromatide} (une seule molécule d'ADN double brin). C'est l'état du chromosome en phase G1 (après la mitose) et dans les gamètes (état haploïde $n$).
    \item \textbf{Chromosome bichromatidien :} Composé de \textbf{deux chromatides sœurs} (deux molécules d'ADN double brin identiques, issues de la réplication de l'ADN lors de la phase S de l'interphase). Elles sont jointes au niveau du \textbf{centromère}. C'est l'état visible en métaphase de la mitose et en métaphase II de la méiose.
    \item \textbf{Centromère :} Région constrictée où les deux chromatides sœurs sont unies. C'est le point d'attache des fibres du fuseau mitotique (via le kinétochore). Il détermine la position des bras : bras court (\emph{p} pour \emph{petit}) et bras long (\emph{q} pour \emph{queue}).
\end{itemize}
\end{propriete}

\begin{attention}[Piège fréquent : confusion chromatide / chromosome]
\textbf{Ne pas confondre le nombre de chromosomes et le nombre de chromatides.}
\begin{itemize}
    \item On compte le nombre de \textbf{centromères} pour connaître le nombre de chromosomes.
    \item Un chromosome bichromatidien = \textbf{1 chromosome} (1 centromère) mais \textbf{2 chromatides} (2 molécules d'ADN).
    \item Exemple : En métaphase de mitose, une cellule humaine a 46 chromosomes (46 centromères) mais 92 chromatides (92 molécules d'ADN).
\end{itemize}
\end{attention}

\textbf{Le nombre de chromosomes : une constante d'espèce}\par

\paragraph{Dénombrement.} L'observation de nombreuses métaphases chez un même individu, et la comparaison entre individus d'une même espèce, révèle une règle fondamentale : \textbf{le nombre de chromosomes est constant et caractéristique de l'espèce}.

\begin{center}
\begin{tabularx}{\linewidth}{|l|X|c|}\hline
\rowcolor{popblueL}
\textbf{Espèce} & \textbf{Nom scientifique} & \textbf{Nombre de chromosomes (2n)} \\ \hline
Homme & \emph{Homo sapiens} & 46 \\ \hline
Chimpanzé & \emph{Pan troglodytes} & 48 \\ \hline
Gorille & \emph{Gorilla gorilla} & 48 \\ \hline
Chien & \emph{Canis lupus familiaris} & 78 \\ \hline
Chat & \emph{Felis catus} & 38 \\ \hline
Souris & \emph{Mus musculus} & 40 \\ \hline
Maïs & \emph{Zea mays} & 20 \\ \hline
Blé tendre & \emph{Triticum aestivum} & 42 \\ \hline
Mouche du vinaigre & \emph{Drosophila melanogaster} & 8 \\ \hline
\end{tabularx}
\end{center}

\begin{aretenir}[Nombre chromosomique]
\begin{itemize}
    \item \textbf{2n (diploïde) :} Nombre de chromosomes dans les cellules somatiques (non sexuées). Il est \textbf{pair} car les chromosomes sont organisés en paires. Pour l'Homme : \textbf{2n = 46}.
    \item \textbf{n (haploïde) :} Nombre de chromosomes dans les gamètes (spermatozoïdes, ovocytes). Il correspond à une seule série de chromosomes. Pour l'Homme : \textbf{n = 23}.
    \item \textbf{Règle :} $2n = 2 \times n$. La fécondation (fusion de deux gamètes $n$) restaure le nombre diploïde $2n$ chez le zygote.
\end{itemize}
\end{aretenir}

\begin{attention}[Nombre de chromosomes $\neq$ Complexité de l'organisme]
Le nombre de chromosomes \textbf{ne dépend ni de la taille, ni de la complexité apparente, ni du « rang » évolutif} de l'espèce.
\begin{itemize}
    \item L'homme (46) en a moins que le chien (78) ou le blé (42).
    \item La mouche \emph{Drosophila} n'en a que 8.
    \item Le record animal est détenu par un papillon (\emph{Atlas}) avec 2n $\approx$ 448-450 ; côté plantes, la fougère \emph{Ophioglossum} atteint 2n = 1260 !
\end{itemize}
C'est la \textbf{nature de l'information} (les gènes portés), et non le nombre de « boîtes » (chromosomes) qui détermine l'organisme.
\end{attention}

\textbf{Les chromosomes homologues : l'origine biparentale de l'hérédité}\par

Chez les organismes diploïdes (comme l'Homme), les chromosomes ne sont pas tous différents deux à deux. Ils s'apparentent par \textbf{paires}.

\begin{definition}[Chromosomes homologues]
Deux \textbf{chromosomes homologues} (ou \textbf{homologues}) sont deux chromosomes d'une même paire qui :
\begin{enumerate}
    \item Ont la \textbf{même forme} (même taille, même position du centromère, même bandeau de coloration/G-banding).
    \item Portent les \textbf{mêmes gènes} aux \textbf{mêmes loci} (emplacements), c'est-à-dire qu'ils déterminent les \textbf{mêmes caractères} (ex : groupe sanguin, couleur des yeux, maladie génétique...).
    \item Ont une \textbf{origine parentale différente} : l'un vient du \textbf{père} (chromosome paternel), l'autre de la \textbf{mère} (chromosome maternel).
\end{enumerate}
\end{definition}

\begin{exemplebox}[Allèles sur chromosomes homologues]
Prenons le gène du groupe sanguin ABO (localisé sur le chromosome 9).
\begin{itemize}
    \item Le chromosome 9 paternel porte peut-être l'allèle $I^A$ (groupe A).
    \item Le chromosome 9 maternel porte peut-être l'allèle $i$ (groupe O).
    \item L'individu est de génotype $[I^A // i]$ (phénotype groupe A).
    \item Les deux chromosomes 9 sont \textbf{homologues} : même gène ABO au même locus, mais \textbf{allèles différents}.
\end{itemize}
\end{exemplebox}

\paragraph{La paire 23 : les chromosomes sexuels (gonosomes).}
Sur les 23 paires humaines, 22 paires sont des \textbf{autosomes} (identiques chez l'homme et la femme). La 23\textsuperscript{e} paire détermine le \textbf{sexe chromosomique} :
\begin{itemize}
    \item \textbf{Femme (XX) :} Deux chromosomes X homologues (même taille, mêmes gènes).
    \item \textbf{Homme (XY) :} Un chromosome X et un chromosome Y. Ils sont \textbf{non homologues} (hétérologues) : le Y est beaucoup plus petit, porte des gènes différents (dont \emph{SRY}, déterminant la masculinité), mais ils s'apparient tout de même lors de la méiose sur une petite région pseudo-autosomique.
\end{itemize}

\textbf{Le caryotype : la « carte d'identité » chromosomique}\par

Pour analyser l'ensemble des chromosomes, les biologistes réalisent un \textbf{caryotype}.

\begin{methode}[Réalisation et analyse d'un caryotype (méthode standard au laboratoire de cytogénétique, ex: Centre Pasteur du Cameroun, Hôpital Central de Yaoundé)]
\begin{enumerate}
    \item \textbf{Culture cellulaire :} Prélèvement de sang (lymphocytes T stimulés par PHA), de liquide amniotique (amniocentèse) ou de chorion (biopsie trophoblastique). Culture \SIrange{48}{72}{h}.
    \item \textbf{Blocage en métaphase :} Ajout de \textbf{colchicine} (inhibiteur de la polymérisation des microtubules du fuseau) $\rightarrow$ accumulation des cellules en métaphase (chromosomes maximalement condensés).
    \item \textbf{Traitement hypotonique :} Gonflement des cellules (éclatement de la membrane plasmique, étalement des chromosomes).
    \item \textbf{Fixation :} Méthanol / Acide acétique (3/1).
    \item \textbf{Étalement sur lame :} Goutte de suspension cellulaire sur lame froide/humide $\rightarrow$ éclatement du noyau, chromosomes dispersés.
    \item \textbf{Coloration par bandage (G-banding au trypsin-Giemsa) :} Révèle des bandes claires et sombres spécifiques pour chaque chromosome (empreinte digitale).
    \item \textbf{Photographie / Numérisation :} Microscope $\times 1000$, caméra CCD.
    \item \textbf{Classement (karyotyping) :} Logiciel de tri automatique ou manuel : appariement des homologues par taille, position du centromère, profil de bandes. Classement par groupes (A à G) puis par numéro décroissant de taille. Paire 23 à part.
\end{enumerate}
\end{methode}

\begin{popfigure}
\begin{tikzpicture}[scale=0.45, every node/.style={font=\tiny}]
    % Paramètres géométriques
    \def\chrW{0.3}
    \def\gapX{0.55}
    \def\pairGapX{0.8}
    \def\rowGapY{2.0}
    
    % Style chromosome
    \tikzset{
        chr/.style={thick, fill=#1!20, rounded corners=0.5pt},
        centro/.style={poppurple, very thick},
        sat/.style={thick}
    }

    % --- GROUPE A (1-3) : Ligne 1 (y=5) ---
    \foreach \n/\h/\c/\col in {1/1.6/0.4/popblue, 2/1.5/0.4/popblue, 3/1.4/0.4/popblue} {
        \pgfmathsetmacro{\xi}{(\n-1)*\gapX}
        \pgfmathsetmacro{\cy}{5 + \h*\c}
        % Paire
        \draw[\col, chr=\col] (\xi, 5) rectangle (\xi+\chrW, 5+\h);
        \draw[centro] (\xi, \cy) -- (\xi+\chrW, \cy);
        \draw[\col, chr=\col] (\xi+\chrW+0.05, 5) rectangle (\xi+2*\chrW+0.05, 5+\h);
        \draw[centro] (\xi+\chrW+0.05, \cy) -- (\xi+2*\chrW+0.05, \cy);
        \node[\col] at (\xi+\chrW/2, 4.6) {\n};
        \node[\col] at (\xi+1.5*\chrW+0.025, 4.6) {\n};
    }
    \node[popblue, anchor=east] at (-0.5, 5.8) {Groupe A (1-3):};

    % --- GROUPE B (4-5) : Ligne 1 (y=5) suite ---
    \foreach \n/\h/\c/\col [count=\i from 0] in {4/1.3/0.3/popgreen, 5/1.25/0.3/popgreen} {
        \pgfmathsetmacro{\xi}{3.5 + \i*\gapX}
        \pgfmathsetmacro{\cy}{5 + \h*\c}
        \draw[\col, chr=\col] (\xi, 5) rectangle (\xi+\chrW, 5+\h);
        \draw[centro] (\xi, \cy) -- (\xi+\chrW, \cy);
        \draw[\col, chr=\col] (\xi+\chrW+0.05, 5) rectangle (\xi+2*\chrW+0.05, 5+\h);
        \draw[centro] (\xi+\chrW+0.05, \cy) -- (\xi+2*\chrW+0.05, \cy);
        \node[\col] at (\xi+\chrW/2, 4.6) {\n};
        \node[\col] at (\xi+1.5*\chrW+0.025, 4.6) {\n};
    }
    \node[popgreen, anchor=east] at (3.0, 5.8) {Groupe B (4-5):};

    % --- GROUPE C (6-12) : Lignes 2-3 (y=3, y=1) ---
    \foreach \n/\h/\c/\col [count=\i] in {6/1.1/0.35/poporange, 7/1.05/0.35/poporange, 8/1.0/0.35/poporange, 9/0.95/0.35/poporange, 10/0.9/0.35/poporange, 11/0.85/0.35/poporange, 12/0.8/0.35/poporange} {
        \pgfmathsetmacro{\row}{int((\i-1)/4)} 
        \pgfmathsetmacro{\ccol}{mod(\i-1,4)}
        \pgfmathsetmacro{\xi}{\ccol*\gapX}
        \pgfmathsetmacro{\yi}{3 - \row*\rowGapY}
        \pgfmathsetmacro{\cy}{\yi + \h*\c}
        \draw[\col, chr=\col] (\xi, \yi) rectangle (\xi+\chrW, \yi+\h);
        \draw[centro] (\xi, \cy) -- (\xi+\chrW, \cy);
        \draw[\col, chr=\col] (\xi+\chrW+0.05, \yi) rectangle (\xi+2*\chrW+0.05, \yi+\h);
        \draw[centro] (\xi+\chrW+0.05, \cy) -- (\xi+2*\chrW+0.05, \cy);
        \node[\col] at (\xi+\chrW/2, \yi-0.4) {\n};
        \node[\col] at (\xi+1.5*\chrW+0.025, \yi-0.4) {\n};
    }
    \node[poporange, anchor=east] at (-0.5, 3.6) {Groupe C (6-12):};

    % --- GROUPE D (13-15) Acrocentriques : Ligne 1 (y=5) droite ---
    \foreach \n/\h/\c/\col [count=\i from 0] in {13/0.75/0.15/poppink, 14/0.7/0.15/poppink, 15/0.65/0.15/poppink} {
        \pgfmathsetmacro{\xi}{5.5 + \i*\gapX}
        \pgfmathsetmacro{\cy}{5 + \h*\c}
        \pgfmathsetmacro{\satY}{5 + \h}
        \draw[\col, chr=\col] (\xi, 5) rectangle (\xi+\chrW, 5+\h);
        \draw[centro] (\xi, \cy) -- (\xi+\chrW, \cy);
        \draw[\col, sat] (\xi, \satY) -- (\xi+\chrW, \satY); % Satellite
        \draw[\col, chr=\col] (\xi+\chrW+0.05, 5) rectangle (\xi+2*\chrW+0.05, 5+\h);
        \draw[centro] (\xi+\chrW+0.05, \cy) -- (\xi+2*\chrW+0.05, \cy);
        \draw[\col, sat] (\xi+\chrW+0.05, \satY) -- (\xi+2*\chrW+0.05, \satY); % Satellite
        \node[\col] at (\xi+\chrW/2, 4.6) {\n};
        \node[\col] at (\xi+1.5*\chrW+0.025, 4.6) {\n};
    }
    \node[poppink, anchor=east] at (5.0, 5.8) {Groupe D (13-15):};

    % --- GROUPE E (16-18) : Ligne 2 (y=3) droite ---
    \foreach \n/\h/\c/\col [count=\i from 0] in {16/0.6/0.45/popgold, 17/0.55/0.45/popgold, 18/0.5/0.45/popgold} {
        \pgfmathsetmacro{\xi}{5.5 + \i*\gapX}
        \pgfmathsetmacro{\cy}{3 + \h*\c}
        \draw[\col, chr=\col] (\xi, 3) rectangle (\xi+\chrW, 3+\h);
        \draw[centro] (\xi, \cy) -- (\xi+\chrW, \cy);
        \draw[\col, chr=\col] (\xi+\chrW+0.05, 3) rectangle (\xi+2*\chrW+0.05, 3+\h);
        \draw[centro] (\xi+\chrW+0.05, \cy) -- (\xi+2*\chrW+0.05, \cy);
        \node[\col] at (\xi+\chrW/2, 2.6) {\n};
        \node[\col] at (\xi+1.5*\chrW+0.025, 2.6) {\n};
    }
    \node[popgold, anchor=east] at (5.0, 3.6) {Groupe E (16-18):};

    % --- GROUPE F (19-20) : Ligne 3 (y=1) droite ---
    \foreach \n/\h/\c/\col [count=\i from 0] in {19/0.45/0.5/poppurple, 20/0.4/0.5/poppurple} {
        \pgfmathsetmacro{\xi}{7.5 + \i*\gapX}
        \pgfmathsetmacro{\cy}{1 + \h*\c}
        \draw[\col, chr=\col] (\xi, 1) rectangle (\xi+\chrW, 1+\h);
        \draw[centro] (\xi, \cy) -- (\xi+\chrW, \cy);
        \draw[\col, chr=\col] (\xi+\chrW+0.05, 1) rectangle (\xi+2*\chrW+0.05, 1+\h);
        \draw[centro] (\xi+\chrW+0.05, \cy) -- (\xi+2*\chrW+0.05, \cy);
        \node[\col] at (\xi+\chrW/2, 0.6) {\n};
        \node[\col] at (\xi+1.5*\chrW+0.025, 0.6) {\n};
    }
    \node[poppurple, anchor=east] at (7.0, 1.6) {Groupe F (19-20):};

    % --- GROUPE G (21-22) Acrocentriques : Ligne 3 (y=1) suite ---
    \foreach \n/\h/\c/\col [count=\i from 0] in {21/0.35/0.1/popmuted, 22/0.3/0.1/popmuted} {
        \pgfmathsetmacro{\xi}{9.0 + \i*\gapX}
        \pgfmathsetmacro{\cy}{1 + \h*\c}
        \pgfmathsetmacro{\satY}{1 + \h}
        \draw[\col, chr=\col] (\xi, 1) rectangle (\xi+\chrW, 1+\h);
        \draw[centro] (\xi, \cy) -- (\xi+\chrW, \cy);
        \draw[\col, sat] (\xi, \satY) -- (\xi+\chrW, \satY);
        \draw[\col, chr=\col] (\xi+\chrW+0.05, 1) rectangle (\xi+2*\chrW+0.05, 1+\h);
        \draw[centro] (\xi+\chrW+0.05, \cy) -- (\xi+2*\chrW+0.05, \cy);
        \draw[\col, sat] (\xi+\chrW+0.05, \satY) -- (\xi+2*\chrW+0.05, \satY);
        \node[\col] at (\xi+\chrW/2, 0.6) {\n};
        \node[\col] at (\xi+1.5*\chrW+0.025, 0.6) {\n};
    }
    \node[popmuted, anchor=east] at (8.5, 1.6) {Groupe G (21-22):};

    % --- PAIRE 23 : Sex chromosomes (Homme XY) : Ligne 2 (y=3) droite ---
    \pgfmathsetmacro{\xi}{10.5}
    % X
    \draw[popdark, chr=popdark] (\xi, 3) rectangle (\xi+\chrW, 3+1.2);
    \draw[centro] (\xi, 3+1.2*0.35) -- (\xi+\chrW, 3+1.2*0.35);
    \node[popdark] at (\xi+\chrW/2, 2.6) {X};
    % Y
    \pgfmathsetmacro{\xiY}{\xi+\chrW+0.05}
    \draw[popdark, chr=popdark] (\xiY, 3) rectangle (\xiY+\chrW, 3+0.5);
    \draw[centro] (\xiY, 3+0.5*0.1) -- (\xiY+\chrW, 3+0.5*0.1);
    \node[popdark] at (\xiY+\chrW/2, 2.6) {Y};
    \node[popdark, align=center] at (\xi+\chrW+0.025, 4.2) {\textbf{Paire 23}\\ \textbf{(Sexes)}};
    \node[popdark, align=center] at (\xi+\chrW+0.025, 2.2) {Homme (XY)};

    % Légende centromère / satellite
    \draw[centro] (0, -0.5) -- (0.5, -0.5) node[midway, below] {Centromère};
    \draw[popmuted, sat] (0.7, -0.7) -- (0.7, -0.3) node[midway, right] {Satellite (acrocentrique)};
\end{tikzpicture}
\legende{Caryotype humain schématisé (homme, 46,XY). Les 22 paires d'autosomes sont classées par taille décroissante et position du centromère en 7 groupes (A à G). La paire 23 (chromosomes sexuels) montre l'hétérologie XY. Les chromosomes acrocentriques (groupes D et G) portent des satellites (régions organisatrices des nucléoles). Bandes G non représentées pour la clarté.}
\end{popfigure}

\begin{definition}[Caryotype]
Le \textbf{caryotype} est la représentation photographique ordonnée de l'ensemble des chromosomes d'une cellule (généralement en métaphase), classés par paires d'homologues selon leur taille, la position de leur centromère et leur profil de bandage (G-banding). Il se note sous la forme : \textbf{46,XX} (femme) ou \textbf{46,XY} (homme).
\end{definition}

\textbf{Exploitation d'un caryotype : dénombrer, apparier, identifier le sexe}\par

L'analyse d'un caryotype suit une démarche rigoureuse en trois étapes.

\begin{methode}[Analyse d'un caryotype (fiche pratique)]
\begin{enumerate}
    \item \textbf{Dénombrer :} Compter le nombre total de centromères (donc de chromosomes). Vérifier si c'est $2n$ (ex: 46) ou anormal (ex: 47, 45).
    \item \textbf{Apparier :} Regrouper les chromosomes deux par deux selon : taille, position du centromère (médian, sous-médian, acrocentrique), profil de bandes. On obtient $n$ paires (23 pour l'homme). Numéroter les paires de 1 à 22 (autosomes) + paire 23 (sexuelles).
    \item \textbf{Identifier le sexe chromosomique :} Observer la paire 23.
    \begin{itemize}
        \item Deux chromosomes X homologues (grands, sous-médio-centriques) $\rightarrow$ \textbf{Femelle (46,XX)}.
        \item Un chromosome X + un chromosome Y (petit, acrocentrique) $\rightarrow$ \textbf{Mâle (46,XY)}.
    \end{itemize}
\end{enumerate}
\end{methode}

\begin{exoresolu}[Analyse d'un caryotype anormal : Trisomie 21]
\textbf{Document :} Caryotype obtenu sur cellules de liquide amniotique (amniocentèse à 16 SA). On compte 47 chromosomes. L'appariement montre 22 paires d'autosomes normales, mais pour la paire 21 (groupe G, très petits acrocentriques), on observe \textbf{trois} chromosomes 21 au lieu de deux. La paire 23 est XX.
\textbf{Questions :}
\begin{enumerate}
    \item Écrire la formule caryotypique de cet individu.
    \item Préciser le sexe chromosomique.
    \item Quelle est l'anomalie numérique ? Quelle en est la conséquence phénotypique connue ?
\end{enumerate}
\tcblower
\textbf{Solution détaillée :}
\begin{enumerate}
    \item \textbf{Formule : 47,XX,+21} (ou 47,XX,trisomie 21). 47 = nombre total ; XX = sexe ; +21 = chromosome 21 en excès.
    \item \textbf{Sexe : Femelle} (présence de deux chromosomes X, pas de Y).
    \item \textbf{Anomalie : Trisomie 21} (trois exemplaires du chromosome 21 au lieu de deux). C'est une anomalie numérique (non-disjonction lors de la méiose I ou II, le plus souvent maternelle).
    \item \textbf{Conséquence :} Syndrome de Down (trisomie 21). Retard mental variable, traits faciaux caractéristiques (fente palpébrale oblique, épicanthus), hypotonie, cardiopathies congénitales fréquentes (communication interventriculaire), risque accru de leucémie, maladie d'Alzheimer précoce. C'est la première cause de déficience intellectuelle d'origine génétique en France et au Cameroun (incidence $\approx$ 1/700 naissances).
\end{enumerate}
\end{exoresolu}

\begin{savaistu}[Trisomie 21 et dépistage prénatal au Cameroun]
Au Cameroun, le dépistage prénatal de la trisomie 21 se développe dans les grands centres (Hôpital Central, Hôpital Général, Centre Pasteur, cliniques privées à Yaoundé et Douala). Il repose sur :
\begin{itemize}
    \item \textbf{Marqueurs sériques maternels} (PAPP-A, $\beta$-hCG libre) + \textbf{clarté nucale} (échographie 11-13 SA) $\rightarrow$ calcul de risque.
    \item Si risque élevé ($> 1/250$) : proposition de \textbf{diagnostic prénatal invasif} : amniocentèse (15-18 SA) ou biopsie trophoblastique (11-13 SA) pour caryotypage.
    \item Depuis peu, le \textbf{DPNI} (Diagnostic Prénatal Non Invasif) sur ADN fœtal circulant dans le sang maternel (test NIPT/Harmony/Verifi) est disponible dans quelques laboratoires privés (coût élevé, non remboursé par la CNPS).
\end{itemize}
La prise en charge multidisciplinaire (pédiatre, cardiologue, orthophoniste, éducateur spécialisé) permet aujourd'hui une bien meilleure qualité de vie et une scolarisation en milieu ordinaire pour de nombreux enfants trisomiques 21.
\end{savaistu}

\textbf{Le chromosome, support physique de l'information génétique : synthèse}\par

L'ensemble de ces observations (constance du nombre, organisation en paires homologues biparentales, ségrégation lors de la méiose, anomalies numériques corrélées à des phénotypes précis) conduit à la conclusion majeure de la génétique classique (Théorie chromosomique de l'hérédité, Sutton-Boveri, 1902-1903) :

\begin{propriete}[Le chromosome, support de l'information génétique]
\begin{itemize}
    \item Les \textbf{chromosomes} sont les structures physiques qui portent l'information génétique (les \textbf{gènes}).
    \item Un \textbf{gène} est un segment d'un chromosome (un locus défini) qui détermine un caractère héréditaire.
    \item Les \textbf{allèles} d'un même gène occupent le \textbf{même locus} sur deux chromosomes homologues.
    \item La \textbf{méiose} assure la ségrégation des chromosomes homologues (donc des allèles) et le brassage interchromosomique, fondement de la variabilité génétique.
\end{itemize}
\end{propriete}

Mais de quoi est fait le chromosome ? La section suivante (« L'ADN, molécule de l'information génétique ») montrera que le chromosome n'est que l'emballage condensé d'une molécule chimique unique : l'acide désoxyribonucléique (ADN).

\begin{exercice}[Entraînement : Analyse de caryotypes]
\textbf{Doc 1 :} Caryotype 46,XY. \textbf{Doc 2 :} Caryotype 47,XX,+21. \textbf{Doc 3 :} Caryotype 45,X. \textbf{Doc 4 :} Caryotype 47,XXY.
\begin{enumerate}
    \item Pour chaque document, préciser : le nombre total de chromosomes, le sexe chromosomique, et nommer l'anomalie éventuelle (si 2n $\neq$ 46).
    \item Rappeler la définition de chromosomes homologues. Les chromosomes X et Y de l'homme sont-ils homologues ? Justifier.
    \item Une cellule en métaphase de mitose d'un homme (46,XY) contient-elle 46 ou 92 molécules d'ADN ? Expliquer.
\end{enumerate}
\end{exercice}

\begin{corrige}[Exercice 1]
\begin{enumerate}
    \item \textbf{Doc 1 (46,XY) :} 46 chromosomes, Mâle, Normal. \textbf{Doc 2 (47,XX,+21) :} 47 chromosomes, Femelle, Trisomie 21. \textbf{Doc 3 (45,X) :} 45 chromosomes, Femelle (un seul X), Monosomie X (Syndrome de Turner). \textbf{Doc 4 (47,XXY) :} 47 chromosomes, Mâle (présence d'un Y), Disomie X / Trisomie des gonosomes (Syndrome de Klinefelter).
    \item Chromosomes homologues = même taille, même centromère, mêmes gènes aux mêmes loci, origine biparentale (père/mère). X et Y ne sont \textbf{PAS homologues} (hétérologues) : tailles très différentes, gènes très différents (peu de régions d'homologie = régions pseudo-autosomiques PAR1/PAR2 nécessaires pour l'appariement méiotique).
    \item \textbf{92 molécules d'ADN.} En métaphase de mitose, chaque chromosome est bichromatidien (2 chromatides sœurs = 2 molécules d'ADN répliquées). 46 chromosomes $\times$ 2 = 92 molécules d'ADN. (En G1, il y en aurait 46).
\end{enumerate}
\end{corrige}

\coursec{Problème et hypothèses : où se trouve l'information génétique ?}

\courssub{Une situation concrète : la transmission des caractères}

Observons une épi de maïs cultivé dans les champs de l'Ouest-Cameroun : les grains sont jaunes ou blancs, lisses ou ridés. Chez l'homme, au sein d'une même famille à Yaoundé ou Garoua, certains enfants ont les yeux marrons, d'autres bleus ; certains sont porteurs du trait drépanocytaire (résistance au paludisme), d'autres non. Ces \cle{caractères héréditaires} se transmettent de génération en génération selon des régularités statistiques (lois de Mendel, programme de 2nde).

\textbf{Question centrale :} Quel est le support physique de cette information ? Où est-elle stockée dans la cellule ? De quelle nature chimique est-elle faite ?

\begin{itemize}
    \item \textbf{Hypothèse 1 (Cytoplasme) :} L'information serait dans le cytoplasme (mitochondries, chloroplastes, inclusions).
    \item \textbf{Hypothèse 2 (Noyau) :} L'information serait dans le noyau, sur les chromosomes (observés en division cellulaire).
    \item \textbf{Hypothèse 3 (Protéines vs ADN) :} Si c'est dans le noyau/chromosomes, le support est-il les \textbf{protéines} (diversité grande = 20 acides aminés) ou l'\textbf{ADN} (structure apparemment répétitive, 4 bases) ?
\end{itemize}

Pour lever ces hypothèses, l'histoire des sciences a utilisé une démarche expérimentale rigoureuse : \textbf{le transfert de noyau}.

\coursec{Localisation de l'information génétique : les expériences de transfert de noyau}

\courssub{L'expérience fondatrice de Hammerling (1930-1953) sur l'Acétabulaire}

L'algue unicellulaire géante \emph{Acetabularia mediterranea} (pied, tige, calice en forme de parapluie) et \emph{Acetabularia crenulata} (calice découpé) sont des modèles idéaux : un seul noyau volumineux situé dans le \textbf{pied} (rhizoïde), loin du calice.

\begin{experience}[Transfert de noyau chez Acetabularia — J. Hammerling]
\textbf{Observation :} L'espèce \emph{mediterranea} a un calice entier ; l'espèce \emph{crenulata} a un calice découpé. Le noyau est dans le pied.
\textbf{Protocole :}
\begin{enumerate}
    \item Couper le pied (contenant le noyau) d'une \emph{mediterranea} et le greffer sur le tronçon apical (sans noyau) d'une \emph{crenulata} $\rightarrow$ Chimère M-pied / C-tête.
    \item Couper le pied d'une \emph{crenulata} et le greffer sur le tronçon apical d'une \emph{mediterranea} $\rightarrow$ Chimère C-pied / M-tête.
    \item Laisser régénérer le calice.
\end{enumerate}
\textbf{Résultats :}
\begin{itemize}
    \item Chimère M-pied / C-tête $\rightarrow$ Nouveau calice \textbf{entier} (type \emph{mediterranea}, donneur du noyau).
    \item Chimère C-pied / M-tête $\rightarrow$ Nouveau calice \textbf{découpé} (type \emph{crenulata}, donneur du noyau).
    \item Si on greffe les \textbf{deux noyaux} ensemble $\rightarrow$ Calice \textbf{intermédiaire} (caractères des deux espèces).
\end{itemize}
\textbf{Interprétation :} Le noyau (et non le cytoplasme) dirige la morphogenèse du calice. L'information héréditaire est dans le \textbf{noyau}.
\textbf{Conclusion :} Le noyau est le siège de l'information génétique déterminant les caractères de l'espèce.
\end{experience}

\courssub{Confirmation chez l'animal : l'expérience de Gurdon (1962) — Prélude au clonage}

John Gurdon (Prix Nobel 2012) utilise la grenouille \emph{Xenopus laevis}.

\begin{experience}[Transfert de noyau différencié — J. Gurdon]
\textbf{Protocole :}
\begin{enumerate}
    \item Prélever un noyau d'une cellule intestinale différenciée d'un têtard (noyau « spécialisé »).
    \item L'injecter dans un ovule énucléé (dont on a détruit le noyau par UV).
    \item Observer le développement.
\end{enumerate}
\textbf{Résultat :} L'ovule se développe en un têtard normal, puis une grenouille adulte, \textbf{génétiquement identique} au donneur du noyau intestinal.
\textbf{Interprétation :} Un noyau de cellule différenciée conserve \textbf{toute l'information génétique} nécessaire à la construction de l'organisme entier (totipotence). L'information est stable et complète dans le noyau de \textbf{toute} cellule somatique.
\end{experience}

\begin{aretenir}[Localisation de l'information génétique]
\centering
\textbf{L'information génétique est localisée dans le \cle{noyau} de la cellule, supportée par les \cle{chromosomes}.}
\end{aretenir}

\coursec{Les chromosomes, supports de l'information génétique}

\courssub{Observation microscopique : le caryotype}

Au microscope optique, après coloration (Giemsa, Feulgen), les chromosomes apparaissent lors de la division cellulaire (mitose) sous forme de bâtonnets compacts. En interphase, ils se décondensent en \cle{chromatine} (invisible individuellement).

\begin{definition}[Chromosome]
Un \textbf{chromosome} est une structure filamenteuse, visible au microscope optique lors de la division cellulaire, constituée d'une \textbf{molécule d'ADN continue} enroulée autour de protéines (histones). Il porte une partie de l'information génétique.
\end{definition}

\begin{popfigure}
\begin{tikzpicture}[scale=0.9, every node/.style={font=\small}]
    % Chromosome monochromatidien (G1)
    \begin{scope}[shift={(-4,0)}]
        \draw[popblue, very thick, fill=popblue!10] (-0.6, -2) -- (-0.6, -0.5) .. controls (-0.6,0) and (0.6,0) .. (0.6, -0.5) -- (0.6, -2) .. controls (0.6,-2.5) and (-0.6,-2.5) .. (-0.6,-2);
        \draw[poporange, thick, fill=poporange!30] (-0.2, -1.3) rectangle (0.2, -1.0);
        \node[poporange] at (0, -1.15) {Centromère};
        \node[align=center] at (0, -2.8) {Monochromatidien\\(1 chromatide = 1 molécule ADN)};
        \node[align=center, popblue] at (0, 0.5) {Phase G1 / Interphase};
    \end{scope}
    
    % Chromosome bichromatidien (G2 / Mitose)
    \begin{scope}[shift={(4,0)}]
        % Chromatide 1
        \draw[popblue, very thick, fill=popblue!10] (-1.0, -2) -- (-1.0, -0.5) .. controls (-1.0,0) and (-0.2,0) .. (-0.2, -0.5) -- (-0.2, -1.2);
        % Chromatide 2
        \draw[poppurple, very thick, fill=poppurple!10] (0.2, -1.2) -- (0.2, -0.5) .. controls (0.2,0) and (1.0,0) .. (1.0, -0.5) -- (1.0, -2);
        % Centromère commun
        \draw[poporange, thick, fill=poporange!30] (-0.2, -1.3) rectangle (0.2, -1.0);
        \node[poporange] at (0, -1.15) {Centromère};
        % Bras
        \node[left] at (-1.1, -1.5) {Bras court (p)};
        \node[right] at (1.1, -1.5) {Bras long (q)};
        \node[align=center] at (0, -2.8) {Bichromatidien\\(2 chromatides sœurs = 2 molécules ADN identiques)};
        \node[align=center, popblue] at (0, 0.5) {Phase G2 / Mitose (après réplication)};
    \end{scope}
\end{tikzpicture}
\legende{Structure d'un chromosome. \textbf{Gauche :} Chromosome monochromatidien (une chromatide, une molécule d'ADN). \textbf{Droite :} Chromosome bichromatidien (deux chromatides sœurs issues de la réplication de l'ADN, jointes au centromère). Chaque chromatide contient une molécule d'ADN double brin.}
\end{popfigure}

\courssub{Les paires de chromosomes homologues}

Chez l'homme (2n = 46), les chromosomes s'organisent en \textbf{23 paires homologues} :
\begin{itemize}
    \item \textbf{22 paires d'autosomes} (identiques en taille, forme, position du centromère, gènes aux mêmes loci).
    \item \textbf{1 paire de gonosomes (chromosomes sexuels) :} XX (femelle) ou XY (mâle).
\end{itemize}
\begin{itemize}
    \item Un chromosome de la paire vient du \textbf{père} (via le spermatozoïde), l'autre de la \textbf{mère} (via l'ovocyte).
    \item Ils portent les \textbf{mêmes gènes} aux \textbf{mêmes loci}, mais pas nécessairement les \textbf{mêmes allèles} (versions du gène).
\end{itemize}

\begin{propriete}[Vocabulaire essentiel pour le BEPC]
\begin{itemize}
    \item \cle{Chromatide} : Un des deux brins d'un chromosome bichromatidien (une molécule d'ADN).
    \item \cle{Chromatides sœurs} : Les deux chromatides d'un chromosome bichromatidien, \textbf{identiques} (issues de la réplication).
    \item \cle{Chromosomes homologues} : Deux chromosomes (un paternel, un maternel) de même forme/taille, portant les mêmes gènes aux mêmes loci.
    \item \cle{Locus} : Position précise d'un gène sur un chromosome.
    \item \cle{Allèle} : Version alternative d'un gène (ex : allèle $A$ = groupe sanguin A, allèle $B$ = groupe sanguin B).
\end{itemize}
\end{propriete}

\coursec{L'ADN, molécule de l'information génétique}

\courssub{De la structure visible à la molécule invisible : une question chimique}

Nous savons que les chromosomes (dans le noyau) sont les supports physiques de l'hérédité. Mais de quoi sont faits exactement ces chromosomes ? Quelle est la nature chimique de cette « matière héréditaire » capable de coder la couleur des yeux, la forme des graines de maïs ou la résistance au paludisme ?

Dès le début du XX\textsuperscript{e} siècle, les biochimistes analysent la composition chimique des noyaux isolés (thymus de veau, spermatozoïdes de saumon, levure). Les résultats sont nets : les chromosomes sont constitués de deux types de macromolécules majoritaires : des \cle{protéines} (histones et non-histones) et de l'\cle{ADN} (acide désoxyribonucléique).

Pendant longtemps, les protéines, par leur grande diversité (20 acides aminés), ont paru les candidates idéales. L'ADN, avec seulement quatre bases, semblait trop « monotone ». L'histoire a tranché : c'est bien l'ADN, par la \emph{séquence} de ses quatre bases, qui code l'information.

\begin{experience}[Analyse chimique des chromosomes — Démarche historique]
\textbf{Observation :} Au microscope optique, après coloration spécifique (Feulgen pour l'ADN, Giemsa pour l'ADN+protéines), les chromosomes apparaissent comme des fils ou bâtonnets compacts dans le noyau en division.
\textbf{Question :} Quelle est la nature chimique du matériel coloré par la réaction de Feulgen (spécifique de l'ADN) ?
\textbf{Protocole (simplifié) :}
\begin{enumerate}
    \item Isoler des noyaux en grande quantité (ex. thymus de veau, levure de boulanger, spermatozoïdes de hareng).
    \item Extraire la chromatine par traitement acide/basique et enzymatique (protéases pour digérer les protéines, RNAses pour éliminer l'ARN).
    \item Dosage élémentaire (C, H, O, N, P) et identification des monomères après hydrolyse complète.
\end{enumerate}
\textbf{Résultats :}
\begin{itemize}
    \item Présence de \textbf{phosphore (P)} en quantité significative $\rightarrow$ groupe phosphate.
    \item Présence d'un \textbf{sucre à 5 atomes de carbone (pentose)} : le \textbf{désoxyribose} (il manque un groupement OH sur le carbone 2' par rapport au ribose de l'ARN).
    \item Présence de \textbf{quatre bases azotées} : \textbf{Adénine (A)}, \textbf{Guanine (G)} (purines, double cycle), \textbf{Cytosine (C)}, \textbf{Thymine (T)} (pyrimidines, cycle simple). \emph{Note : l'Uracile (U) est absent de l'ADN, remplacé par la Thymine.}
    \item Présence de \textbf{protéines} (histones riches en lysine/arginine, basiques).
\end{itemize}
\textbf{Interprétation :} Le chromosome est un complexe \textbf{ADN-protéines} appelé \textbf{chromatine}. L'ADN est un polymère de \cle{nucléotides}.
\textbf{Conclusion :} L'ADN est un constituant majeur, constant et universel des chromosomes de tous les êtres vivants (sauf certains virus à ARN).
\end{experience}

\courssub{La structure de l'ADN : le modèle de la double hélice (Watson, Crick, 1953)}

\cle{Comment une molécule « simple » peut-elle porter une information complexe ?} La réponse est venue de la détermination de sa structure tridimensionnelle. En 1953, James Watson et Francis Crick, s'appuyant sur les données de diffraction aux rayons X de \textbf{Rosalind Franklin} et Maurice Wilkins (cliché 51), et sur les règles de proportionnalité des bases d'Erwin Chargaff, proposent le modèle universel de l'ADN : la \cle{double hélice}.

\begin{definition}[ADN — Acide DésoxyriboNucléique]
L'\textbf{ADN} est une macromolécule biologique, support de l'information génétique, formée de \textbf{deux brins polynucléotidiques} enroulés l'un autour de l'autre en une \textbf{double hélice} droite (hélice B). Les deux brins sont \textbf{antiparallèles} (sens 5' $\rightarrow$ 3' opposé) et \textbf{complémentaires} par appariement spécifique des bases azotées.
\end{definition}

\begin{popfigure}
\begin{tikzpicture}[scale=0.85, every node/.style={font=\small}]
    % Représentation 2D "échelle tordue" (vue de côté simplifiée)
    \def\h{0.34} % hauteur par paire (nm)
    \def\w{2.0}  % largeur hélice (nm)
    \def\n{14}   % nombre de paires
    
    \colorlet{backbone}{popblue}
    \colorlet{baseA}{poporange}
    \colorlet{baseT}{popgreen}
    \colorlet{baseG}{poppurple}
    \colorlet{baseC}{poppink}
    \colorlet{hbond}{popmuted}
    
    % Squelettes (deux courbes sinusoïdales décalées)
    \draw[backbone, very thick, smooth, variable=\t, domain=0:\n*\h, samples=200]
        plot ({-\w/2 * sin(360*\t/\h/10)}, {\t}); % Brin 1
    \draw[backbone, very thick, smooth, variable=\t, domain=0:\n*\h, samples=200]
        plot ({\w/2 * sin(360*\t/\h/10)}, {\t});  % Brin 2
    
    % Paires de bases (barres horizontales au centre)
    \foreach \i in {1,...,\n} {
        \pgfmathsetmacro{\y}{\i*\h - \h/2};
        \pgfmathsetmacro{\r}{mod(\i*7, 4)}; % pseudo-aléatoire fixe
        \ifdim\r pt < 1pt \def\colL{baseA}\def\colR{baseT}\def\labL{A}\def\labR{T}\def\hb{2}
        \else\ifdim\r pt < 2pt \def\colL{baseT}\def\colR{baseA}\def\labL{T}\def\labR{A}\def\hb{2}
        \else\ifdim\r pt < 3pt \def\colL{baseG}\def\colR{baseC}\def\labL{G}\def\labR{C}\def\hb{3}
        \else \def\colL{baseC}\def\colR{baseG}\def\labL{C}\def\labR{G}\def\hb{3}
        \fi\fi\fi
        
        % Liaisons H (pointillés au centre)
        \draw[hbond, dashed, thick] (0, \y) -- (0, \y); % point central pour repère
        \draw[hbond, dashed, thick] (-0.15, \y) -- (0.15, \y);
        \ifnum\hb=3 \draw[hbond, dashed, thick] (-0.1, \y+0.02) -- (0.1, \y+0.02); \fi % 3ème trait simulé
        
        % Bases (rectangles attachés aux squelettes)
        \pgfmathsetmacro{\xL}{-\w/2 * sin(360*\y/\h/10)};
        \pgfmathsetmacro{\xR}{\w/2 * sin(360*\y/\h/10)};
        
        \fill[\colL!80!black] (\xL-0.15, \y-0.07) rectangle (\xL+0.15, \y+0.07);
        \fill[\colR!80!black] (\xR-0.15, \y-0.07) rectangle (\xR+0.15, \y+0.07);
        
        % Étiquettes quelques bases
        \ifnum\i=3 \node[left=2pt, \colL!80!black] at (\xL, \y) {\labL}; \node[right=2pt, \colR!80!black] at (\xR, \y) {\labR}; \fi
        \ifnum\i=8 \node[left=2pt, \colL!80!black] at (\xL, \y) {\labL}; \node[right=2pt, \colR!80!black] at (\xR, \y) {\labR}; \fi
    }
    
    % Flèches antiparallèles
    \draw[poporange, thick, -{Stealth[length=3mm]}] (-\w/2-0.6, -0.2) -- (-\w/2-0.6, 0.6) node[left, pos=0.5] {5' $\rightarrow$ 3'};
    \draw[popgreen, thick, -{Stealth[length=3mm]}] (\w/2+0.6, \n*\h+0.2) -- (\w/2+0.6, \n*\h-0.6) node[right, pos=0.5] {3' $\leftarrow$ 5'};
    
    % Légende
    \begin{scope}[shift={(-4.5, -0.5)}]
        \fill[baseA] (0,0) rectangle (0.35,0.35); \node[right] at (0.45,0.17) {Adénine (A)};
        \fill[baseT] (0,-0.5) rectangle (0.35,-0.15); \node[right] at (0.45,-0.32) {Thymine (T)};
        \fill[baseG] (0,-1.0) rectangle (0.35,-0.65); \node[right] at (0.45,-0.82) {Guanine (G)};
        \fill[baseC] (0,-1.5) rectangle (0.35,-1.15); \node[right] at (0.45,-1.32) {Cytosine (C)};
        \draw[hbond, dashed] (1.5, -0.15) -- (2.3, -0.15) node[right] {2 liaisons H};
        \draw[hbond, dashed] (1.5, -1.15) -- (2.3, -1.15) node[right] {3 liaisons H};
    \end{scope}
    
    % Cotes
    \draw[<->, popdark] (-\w/2, \n*\h+0.3) -- (\w/2, \n*\h+0.3) node[midway, above] {2 nm};
    \draw[<->, popdark] (-\w/2-0.3, 0) -- (-\w/2-0.3, \h) node[midway, left] {0,34 nm};
    \node[popdark, rotate=90] at (-\w/2-0.8, \h/2) {10 paires/bases par tour};
\end{tikzpicture}
\legende{Schéma de la double hélice d'ADN (modèle B). Les deux brins (squelette sucre-phosphate en bleu) sont antiparallèles. Les bases (A/T orange/vert, G/C violet/rose) s'apparient au centre par liaisons hydrogène (pointillés) : A=T (2 liaisons), G$\equiv$C (3 liaisons). Largeur constante 2 nm.}
\end{popfigure}

\courssub{Le nucléotide : l'unité de base de l'ADN}

Chaque brin de l'ADN est un polymère linéaire : un enchaînement de \cle{nucléotides}. Un nucléotide d'ADN est constitué de trois éléments covalemment liés :

\begin{enumerate}
    \item Un \textbf{sucre pentose} : le \textbf{désoxyribose} (C$_5$H$_{10}$O$_4$). Ses atomes de carbone sont numérotés 1' à 5'. L'absence de groupement OH sur le C2' (d'où « désoxy ») confère une plus grande stabilité chimique à l'ADN par rapport à l'ARN.
    \item Un \textbf{groupe phosphate} (acide phosphorique, PO$_4^{3-}$) estérifié sur le carbone 5' du sucre.
    \item Une \textbf{base azotée} (cyclique, basique) liée sur le carbone 1' du sucre par une \textbf{liaison osidique (N-glycosidique)}.
\end{enumerate}

\begin{center}
\begin{tikzpicture}[scale=1.0, every node/.style={font=\small}]
    % Cycle furanose (désoxyribose) - représentation standard
    \coordinate (C1) at (0,0);
    \coordinate (C2) at (1.1,0);
    \coordinate (C3) at (1.6,0.9);
    \coordinate (C4) at (0.5,1.6);
    \coordinate (O4) at (-0.6,0.9);
    \coordinate (C5) at (-0.9,-0.3);
    
    % Cycle
    \draw[popblue, thick, fill=popblue!10] (C1) -- (C2) -- (C3) -- (C4) -- (O4) -- cycle;
    \node[popblue] at (0.5, 0.7) {Désoxyribose};
    
    % Base sur C1 (haut)
    \draw[poporange, thick] (C1) -- ++(0, 1.1) node[above] {Base azotée (A, T, G ou C)};
    \node[poporange, font=\tiny] at (0, 1.3) {Liaison osidique};
    
    % Phosphate sur C5 (bas gauche)
    \draw[popgreen, thick] (C5) -- ++(-1.2, -0.8) node[left] {Phosphate (sur C5')};
    \node[popgreen, font=\tiny] at (-1.4, -0.4) {Liaison ester};
    
    % OH libre sur C3 (liaison vers le suivant)
    \draw[poppurple, thick, dashed] (C3) -- ++(1.0, 0.5) node[right] {OH libre sur C3'};
    
    % Numérotation carbones
    \foreach \pos/\label in {C1/1', C2/2', C3/3', C4/4', O4/O, C5/5'}
        \node[font=\tiny, popdark] at (\pos) {\label};
        
    % Sens 5' -> 3'
    \draw[popdark, thick, -{Stealth[length=2mm]}] (-3.0, -1.8) -- (-1.8, -1.8) node[midway, below] {Sens 5' $\rightarrow$ 3'};
\end{tikzpicture}
\end{center}

\medskip
La polymérisation se fait par \textbf{liaison phosphodiester} entre le groupe phosphate du carbone 5' d'un nucléotide et le groupe hydroxyle (OH) libre du carbone 3' du nucléotide précédent. Cela crée une \textbf{asymétrie directionnelle} : chaque brin a une extrémité \textbf{5'} (phosphate libre) et une extrémité \textbf{3'} (OH libre). Les deux brins de la double hélice sont \textbf{antiparallèles} : l'un court 5' $\rightarrow$ 3', l'autre 3' $\rightarrow$ 5'.

\courssub{Les règles de Chargaff : la clé de l'appariement des bases}

Avant même la découverte de la structure, le biochimiste Erwin Chargaff (années 1940-50) a dosé les bases dans l'ADN de nombreuses espèces (bactéries, levures, mammifères, plantes). Ses résultats, d'une régularité frappante, sont des \textbf{faits expérimentaux} fondamentaux.

\begin{propriete}[Règles de Chargaff — Fait expérimental]
Dans l'ADN de \textbf{tout organisme} (quel que soit l'espèce, le tissu, l'âge, pour de l'ADN double brin) :
\begin{enumerate}
    \item La quantité d'\textbf{Adénine (A)} est \textbf{égale} à la quantité de \textbf{Thymine (T)} : $\%A = \%T$.
    \item La quantité de \textbf{Guanine (G)} est \textbf{égale} à la quantité de \textbf{Cytosine (C)} : $\%G = \%C$.
    \item \textbf{Corollaire :} La somme des purines (A+G) égale la somme des pyrimidines (T+C) : $\%A + \%G = \%T + \%C = 50\%$.
    \item Le rapport $\frac{A+G}{T+C} = 1$ est constant. En revanche, le rapport $\frac{A+T}{G+C}$ \textbf{varie selon les espèces} (ex. : $\approx 1,5$ chez l'homme, $\approx 1,0$ chez le maïs, $\approx 0,7$ chez certaines bactéries). Ce rapport est une \textbf{signature spécifique de l'espèce}.
\end{enumerate}
\end{propriete}

\begin{popfigure}
\begin{tikzpicture}
\begin{axis}[
    ybar,
    bar width=10pt,
    width=\linewidth,
    height=7cm,
    ylabel={Pourcentage (\% en moles)},
    xlabel={Bases azotées},
    symbolic x coords={A, T, G, C},
    xtick=data,
    ymin=0, ymax=35,
    ymajorgrids=true,
    grid style=dashed,
    nodes near coords,
    nodes near coords align={vertical},
    every node near coord/.append style={font=\small, popdark},
    enlarge x limits=0.25,
    legend style={at={(0.5,-0.18)}, anchor=north, legend columns=-1},
    tick label style={font=\small},
    label style={font=\small}
]
\addplot[fill=poporange, draw=poporange!80!black] coordinates {(A, 30.4) (T, 29.6) (G, 19.9) (C, 20.1)};
\addlegendentry{ADN Humain (données Chargaff)}
\addplot[fill=popblue, draw=popblue!80!black, opacity=0.4] coordinates {(A, 25) (T, 25) (G, 25) (C, 25)};
\addlegendentry{Distribution équitable (théorique)}
\end{axis}
\end{tikzpicture}
\legende{Pourcentages molaires des bases dans l'ADN humain. On observe $\%A \approx \%T \approx 30\%$ et $\%G \approx \%C \approx 20\%$. La somme A+T $\approx$ 60\%, G+C $\approx$ 40\%. Le rapport (A+T)/(G+C) $\approx$ 1,5 est caractéristique de l'espèce humaine.}
\end{popfigure}

\noindent\textbf{Interprétation (Watson \& Crick) :} L'égalité stricte A=T et G=C ne peut s'expliquer que si A et T, d'une part, G et C, d'autre part, sont \textbf{structuralement appariés} deux à deux au centre de l'hélice. La géométrie de l'hélice impose qu'une purine (grosse, double cycle) s'apparie avec une pyrimidine (petite, cycle simple) pour garder une largeur constante de 2 nm.
\begin{itemize}
    \item \textbf{A = T} : 2 liaisons hydrogène (plus faible, fond plus facilement à la chaleur).
    \item \textbf{G $\equiv$ C} : 3 liaisons hydrogène (plus fort, plus stable thermiquement).
\end{itemize}
C'est la \cle{complémentarité des bases} : la séquence d'un brin dicte obligatoirement la séquence de l'autre.

\begin{methode}[Déterminer la séquence du brin complémentaire]
\textbf{Donné :} La séquence d'un brin d'ADN (sens 5' $\rightarrow$ 3').
\textbf{But :} Écrire la séquence du brin opposé (sens 3' $\rightarrow$ 5' ou 5' $\rightarrow$ 3').
\textbf{Étapes :}
\begin{enumerate}
    \item Écrire la séquence donnée (ex. : 5' - \textbf{A T G C C A} - 3').
    \item Appliquer les règles d'appariement \textbf{base par base} : A$\leftrightarrow$T, G$\leftrightarrow$C.
    \item Écrire les bases complémentaires \textbf{en face} (alignées).
    \item \textbf{Respecter l'antiparallélisme :} Si le brin donné est 5' $\rightarrow$ 3', le brin complémentaire s'écrit 3' $\leftarrow$ 5' (ou 5' $\rightarrow$ 3' en le lisant à l'envers).
\end{enumerate}
\textbf{Exemple :}
\begin{align*}
\text{Brin 1 (donné) :} \quad & 5' - \textbf{A} - \textbf{T} - \textbf{G} - \textbf{C} - \textbf{C} - \textbf{A} - 3' \\
                              & \quad | \quad | \quad | \quad | \quad | \quad | \\
\text{Brin 2 (complémentaire) :} \quad & 3' - \textbf{T} - \textbf{A} - \textbf{C} - \textbf{G} - \textbf{G} - \textbf{T} - 5' \\
\text{Écrit en 5' $\rightarrow$ 3' :} \quad & 5' - \textbf{T} - \textbf{G} - \textbf{G} - \textbf{C} - \textbf{A} - \textbf{T} - 3'
\end{align*}
\end{methode}

\begin{exoresolu}[Application des règles de Chargaff et complémentarité]
\textbf{Énoncé :} On analyse l'ADN d'une bactérie isolée d'un sol volcanique au Cameroun (flanc du Mont Cameroun). Le dosage montre que la cytosine (C) représente 18 \% des bases.
\begin{enumerate}
    \item Déduire le pourcentage de chaque base (A, T, G).
    \item Calculer le rapport (A+T)/(G+C).
    \item Écrire la séquence complémentaire (en 5' $\rightarrow$ 3') du fragment : 5' - A A T G C C - 3'.
\end{enumerate}
\textbf{Solution détaillée :}
\begin{enumerate}
    \item \textbf{Règles de Chargaff :} $\%C = \%G = 18\%$.
    $\%G + \%C = 36\%$.
    $\%A + \%T = 100\% - 36\% = 64\%$.
    Comme $\%A = \%T$, alors $\%A = \%T = 32\%$.
    \textbf{Bilan :} A = 32\%, T = 32\%, G = 18\%, C = 18\%.
    \item \textbf{Rapport (A+T)/(G+C) = 64 / 36 = 1,78.} Ce rapport élevé suggère un ADN riche en A+T (fréquent chez certaines bactéries thermophiles).
    \item \textbf{Complémentarité (antiparallèle) :}
    \begin{align*}
    5' - \textbf{A} - \textbf{A} - \textbf{T} - \textbf{G} - \textbf{C} - \textbf{C} - 3' \\
    \quad | \quad | \quad | \quad | \quad | \quad | \\
    3' - \textbf{T} - \textbf{T} - \textbf{A} - \textbf{C} - \textbf{G} - \textbf{G} - 5'
    \end{align*}
    \textbf{Réponse (5' $\rightarrow$ 3') :} 5' - \textbf{G G C A T T} - 3'.
\end{enumerate}
\end{exoresolu}

\courssub{L'information génétique : l'ordre des bases}

Si la structure (squelette, double hélice) est \textbf{identique} chez tous les êtres vivants (universalité de la molécule), l'\textbf{information} réside uniquement dans la \textbf{séquence linéaire} des quatre bases le long d'un brin. C'est un \textbf{langage à 4 lettres} (A, T, G, C).

\begin{aretenir}[Chaîne logique de localisation de l'information génétique]
\centering
\textbf{Cellule} $\xrightarrow{\text{contient}}$ \textbf{Noyau} $\xrightarrow{\text{contient}}$ \textbf{Chromosomes} $\xrightarrow{\text{constitués de}}$ \textbf{ADN (double hélice)} $\xrightarrow{\text{information =}}$ \textbf{Séquence des bases (A, T, G, C)}
\end{aretenir}

\begin{attention}[Pièges fréquents à éviter au BEPC]
\begin{itemize}
    \item \textbf{Confondre Chromosome et ADN :} Le chromosome est la \textbf{structure microscopique visible} (condensée en mitose), faite d'ADN \emph{enroulé autour de protéines (histones)}. L'ADN est la \textbf{molécule chimique} (le fil continu). \textit{Analogie : le chromosome est le peloton de laine ; l'ADN est le fil de laine.}
    \item \textbf{Oublier l'antiparallélisme :} Écrire le brin complémentaire dans le même sens (5' $\rightarrow$ 3') sans l'inverser. \textbf{Toujours vérifier :} 5' face à 3'.
    \item \textbf{Confondre Uracile et Thymine :} L'ADN contient de la \textbf{Thymine (T)}. L'ARN contient de l'\textbf{Uracile (U)}. Pas de U dans l'ADN (sauf exceptions rares post-réplication).
    \item \textbf{Croire que A=T et G=C en \emph{nombre d'atomes} ou en \emph{masse} :} C'est en \textbf{pourcentage molaire} (nombre de moles de bases).
    \item \textbf{Appliquer Chargaff à l'ADN simple brin :} Les règles ne valent que pour l'ADN \textbf{double brin} (duplexe).
\end{itemize}
\end{attention}

\courssub{Du brin d'ADN au gène : notion introductive}

Un chromosome unique (ex. chromosome 1 humain) contient une seule molécule d'ADN continue de $\approx$ 250 millions de paires de bases. Cette immense molécule est découpée en \textbf{unités fonctionnelles} : les \cle{gènes}.

\begin{definition}[Gène]
Un \textbf{gène} est un \textbf{segment d'ADN} (une portion de la molécule) qui porte l'information nécessaire à la synthèse d'un \textbf{produit fonctionnel} (généralement une \textbf{protéine} : enzyme, anticorps, hémoglobine, récepteur, ou un ARN fonctionnel). Un gène correspond à \textbf{un caractère héréditaire} (au sens large : structure d'une protéine).
\end{definition}

\begin{savaistu}[L'ADN, une molécule géante microscopique]
\begin{itemize}
    \item Dans chaque noyau de vos cellules (sauf globules rouges matures), les 46 chromosomes (23 paires) contiennent au total \textbf{$\approx$ 2 mètres d'ADN} déroulé (diamètre 2 nm).
    \item Cet ADN est compacté $\approx$ 10 000 fois grâce aux histones (nucléosomes $\rightarrow$ fibre de chromatine $\rightarrow$ chromosome mitotique) pour tenir dans un noyau de $\approx$ 5 à 10 $\mu$m.
    \item Le génome humain = $\approx$ 3,2 milliards de paires de bases (3,2 Gb), pour environ \textbf{20 000 à 25 000 gènes} codant des protéines. Le reste (98 \%) est de l'ADN non codant (régulation, structure, éléments transposables).
    \item Au Cameroun, la diversité génétique humaine est la plus grande au monde. L'étude de l'ADN mitochondrial (transmis par la mère) et du chromosome Y (transmis par le père) a permis de retracer les migrations des populations Bantoues, Pygmées, Sahéliennes depuis la vallée du Rift.
\end{itemize}
\end{savaistu}

\coursec{TP N°1 : Mise en évidence du matériel génétique d'une cellule}

\courssub{Objectif}
Mettre en évidence la présence d'ADN dans le noyau des cellules et sa localisation sur les chromosomes, par deux approches complémentaires : coloration spécifique (microscopie) et extraction (biochimie).

\courssub{Activité A : Réaction de Feulgen (Localisation microscopique de l'ADN)}

\begin{experience}[TP1-A : Coloration de Feulgen sur pointe de racine d'oignon ou d'ail]
\textbf{Matériel :} Pointes de racines d'oignon/ail (mérostème en division active), acide chlorhydrique 1N, réactif de Feulgen (acide sulfureux + fuchsine décolorée), bleu de méthylène (contre-coloration), lames, lamelles, microscope optique, bain-marie à 60°C.
\textbf{Protocole :}
\begin{enumerate}
    \item \textbf{Fixation :} Pointes de racines dans éthanol/acide acétique (3/1) 24h (ou fraîches).
    \item \textbf{Hydrolyse acide :} Plonger 8-10 min dans HCl 1N à 60°C (élimine l'ARN, dépurine l'ADN $\rightarrow$ aldéhydes sur désoxyribose).
    \item \textbf{Rinçage :} Eau distillée abondante.
    \item \textbf{Coloration :} Plonger 30-60 min dans le réactif de Feulgen (obscurité, température ambiante). La fuchsine se fixe sur les aldéhydes $\rightarrow$ couleur \textbf{rose-magenta vif}.
    \item \textbf{Rinçage :} Eau distillée (3 bains).
    \item \textbf{Contre-coloration (optionnel) :} Bleu de méthylène 1 min (colorie le cytoplasme/noyau en bleu, chromosomes en rose).
    \item \textbf{Montage :} Écraser une pointe de racine sur lame (goutte d'eau/gel), poser lamelle, observer $\times$400.
\end{enumerate}
\textbf{Observations attendues :}
\begin{itemize}
    \item Noyaux en interphase : taches roses diffuses (chromatine).
    \item Cellules en mitose (prophase, métaphase, anaphase, télophase) : \textbf{chromosomes roses magenta} nets, bien individualisés.
    \item Cytoplasme : incolore (ou bleu si contre-coloration).
\end{itemize}
\textbf{Interprétation :} La coloration spécifique de l'ADN (réaction de Feulgen) se localise \textbf{exclusivement dans le noyau}, sur les \textbf{chromosomes}. L'ADN est bien le constituant chromatique des chromosomes.
\textbf{Témoin de contrôle :} Une lame traitée par la \textbf{DNase} (digestion de l'ADN) avant Feulgen $\rightarrow$ \textbf{pas de coloration rose}. Preuve de la spécificité.
\end{experience}

\courssub{Activité B : Extraction grossière de l'ADN (Nature chimique)}

\begin{experience}[TP1-B : Extraction d'ADN à partir de banane plantain ou de foie de bœuf (abattoir de Ngaoundéré/Yaoundé)]
\textbf{Matériel :} Banane plantain mûre (ou foie frais), eau, sel (NaCl), liquide vaisselle (détergent), alcool éthylique 96\% (glacé au congélateur), mixeur, passoire, bécher, baguette de verre, tube à essai.
\textbf{Protocole :}
\begin{enumerate}
    \item \textbf{Lysis :} Mixer 50 g de tissu + 200 mL eau + 1 c.à.c sel + 2 c.à.s détergent. Le détergent solubilise les membranes (lipides), le sel neutralise les charges négatives de l'ADN.
    \item \textbf{Filtration :} Filtrer sur gaze/passoire $\rightarrow$ filtrat (ADN + protéines + sucres en solution).
    \item \textbf{Précipitation :} Verser délicatement l'alcool glacé le long de la paroi du tube (2 volumes). L'ADN est insoluble dans l'alcool froid.
    \item \textbf{Observation :} Au contact des deux liquides, un \textbf{filament blanchâtre, visqueux, gélatineux} (l'ADN) précipite et remonte dans l'alcool. On peut l'enrouler sur la baguette (spooling).
\end{enumerate}
\textbf{Interprétation :} On isole une substance macromoléculaire, filamenteuse, insoluble dans l'alcool, présente dans le noyau. C'est l'ADN.
\textbf{Lien avec le cours :} Cette substance est le support chimique de l'information observée sur les chromosomes au TP1-A.
\end{experience}

\begin{aretenir}[Bilan TP N°1]
\centering
La réaction de Feulgen localise l'ADN \textbf{dans le noyau, sur les chromosomes}. L'extraction montre que l'ADN est une \textbf{macromolécule filamenteuse, insoluble dans l'alcool}, extraite du noyau. L'ADN est bien le matériel génétique.
\end{aretenir}

\coursec{TP N°2 : Maquettes de paires de chromosomes homologues et applications}

\courssub{Objectif}
Construire des modèles physiques pour visualiser la structure des chromosomes (monochromatidiens / bichromatidiens), la notion de paire homologue, d'allèles et de locus.

\courssub{Matériel}
Fil de fer rigide ou cure-pipes (chenille) de \textbf{deux couleurs différentes} (ex. \textbf{rouge} = origine paternelle, \textbf{bleu} = origine maternelle), perles ou morceaux de pâte à modeler de couleurs distinctes (représentant \textbf{allèles différents} d'un même gène), étiquettes « Centromère », ciseaux.

\courssub{Protocole et Observations guidées}

\begin{enumerate}
    \item \textbf{Construction d'une paire homologue monochromatidienne (Phase G1 / Interphase) :}
    \begin{itemize}
        \item Prendre un fil \textbf{rouge} et un fil \textbf{bleu} de \textbf{même longueur} (même taille de chromosome).
        \item Placer une étiquette « Centromère » à la \textbf{même position} sur les deux (même morphologie).
        \item Fixer \textbf{3 perles} (gènes) aux \textbf{mêmes positions} (loci) sur les deux fils.
        \item \textbf{Choix des allèles :} Ex. Gène 1 : Rouge=Allèle $A$ (perle jaune), Bleu=Allèle $a$ (perle verte). Gène 2 : Rouge=$B$ (rouge), Bleu=$b$ (bleu). Gène 3 : Rouge=$C$ (blanc), Bleu=$C$ (blanc) $\rightarrow$ homozygote.
        \item \textbf{Observation :} Deux chromosomes \textbf{homologues} (même taille, même centromère, mêmes loci), mais \textbf{allèles différents} possible à chaque locus. Représente l'état \textbf{avant réplication} (2n, 2 molécules d'ADN par paire).
    \end{itemize}
    
    \item \textbf{Simulation de la réplication (Synthèse d'ADN - Phase S) :}
    \begin{itemize}
        \item Doubler chaque fil : créer une copie \textbf{identique} (même couleur, mêmes perles aux mêmes positions) pour le rouge $\rightarrow$ 2 fils rouges identiques. Idem pour le bleu $\rightarrow$ 2 fils bleus identiques.
        \item Assembler les deux copies rouges par leur centromère $\rightarrow$ \textbf{Chromosome bichromatidien paternel} (2 chromatides sœurs \textbf{identiques}).
        \item Assembler les deux copies bleues $\rightarrow$ \textbf{Chromosome bichromatidien maternel} (2 chromatides sœurs \textbf{identiques}).
        \item \textbf{Observation :} On a maintenant \textbf{2 chromosomes bichromatidiens homologues}. Chaque chromosome = 2 chromatides sœurs (identiques). La paire homologue = 4 chromatides au total (2 paires de sœurs). L'information est dupliquée fidèlement (complémentarité A=T, G=C).
    \end{itemize}
    
    \item \textbf{Application : Ségrégation à la méiose (Préfiguration programme suivant) :}
    \begin{itemize}
        \item \textbf{Méiose I (Réductionnelle) :} Séparation des chromosomes homologues (Rouge vs Bleu). Les chromatides sœurs restent jointes.
        \item \textbf{Méiose II (Équationnelle) :} Séparation des chromatides sœurs.
        \item Manipuler les maquettes pour montrer qu'un gamète reçoit \textbf{soit} le chromosome rouge (avec ses 2 chromatides), \textbf{soit} le bleu. Puis après Méiose II, une seule chromatide (une molécule d'ADN) par gène.
    \end{itemize}
\end{enumerate}

\begin{exemplebox}[Modélisation d'un locus : Groupe sanguin ABO]
Sur le chromosome 9 (autosome), le gène ABO a 3 allèles principaux : $I^A$, $I^B$, $i$.
\begin{itemize}
    \item Chromosome paternel (Rouge) : Allèle $I^A$ (perle rouge) au locus ABO.
    \item Chromosome maternel (Bleu) : Allèle $i$ (perle blanche) au locus ABO.
    \item Individu : Génotype $[I^A // i]$ $\rightarrow$ Phénotype Groupe A.
    \item Après réplication : 2 chromatides sœurs rouges ($I^A$), 2 chromatides sœurs bleues ($i$).
    \item Gamètes possibles : $I^A$ ou $i$ (ségrégation).
\end{itemize}
\end{exemplebox}

\courssub{Synthèse et ouverture}

Nous avons maintenant tous les éléments pour comprendre la structure fine du support de l'hérédité :
\begin{enumerate}
    \item L'information est dans le \textbf{noyau}, sur les \textbf{chromosomes} (prouvé par transfert de noyau).
    \item Le chromosome est fait de \textbf{chromatine} (ADN + protéines histones).
    \item L'\textbf{ADN} est une \textbf{double hélice} de deux brins antiparallèles.
    \item L'unité est le \textbf{nucléotide} (Désoxyribose + Phosphate + Base A/T/G/C).
    \item L'appariement \textbf{A=T et G$\equiv$C} (règles de Chargaff) assure la \textbf{fidélité de la réplication} (chaque brin sert de matrice) et la \textbf{stabilité} de l'information.
    \item L'information = \textbf{séquence des bases}. Un \textbf{gène} = un segment d'ADN codant un caractère (protéine).
    \item Les chromosomes existent par \textbf{paires homologues} (paternel/maternel), monochromatidiens (G1) ou bichromatidiens (G2/M).
\end{enumerate}

\noindent Ces acquis permettent d'aborder le cycle cellulaire, la mitose (conservation du patrimoine) et la méiose (brassage et réduction) au chapitre suivant.

\begin{exercice}[Entraînement complet : Unité 2 - Information génétique]
\begin{enumerate}
    \item \textbf{Transfert de noyau :} Résumer l'expérience de Hammerling sur \emph{Acetabularia} et conclure sur la localisation de l'information génétique.
    \item \textbf{Structure chromosome :} Différencier chromosome monochromatidien et bichromatidien. Définir chromatides sœurs et chromosomes homologues.
    \item \textbf{ADN - Calcul Chargaff :} Un ADN d'une plante du Sahel camerounais contient 24\% d'Adénine. Calculer \%T, \%G, \%C et le rapport (A+T)/(G+C).
    \item \textbf{Complémentarité :} Donner le brin complémentaire (5' $\rightarrow$ 3') de : 5' - T A C G G A T - 3'.
    \item \textbf{TP Feulgen :} Pourquoi l'hydrolyse acide à 60°C est-elle une étape cruciale ? Que se passe-t-il si on traite par la DNase avant coloration ?
    \item \textbf{Maquettes TP2 :} Sur votre maquette de paire homologue bichromatidienne, combien de molécules d'ADN au total ? Combien de chromatides sœurs par chromosome ? Les chromatides sœurs sont-elles identiques ? Les chromosomes homologues sont-ils identiques ?
    \item \textbf{Synthèse :} Compléter : « L'information génétique est portée par la molécule d'\underline{\hspace{2cm}}, située dans le \underline{\hspace{2cm}}, organisée en \underline{\hspace{2cm}}. L'unité de base est le \underline{\hspace{2cm}}. L'appariement des bases suit les règles de \underline{\hspace{2cm}} : A avec \underline{\hspace{1cm}}, G avec \underline{\hspace{1cm}}. Les deux brins sont \underline{\hspace{2cm}}. »
\end{enumerate}
\end{exercice}

\begin{corrige}[Exercice : Unité 2 - Information génétique]
\begin{enumerate}
    \item \textbf{Hammerling :} Greffe de pied (noyau) d'une espèce sur tronçon apical d'une autre. Le calice régénéré a la forme de l'espèce donneuse du \textbf{noyau}. Conclusion : L'information génétique est dans le \textbf{noyau}.
    \item \textbf{Monochromatidien :} 1 chromatide = 1 molécule ADN. \textbf{Bichromatidien :} 2 chromatides sœurs = 2 molécules ADN identiques (issues réplication). \textbf{Chromatides sœurs :} Identiques, même chromosome. \textbf{Homologues :} Pareils taille/forme/loci, origine père/mère, allèles éventuellement différents.
    \item \textbf{Chargaff :} \%A = 24\% $\Rightarrow$ \%T = 24\%. A+T = 48\%. G+C = 52\% $\Rightarrow$ \%G = \%C = 26\%. Rapport (A+T)/(G+C) = 48/52 = 0,92.
    \item \textbf{Complémentaire :} Donné 5'-TACGGAT-3'. Appariement : T$\leftrightarrow$A, A$\leftrightarrow$T, C$\leftrightarrow$G, G$\leftrightarrow$C, G$\leftrightarrow$C, A$\leftrightarrow$T, T$\leftrightarrow$A. Brin opposé 3'-ATGCCTA-5'. Écrit 5'$\rightarrow$3' : \textbf{5' - A T C C G T A - 3'}.
    \item \textbf{Feulgen :} Hydrolyse HCl 60°C = dépurination (création groupements aldéhydes sur désoxyribose) spécifiques de l'ADN (ARN hydrolysé). DNase avant Feulgen $\rightarrow$ destruction ADN $\rightarrow$ \textbf{pas de coloration rose} (témoin négatif prouvant spécificité).
    \item \textbf{Maquettes :} Paire homologue bichromatidienne = 2 chromosomes $\times$ 2 chromatides = \textbf{4 molécules d'ADN}. 2 chromatides sœurs par chromosome. \textbf{Sœurs = identiques}. \textbf{Homologues = non nécessairement identiques} (allèles différents possibles).
    \item \textbf{Synthèse :} ADN / noyau / chromosomes / nucléotide / Chargaff / T / C / antiparallèles.
\end{enumerate}
\end{corrige}

\coursec{TP N°1 : Mise en évidence du matériel génétique d'une cellule}

\courssub{Transition depuis la section précédente}

Nous venons d'établir, grâce aux expériences de transfert de noyau et à l'étude des chromosomes, que l'information génétique est logée dans le noyau des cellules, portée par des filaments d'ADN organisés en chromosomes. Mais l'ADN est une molécule microscopique, invisible à l'œil nu. Comment un biologiste peut-il \emph{voir} et \emph{manipuler} cette molécule pour l'étudier ? C'est tout l'objet de ce premier travail pratique : extraire l'ADN d'un tissu vivant pour le rendre visible macroscopiquement. Ce TP met en œuvre une démarche d'investigation scientifique complète, depuis la formulation d'un problème jusqu'à la rédaction d'un compte rendu rigoureux.

\courssub{Problématique et hypothèses}

\begin{definition}[Problème scientifique]
Comment mettre en évidence la présence d'ADN (matériel génétique) au sein des cellules d'un organisme vivant ?
\end{definition}

\begin{exemplebox}[Situation concrète : l'ADN dans votre assiette]
Chaque cellule de la banane que vous mangez au petit-déjeuner, de l'oignon que vous émincez pour la sauce, ou de votre propre salive contient une copie complète du génome de l'organisme. Si l'on pouvait briser les parois cellulaires et les membranes, puis dissoudre les protéines qui enrobent l'ADN, cette molécule longue de plusieurs mètres (si on la dépliait) se retrouverait libre dans le milieu. Mais elle est soluble dans l'eau... Comment la faire « sortir » de solution pour la voir ?
\end{exemplebox}

\begin{propriete}[Hypothèses de travail]
\begin{enumerate}
    \item Les cellules sont entourées d'une membrane (et d'une paroi pour les végétaux) qu'il faut rompre (lyse cellulaire).
    \item L'ADN est associé à des protéines (histones) dans le noyau ; il faut les dénaturer ou les éliminer.
    \item L'ADN est soluble dans l'eau (grâce à ses groupes phosphate chargés négativement) mais \textbf{insoluble dans l'alcool froid}. Cette propriété physique permet sa précipitation.
\end{enumerate}
\end{propriete}

\courssub{Matériel nécessaire}

\begin{center}
\begin{tabularx}{\linewidth}{|l|X|}\hline
\textbf{Matériel biologique} & Écrasé d'oignon cru, de banane mûre ou prélèvement de salive (cellules buccales) \\ \hline
\textbf{Produits chimiques} & Eau distillée, sel de cuisine (NaCl), liquide vaisselle ou détergent (SDS), alcool éthylique à 96\% \textbf{glacé} (conservé au congélateur au moins 1h) \\ \hline
\textbf{Verreirie \& matériel} & Mortier et pilon (ou sac de congélation + rouleau), bécher, tamis ou gaze stérile / filtre à café, baguette de verre, tubes à essai, support de tubes, gant de protection \\ \hline
\end{tabularx}
\end{center}

\begin{attention}[Sécurité au laboratoire]
\begin{itemize}
    \item L'alcool à 96\% est \textbf{inflammable} et \textbf{irritant} : pas de flamme à proximité, porter des lunettes de protection.
    \item Le détergent est irritant pour les yeux ; se laver les mains après manipulation.
    \item Ne \textbf{jamais} goûter les préparations, même si l'oignon et la banane sont comestibles au départ.
\end{itemize}
\end{attention}

\courssub{Protocole expérimental : démarche pas à pas}

\begin{methode}[Protocole d'extraction de l'ADN (méthode au détergent / alcool glacé)]
\underline{Étape 1 : Broyage et lyse mécanique (Rupture des parois et membranes)}
\begin{enumerate}
    \item Placer \SI{10}{g} de tissu végétal (oignon ou banane) dans le mortier.
    \item Ajouter \SI{10}{mL} d'eau distillée et \SI{1}{g} de sel (NaCl). Le sel neutralise les charges négatives des phosphates de l'ADN, favorisant son agrégation future, et aide à précipiter les protéines.
    \item Broyer énergiquement pendant 2 à 3 min jusqu'à obtenir une bouillie homogène (purée). \textit{Justification : rupture mécanique des parois cellulaires (cellulose) et des membranes plasmatiques (lipides).}
\end{enumerate}

\underline{Étape 2 : Lyse chimique par le détergent (Libération de l'ADN dans le cytoplasme)}
\begin{enumerate}
    \item Transvaser la purée dans un bécher. Ajouter \SI{5}{mL} de liquide vaisselle (détergent). Mélanger délicatement (éviter la mousse excessive) pendant 2 min.
    \item \textit{Justification : le détergent dissout les bicouches phospholipidiques des membranes (nucléaire et plasmique) par action tensioactive, libérant l'ADN dans le milieu. Il dénature aussi une partie des protéines.}
\end{enumerate}

\underline{Étape 3 : Filtration (Séparation du surnageant riche en ADN des débris cellulaires)}
\begin{enumerate}
    \item Filtrer le mélange à travers une gaze ou un filtre à café placé sur un bécher propre.
    \item Recueillir le filtrat (liquide jaunâtre/blanchâtre). Jeter le résidu solide (parois, organites gros, protéines précipitées).
    \item \textit{Justification : élimination des débris cellulaires insolubles (parois, membranes, grosses protéines) pour ne garder que la fraction soluble contenant l'ADN.}
\end{enumerate}

\underline{Étape 4 : Précipitation de l'ADN par l'alcool glacé (Révélation macroscopique)}
\begin{enumerate}
    \item Verser \SI{3}{mL} de filtrat dans un tube à essai.
    \item Incliner le tube à \ang{45}. Verser \textbf{très lentement} \SI{3}{mL} d'alcool éthylique \textbf{glacé} le long de la paroi interne pour former une couche distincte au-dessus du filtrat (ne pas mélanger !).
    \item Laisser reposer 2 à 3 min sans bouger, à la verticale.
    \item \textit{Observation : apparition de filaments blanchâtres, mousseux, à l'interface entre les deux liquides.}
    \item On peut enrouler ces filaments sur une baguette de verre tournée délicatement (spoolage).
\end{enumerate}
\end{methode}

\begin{popfigure}
\begin{tikzpicture}[
    node distance=1.2cm and 1.5cm,
    box/.style={draw=popdark, thick, rounded corners, fill=popcream, minimum width=3.5cm, minimum height=1.8cm, align=center, font=\small},
    arrow/.style={-Stealth, thick, popblue},
    labelbox/.style={draw=poporange, thick, rounded corners, fill=poporangeL, minimum width=4cm, minimum height=1.5cm, align=center, font=\footnotesize, text=popdark}
]
% Étape 1
\node[box, align=center] (e1){\textbf{Étape 1 : Broyage}\\\SI{10}{g} tissu + \SI{10}{mL} eau + \SI{1}{g} sel\\\textit{Purée homogène}};
\node[labelbox, right=of e1, align=center] (j1){Rupture parois \& membranes\\\textit{(Lyse mécanique)}};

% Étape 2
\node[box, below=of e1, align=center] (e2){\textbf{Étape 2 : Détergent}\\\textit{Ajout \SI{5}{mL} liquide vaisselle}\\\textit{Mélange doux 2 min}};
\node[labelbox, right=of e2, align=center] (j2){Dissolution membranes\\\textit{(Lyse chimique)}\\\textit{Libération ADN + dénaturation protéines}};

% Étape 3
\node[box, below=of e2, align=center] (e3){\textbf{Étape 3 : Filtration}\\\textit{Gaze / Filtre café}\\\textit{Récupérer filtrat}};
\node[labelbox, right=of e3, align=center] (j3){Séparation phases\\\textit{Élimination débris insolubles}};

% Étape 4
\node[box, below=of e3, align=center] (e4){\textbf{Étape 4 : Alcool glacé}\\\textit{Incliner tube, verser alcool}\\\textit{\textit{le long paroi}}\\\textit{\textit{Repos 2-3 min}}};
\node[labelbox, right=of e4, align=center] (j4){Précipitation ADN\\\textit{Insoluble dans alcool froid}\\\textit{Filaments visibles à l'interface}};

% Flèches
\draw[arrow] (e1.south) -- (e2.north);
\draw[arrow] (e2.south) -- (e3.north);
\draw[arrow] (e3.south) -- (e4.north);

% Titre global
\node[above=0.3cm of e1, font=\bfseries\large, text=popblue] {Protocole d'extraction de l'ADN (4 étapes)};
\end{tikzpicture}
\legende{Schéma récapitulatif du protocole d'extraction de l'ADN. Chaque étape est justifiée par son rôle biophysique (lyse mécanique, lyse chimique, clarification, précipitation).}
\end{popfigure}

\courssub{Observations et résultats attendus}

\begin{popfigure}
\begin{tikzpicture}[
    scale=1.2,
    tube/.style={draw=popdark, very thick, fill=popcream},
    liquid/.style={fill=popblue!20},
    alcohol/.style={fill=poporange!15},
    dna/.style={draw=poppurple, very thick, fill=poppurple!30, decorate, decoration={snake, amplitude=1pt, segment length=4pt}},
    label/.style={font=\small, text=popdark, anchor=west}
]
% Tube à essai
\draw[tube] (0,0) rectangle (3,6);
\draw[tube] (0,6) .. controls (0,6.5) and (3,6.5) .. (3,6); % haut arrondi
\draw[tube] (0,0) .. controls (0,-0.5) and (3,-0.5) .. (3,0); % bas arrondi

% Niveaux liquides
\fill[liquid] (0.15,0.15) rectangle (2.85,3.5); % Filtrat (bas)
\fill[alcohol] (0.15,3.5) rectangle (2.85,5.5); % Alcool (haut)

% Interface
\draw[dashed, thick, popdark] (0.15,3.5) -- (2.85,3.5);
\node[label] at (3.1, 3.5) {Interface alcool / filtrat};

% Filaments d'ADN à l'interface
\draw[dna] (1.0, 3.3) -- (1.0, 3.7);
\draw[dna] (1.5, 3.2) -- (1.5, 3.8);
\draw[dna] (2.0, 3.4) -- (2.0, 3.6);
\draw[dna] (0.7, 3.35) -- (0.7, 3.75);
\draw[dna] (2.3, 3.3) -- (2.3, 3.7);

% Baguette de verre (spoolage)
\draw[draw=popmuted, thick] (1.5, 5.8) -- (1.5, 3.6);
\draw[draw=popmuted, thick, fill=popmuted!20] (1.3, 5.8) rectangle (1.7, 6.2);
\node[label, rotate=90] at (3.2, 4.5) {Filaments d'ADN (blanchâtres)};

% Légendes zones
\node[label, align=left] at (3.1, 4.5) {Phase supérieure :\\Alcool éthylique \textbf{glacé} (96\%)};
\node[label, align=left] at (3.1, 1.8) {Phase inférieure :\\Filtrat (eau, sel, ADN,\\détergent, protéines solubles)};
\node[label] at (3.1, 0.5) {Tube à essai};
\end{tikzpicture}
\legende{Observation macroscopique dans le tube à essai après l'étape 4. L'ADN, insoluble dans l'alcool froid, précipite sous forme de filaments blanchâtres (mucus) exactement à l'interface entre les deux liquides non miscibles. On peut les récupérer à la baguette de verre (spoolage).}
\end{popfigure}

\begin{aretenir}[Résultat macroscopique]
L'expérience met en évidence la présence de \textbf{filaments blanchâtres, visqueux et élastiques} à l'interface entre le filtrat aqueux et l'alcool glacé. Ces filaments correspondent à de l'\textbf{ADN décondensé} (chromatine) qui s'agrège par milliers de brins pour devenir visible à l'œil nu. C'est la preuve expérimentale que les cellules contiennent une substance macromoléculaire aux propriétés physico-chimiques de l'acide désoxyribonucléique.
\end{aretenir}

\courssub{Interprétation rigoureuse : justification de chaque étape}

\begin{experience}[Analyse biophysique du protocole]
\textbf{Objectif :} Justifier scientifiquement le rôle de chaque réactif et de chaque geste technique.

\begin{center}
\begin{tabularx}{\linewidth}{|l|X|X|}\hline
\rowcolor{popblueL}
\textbf{Étape / Réactif} & \textbf{Action biophysique} & \textbf{Conséquence pour l'ADN} \\ \hline
\textbf{Sel (NaCl)} & Les ions $Na^+$ neutralisent les charges négatives ($PO_4^-$) du squelette phosphodiester de l'ADN. & Réduit la répulsion électrostatique entre brins d'ADN $\rightarrow$ favorise l'agrégation (précipitation) future. Précipite aussi certaines protéines. \\ \hline
\textbf{Détergent (Liquide vaisselle / SDS)} & Molécules amphiphiles : têtes polaires (eau), queues hydrophobes (lipides). Insertion dans les bicouches phospholipidiques. & \textbf{Lyse membranaire} : rupture membrane plasmique et enveloppe nucléaire. Libération de l'ADN dans le milieu. Dénaturation partielle des protéines (histones). \\ \hline
\textbf{Broyage (Mécanique)} & Force de cisaillement. & Rupture de la \textbf{paroi cellulosique} (végétaux) inaccessible au détergent seul. Augmente la surface de contact pour la lyse chimique. \\ \hline
\textbf{Filtration} & Tamisage physique (pores gaze/filtre). & Séparation : \textit{Résidu} = parois, membranes, gros organites, protéines précipitées. \textit{Filtrat} = ADN soluble, sels, détergent, petites protéines. \\ \hline
\textbf{Alcool éthylique \textbf{GLACÉ} (96\%)} & Diminue la constante diélectrique du milieu. L'ADN est insoluble dans l'éthanol $> 70\%$, surtout à basse température ($< 4^\circ\mathrm{C}$). & \textbf{Précipitation sélective} : L'ADN « sort » de solution (devient visible). Les sels, détergents et la plupart des protéines restent solubles dans l'alcool. Le froid ralentit l'action des DNases (enzymes dégradant l'ADN) résiduelles. \\ \hline
\textbf{Verser le long de la paroi (inclinaison)} & Évite la turbulence et le mélange brutal des deux phases. & Maintient une \textbf{interface nette}. L'ADN précipite \emph{uniquement} à l'interface où la concentration en alcool devient critique, formant des filaments continus faciles à récupérer. \\ \hline
\end{tabularx}
\end{center}
\end{experience}

\begin{attention}[Pièges fréquents et erreurs à éviter au BEPC]
\begin{itemize}
    \item \textbf{Alcool non glacé :} L'ADN reste soluble ou précipite en « poudre » fine invisible, pas en filaments. \textit{Consigne : sortir l'alcool du congélateur au dernier moment.}
    \item \textbf{Verser l'alcool brutalement / mélanger :} Les deux phases se mélangent $\rightarrow$ précipitation diffuse, pas de filaments à l'interface, perte de rendement.
    \item \textbf{Trop de mousse au détergent :} Piège l'ADN, rend la filtration lente et le filtrat trouble. \textit{Astuce : mélanger délicatement, laisser reposer 1 min avant filtration.}
    \item \textbf{Confondre résidu et filtrat :} L'ADN est dans le \textbf{filtrat} (liquide qui passe), pas dans le résidu (débris solides sur le filtre).
    \item \textbf{Rédaction :} Ne pas écrire « on voit l'ADN ». Écrire : « on observe des filaments blanchâtres \textbf{correspondant à de l'ADN précipité} ». L'ADN pur est incolore ; l'aspect blanchâtre vient de l'agrégation massive et des impuretés protéiques résiduelles.
\end{itemize}
\end{attention}

\courssub{Rédaction du compte rendu de TP (Modèle BEPC)}

Un compte rendu scientifique suit une structure immuable. Voici le canevas attendu à l'examen, illustré par notre expérience.

\begin{exoresolu}[Compte rendu type : Extraction de l'ADN d'oignon]
\textbf{1. Titre} : Mise en évidence du matériel génétique (ADN) dans les cellules d'oignon.

\textbf{2. Problème} : Comment rendre visible l'ADN contenu dans les cellules végétales ?

\textbf{3. Hypothèses} :
\begin{itemize}
    \item H1 : La rupture des parois et membranes libère l'ADN dans le milieu extracellulaire.
    \item H2 : L'ADN, soluble dans l'eau, devient insoluble dans l'alcool éthylique froid, ce qui permet sa précipitation macroscopique.
\end{itemize}

\textbf{4. Matériel et Méthode} : (Voir protocole détaillé ci-dessus. Citer les volumes/masses exactes).

\textbf{5. Résultats (Observations)} :
\begin{itemize}
    \item Après broyage : obtention d'une purée homogène.
    \item Après ajout détergent : mélange mousseux, aspect laiteux.
    \item Après filtration : filtrat limpide, jaunâtre ; résidu solide sur le filtre.
    \item \textbf{Observation clé :} Après ajout lent d'alcool glacé et repos de 2 min, \textbf{apparition de filaments blanchâtres, visqueux, à l'interface des deux liquides}. Ces filaments s'enroulent sur une baguette de verre.
\end{itemize}

\textbf{6. Interprétation} :
\begin{itemize}
    \item Le broyage + sel + détergent $\rightarrow$ lyse complète (paroi + membranes) $\rightarrow$ libération de l'ADN nucléaire dans le filtrat.
    \item L'alcool glacé $\rightarrow$ précipitation sélective de l'ADN (insolubilité dans l'éthanol froid) sous forme de filaments agrégés.
    \item Les filaments observés sont bien du matériel génétique car : (1) ils proviennent du noyau (lyse nucléaire), (2) ils ont les propriétés physico-chimiques de l'ADN (précipitation alcool froid), (3) l'expérience de transfert de noyau (leçon précédente) a prouvé que le noyau contient l'information héréditaire.
\end{itemize}

\textbf{7. Conclusion} :
L'expérience permet de \textbf{mettre en évidence macroscopiquement la présence d'ADN dans les cellules d'oignon}. Le protocole de lyse (mécanique + chimique) suivi de précipitation à l'alcool glacé est une méthode efficace pour extraire le matériel génétique. Cela confirme que l'information génétique est une substance chimique tangible, localisée dans le noyau, que l'on peut isoler et manipuler.
\tcblower
\textbf{Analyse du correcteur (Points de vigilance BEPC) :}
\begin{itemize}
    \item \checkmark Problème et Hypothèses \textbf{distincts} et formulés avant le protocole.
    \item \checkmark Justification de \textbf{chaque} réactif (Sel, Détergent, Alcool glacé) dans l'interprétation.
    \item \checkmark Distinction claire \textbf{Résultats (ce qu'on voit)} vs \textbf{Interprétation (ce qu'on en déduit)}.
    \item \checkmark Conclusion répondant \textbf{exactement} au problème posé.
    \item \checkmark Vocabulaire précis : « filaments blanchâtres », « interface », « précipitation », « lyse », « filtrat/résidu ».
\end{itemize}
\end{exoresolu}

\courssub{Exercices d'application et consolidation}

\begin{exercice}[Analyse critique de protocoles variants]
Un élève propose trois modifications du protocole pour « aller plus vite ». Prédisez le résultat et justifiez.
\begin{enumerate}
    \item Il remplace l'alcool glacé par de l'alcool à température ambiante.
    \item Il omet le sel (NaCl) dans l'eau de broyage.
    \item Il verse l'alcool brutalement au centre du tube sans incliner.
\end{enumerate}
\end{exercice}

\begin{corrige}[Exercice : Analyse critique]
\begin{enumerate}
    \item \textbf{Alcool tiède :} Précipitation faible ou nulle. L'ADN reste soluble dans l'éthanol chaud. Les filaments seront absents ou sous forme de poudre fine difficile à voir. \textit{Justification : la solubilité de l'ADN dans l'éthanol augmente avec la température.}
    \item \textbf{Sans sel :} L'ADN précipite moins bien (répulsion charges négatives non neutralisées), filaments plus fins, plus fragiles, rendement moindre. Les protéines contaminantes précipitent moins $\rightarrow$ filtrat plus trouble. \textit{Justification : rôle écran des ions $Na^+$ sur les phosphates.}
    \item \textbf{Versement brutal :} Mélange des phases $\rightarrow$ précipitation diffuse dans tout le volume, pas d'interface nette $\rightarrow$ pas de filaments continus récupérables à la baguette. Aspect « neige » ou gelée trouble. \textit{Justification : nécessité d'un gradient de concentration d'alcool à l'interface pour une précipitation progressive et fibreuse.}
\end{enumerate}
\end{corrige}

\begin{exercice}[Transfert de compétences : Cas de la salive humaine]
On souhaite extraire l'ADN de cellules buccales (prélèvement à l'écouvillon dans de l'eau salée).
\begin{enumerate}
    \item L'étape de broyage au mortier est-elle nécessaire ? Justifiez.
    \item Le détergent est-il toujours nécessaire ? Justifiez.
    \item L'alcool glacé est-il toujours nécessaire ? Justifiez.
\end{enumerate}
\end{exercice}

\begin{corrige}[Exercice : Salive humaine]
\begin{enumerate}
    \item \textbf{Non.} Les cellules buccales (épithélium) n'ont \textbf{pas de paroi cellulosique}, seulement une membrane plasmique. Le broyage mécanique est inutile (et casserait l'ADN par cisaillement excessif). Un simple mélange par inversion ou agitation douce suffit.
    \item \textbf{Oui.} La membrane plasmique et l'enveloppe nucléaire sont composées de phospholipides. Le détergent est indispensable pour les lyser (lyse chimique).
    \item \textbf{Oui.} La propriété d'insolubilité de l'ADN dans l'alcool froid est une propriété intrinsèque de la molécule, indépendante de l'espèce source. C'est l'étape de révélation universelle.
\end{enumerate}
\end{corrige}

\courssub{Le savais-tu ? : De la classe au laboratoire criminel}

\begin{savaistu}[L'extraction d'ADN : même principe, des enjeux majeurs]
La manipulation que vous venez de réaliser en classe (broyage $\rightarrow$ lyse $\rightarrow$ précipitation alcool) est \textbf{la toute première étape} de \textbf{toutes} les analyses d'ADN modernes :
\begin{itemize}
    \item \textbf{Tests de paternité / maternité :} Comparaison de profils génétiques (empreintes génétiques) entre enfant et père présumé.
    \item \textbf{Identification criminelle (Police scientifique / INTERPOL) :} ADN prélevé sur une trace biologique (cheveu, salive, sang, sperme) sur une scène de crime $\rightarrow$ extraction $\rightarrow$ amplification (PCR) $\rightarrow$ profilage $\rightarrow$ comparaison base de données (FNAEG en France, bases nationales au Cameroun).
    \item \textbf{Diagnostic médical :} Dépistage de mutations (drépanocytose, mucoviscidose, cancers héréditaires) sur prélèvement sanguin.
    \item \textbf{Biodiversité \& Braconnage :} Au Cameroun, analyse d'ADN d'écailles de pangolin, d'ivoire ou de bois précieux saisis pour tracer l'origine géographique et les réseaux de trafic.
    \item \textbf{Agriculture :} Sélection variétale assistée par marqueurs (MAS) pour le cacaoyer, le cotonnier, le manioc (résistance aux virus, rendement).
\end{itemize}
La différence avec le TP ? Les laboratoires utilisent des colonnes à silice ou des billes magnétiques pour purifier l'ADN (enlever protéines, ARN, détergent, inhibiteurs) et obtenir un ADN \textbf{pur, quantifié (ng/µL), intact} (haut poids moléculaire), indispensable pour la PCR (amplification enzymatique). Mais le \textbf{principe physico-chimique de base reste identique} : lyse cellulaire + précipitation sélective de l'acide nucléique.
\end{savaistu}

\courssub{Synthèse des savoir-faire évalués (BEPC)}

\begin{aretenir}[Check-list des compétences TP N°1]
À l'issue de ce TP, vous devez être capable de :
\begin{enumerate}
    \item \textbf{Formuler} un problème et des hypothèses testables à partir d'une question biologique.
    \item \textbf{Nommer} le matériel et les réactifs avec leurs quantités (unités SI).
    \item \textbf{Rédiger} un protocole opératoire logique, numéroté, avec les gestes techniques critiques (inclinaison, lenteur, alcool glacé).
    \item \textbf{Justifier} scientifiquement \textbf{chaque} réactif et \textbf{chaque} étape (rôle biophysique : lyse, dénaturation, précipitation).
    \item \textbf{Distinguer} l'observation brute (filaments blanchâtres à l'interface) de l'interprétation (précipitation de l'ADN).
    \item \textbf{Conclure} en répondant au problème initial, en lien avec la localisation nucléaire de l'information génétique.
    \item \textbf{Transposer} le protocole à d'autres sources cellulaires (salive, banane, foie, levure) en adaptant l'étape de lyse (paroi ou pas).
\end{enumerate}
\end{aretenir}

\begin{exemplebox}[Astuce mnémotechnique pour l'oral ou l'écrit : « S.D.F.A »]
Pour ne rien oublier dans la justification, retenez l'acronyme \textbf{S.D.F.A} (comme \textbf{S}ans \textbf{D}éfaut \textbf{F}ait \textbf{A}fficher l'ADN) :
\begin{itemize}
    \item \textbf{S}el $\rightarrow$ Neutralise charges (agrégation future).
    \item \textbf{D}étergent $\rightarrow$ Dissout membranes (lyse).
    \item \textbf{F}iltration $\rightarrow$ Filtre débris (clarification).
    \item \textbf{A}lcool \textbf{G}lacé $\rightarrow$ Précipite ADN (révélation).
\end{itemize}
\end{exemplebox}

\courssub{Passerelle vers le TP N°2}

Nous venons d'extraire l'ADN « en vrac » : un mélange de millions de fragments de chromosomes de toutes tailles. Mais le programme nous dit que l'information est organisée en \textbf{chromosomes}, structures discrètes, présentes en \textbf{paires homologues} chez l'individu diploïde. Comment passer de la « soupe » d'ADN extraite ici à la compréhension de l'organisation structurée du génome en paires de chromosomes monochromatidiens et bichromatidiens ? C'est l'objet du \textbf{TP N°2 : Conception des maquettes des paires de chromosomes homologues}.

\coursec{TP N°2 : Maquettes de paires de chromosomes homologues et applications}

Après avoir mis en évidence la présence d'ADN dans le noyau lors du TP N°1, nous allons maintenant \textbf{modéliser l'organisation de ce matériel génétique} à l'échelle des chromosomes. Comprendre la différence entre chromosome monochromatidien et bichromatidien, et entre chromatides sœurs et chromosomes homologues, est indispensable pour aborder sereinement la méiose et la génétique mendélienne au programme du BEPC.

\courssub{Objectif et matériel}

\begin{experience}[TP N°2 : Construction de maquettes de chromosomes homologues]
\textbf{Objectif :} Réaliser des maquettes physiques pour visualiser la structure d'une paire de chromosomes homologues à l'état monochromatidien et bichromatidien, et en déduire le nombre de chromosomes et de chromatides.

\textbf{Matériel par binôme :}
\begin{itemize}
    \item Ficelles ou fils de laine de \textbf{deux couleurs distinctes} (ex. : rouge et bleu) — une couleur = origine paternelle, l'autre = origine maternelle.
    \item Pâte à modeler ou papier cartonné de ces deux couleurs (pour le centromère).
    \item Ciseaux, règle, feutres.
    \item Feuilles A3 pour le montage final et l'annotation.
\end{itemize}

\textbf{Consignes de sécurité et de propreté :} Manipuler les ciseaux avec précaution. Ranger le matériel après usage. Se laver les mains après manipulation de la pâte à modeler.
\end{experience}

\begin{methode}[Construire une maquette correcte de chromosomes homologues]
Pour que ta maquette soit scientifiquement valide et te rapporte les points au BEPC, respecte scrupuleusement ces règles :
\begin{enumerate}
    \item \textbf{Code couleur immuable :} Choisis une couleur pour le chromosome d'origine \textbf{paternelle} (ex. rouge) et une autre pour l'origine \textbf{maternelle} (ex. bleu). Ne change pas de couleur en cours de TP.
    \item \textbf{Taille et forme identiques pour les homologues :} Les deux chromosomes d'une paire homologue portent les mêmes gènes aux mêmes loci. Ils doivent avoir la \textbf{même longueur} et le \textbf{centromère à la même position} (métacentrique, sous-métacentrique, acrocentrique).
    \item \textbf{Centromère unique par chromosome :} Matérialise-le par une boule de pâte à modeler ou un pli marqué. C'est le point d'attache des fibres du fuseau.
    \item \textbf{Chromatides sœurs = clones :} Lors de la duplication (étape 2), les deux chromatides d'un même chromosome sont \textbf{identiques} (même couleur, même taille, mêmes gènes). Elles sont jointes par \textbf{le même centromère}.
    \item \textbf{Légende obligatoire :} Sur ton poster final, indique : « Paire d'homologues monochromatidiens (G1) » et « Paire d'homologues bichromatidiens (G2/M) », avec le nombre de chromosomes (2) et de chromatides (2 puis 4).
\end{enumerate}
\end{methode}

\courssub{Étape 1 : Paire de chromosomes homologues monochromatidiens (état G1)}

\emph{Rappel :} Au cours de la phase G1 du cycle cellulaire, avant la réplication de l'ADN (phase S), chaque chromosome ne comporte qu'\textbf{une seule chromatide}. On dit qu'il est \textbf{monochromatidien}.

\begin{enumerate}
    \item Coupe deux fils de laine de \textbf{longueur identique} (ex. \SI{30}{cm}), un rouge (père), un bleu (mère).
    \item Place une boule de pâte à modeler (ou un nœud serré) au même tiers de la longueur sur chaque fil pour représenter le \textbf{centromère}. Tu obtiens deux chromosomes de même morphologie.
    \item Colle-les parallèlement sur ta feuille A3, centromères alignés. Trace un trait d'union entre eux pour symboliser l'homologie.
    \item \textbf{Compte et note :}
    \begin{itemize}
        \item Nombre de \textbf{chromosomes} : 2 (un paternel, un maternel).
        \item Nombre de \textbf{chromatides} : 2 (une par chromosome).
    \end{itemize}
\end{enumerate}

\begin{popfigure}
\begin{tikzpicture}[scale=0.9, every node/.style={font=\small}]
    % Chromosome paternel (rouge) - monochromatidien
    \draw[popred, line width=3pt, rounded corners=2pt] (0,1.2) -- (0,0.2) -- (4,0.2) -- (4,1.2) -- cycle;
    \draw[popred, line width=3pt] (2,0.2) -- (2,-0.3);
    \fill[popred] (2,-0.3) circle (3pt); % centromere
    \node[popred, left] at (-0.5,0.7) {Paternel (Rouge)};
    
    % Chromosome maternel (bleu) - monochromatidien
    \draw[popblue, line width=3pt, rounded corners=2pt] (5,1.2) -- (5,0.2) -- (9,0.2) -- (9,1.2) -- cycle;
    \draw[popblue, line width=3pt] (7,0.2) -- (7,-0.3);
    \fill[popblue] (7,-0.3) circle (3pt); % centromere
    \node[popblue, right] at (9.5,0.7) {Maternel (Bleu)};
    
    % Brace d'homologie (dessin manuel sans librairie decorations.pathreplacing)
    \draw[popdark, thick] (0,1.4) -- (0,1.6) -- (4,1.6) -- (4,1.4);
    \draw[popdark, thick] (5,1.4) -- (5,1.6) -- (9,1.6) -- (9,1.4);
    \node[popdark, above] at (2,1.75) {Paire d'homologues};
    \node[popdark, above] at (7,1.75) {Paire d'homologues};
    
    % Annotations
    \node[align=center, popmuted] at (2, -1.2) {1 Chromosome \\ 1 Chromatide};
    \node[align=center, popmuted] at (7, -1.2) {1 Chromosome \\ 1 Chromatide};
    \node[align=center, popdark, font=\bfseries] at (4.5, -2) {Total : 2 Chromosomes \\ 2 Chromatides};
\end{tikzpicture}
\legende{Schéma d'une paire de chromosomes homologues \textbf{monochromatidiens} (phase G1). Le chromosome paternel (rouge) et le chromosome maternel (bleu) ont même taille et même position du centromère. Chacun ne possède qu'une seule chromatide.}
\end{popfigure}

\courssub{Étape 2 : Transformation en chromosomes bichromatidiens (état G2 / Mitose / Méiose I)}

\emph{Rappel :} Au cours de la phase S (synthèse), l'ADN se réplique. Chaque chromosome est alors constitué de \textbf{deux chromatides sœurs identiques}, réunies par leur centromère commun. C'est l'état \textbf{bichromatidien}.

\begin{enumerate}
    \item Pour \textbf{chaque} chromosome de l'étape 1, coupe un deuxième fil de la \textbf{même couleur} et de la \textbf{même longueur}.
    \item Superpose-le exactement sur le premier fil (même centromère). Fixe-les ensemble au niveau du centromère (pâte à modeler commune ou nœud unique serrant les deux brins).
    \item Écarte légèrement les deux bras (chromatides) pour bien voir la structure en « X » (si centromère médian) ou en « V/L ».
    \item Colle cette nouvelle structure à côté de la première maquette sur ton poster.
    \item \textbf{Compte et note :}
    \begin{itemize}
        \item Nombre de \textbf{chromosomes} : \textbf{2} (toujours un paternel, un maternel — le nombre de centromères n'a pas changé).
        \item Nombre de \textbf{chromatides} : \textbf{4} (deux sœurs rouges identiques, deux sœurs bleues identiques).
    \end{itemize}
\end{enumerate}

\begin{popfigure}
\begin{tikzpicture}[scale=0.9, every node/.style={font=\small}]
    % Chromosome paternel (rouge) - bichromatidien
    \draw[popred, line width=3pt, rounded corners=2pt] (0,1.5) -- (0,0.5) -- (1.8,0.5) -- (2,0) -- (2.2,0.5) -- (4,0.5) -- (4,1.5) -- cycle;
    \draw[popred, line width=3pt, rounded corners=2pt] (0,-0.5) -- (0,-1.5) -- (1.8,-0.5) -- (2,0) -- (2.2,-0.5) -- (4,-0.5) -- (4,-1.5) -- cycle;
    \fill[popred] (2,0) circle (4pt); % centromere commun
    \node[popred, left] at (-0.8,0) {Paternel};
    \node[popred, font=\tiny] at (1,1) {Chromatide sœur 1};
    \node[popred, font=\tiny] at (1,-1) {Chromatide sœur 2 (identique)};
    
    % Chromosome maternel (bleu) - bichromatidien
    \draw[popblue, line width=3pt, rounded corners=2pt] (5,1.5) -- (5,0.5) -- (6.8,0.5) -- (7,0) -- (7.2,0.5) -- (9,0.5) -- (9,1.5) -- cycle;
    \draw[popblue, line width=3pt, rounded corners=2pt] (5,-0.5) -- (5,-1.5) -- (6.8,-0.5) -- (7,0) -- (7.2,-0.5) -- (9,-0.5) -- (9,-1.5) -- cycle;
    \fill[popblue] (7,0) circle (4pt); % centromere commun
    \node[popblue, right] at (9.8,0) {Maternel};
    \node[popblue, font=\tiny] at (6,1) {Chromatide sœur 1};
    \node[popblue, font=\tiny] at (6,-1) {Chromatide sœur 2 (identique)};
    
    % Brace d'homologie (dessin manuel)
    \draw[popdark, thick] (0,-1.7) -- (0,-1.9) -- (4,-1.9) -- (4,-1.7);
    \draw[popdark, thick] (5,-1.7) -- (5,-1.9) -- (9,-1.9) -- (9,-1.7);
    \node[popdark, below] at (2,-2.1) {Paire d'homologues};
    \node[popdark, below] at (7,-2.1) {Paire d'homologues};
    
    % Annotations
    \node[align=center, popmuted] at (2, -2.5) {1 Chromosome \\ 2 Chromatides sœurs};
    \node[align=center, popmuted] at (7, -2.5) {1 Chromosome \\ 2 Chromatides sœurs};
    \node[align=center, popdark, font=\bfseries] at (4.5, -3.3) {Total : 2 Chromosomes \\ 4 Chromatides};
\end{tikzpicture}
\legende{Schéma d'une paire de chromosomes homologues \textbf{bichromatidiens} (phase G2, début mitose/méiose). Chaque chromosome (défini par un centromère) est formé de deux chromatides sœurs \textbf{identiques} (même couleur). La paire compte toujours 2 chromosomes mais 4 chromatides.}
\end{popfigure}

\courssub{Comparaison et synthèse : Mono- vs Bichromatidien}

C'est le cœur de la compréhension pour le BEPC. Ne confonds jamais \textbf{nombre de chromosomes} (compté par le nombre de centromères) et \textbf{nombre de chromatides} (compté par le nombre de brins d'ADN).

\begin{center}
\begin{tabularx}{\linewidth}{|l|X|X|}\hline
\rowcolor{popblueL}
\textbf{Critère} & \textbf{État monochromatidien (G1)} & \textbf{État bichromatidien (G2, M)} \\ \hline
\textbf{Aspect} & Bâtonnet simple (I, V, J) & Double bâtonnet (X, V double) \\ \hline
\textbf{Nombre de centromères / paire} & 2 & 2 (inchangé) \\ \hline
\textbf{Nombre de chromosomes / paire} & \textbf{2} (1 paternel + 1 maternel) & \textbf{2} (1 paternel + 1 maternel) \\ \hline
\textbf{Nombre de chromatides / paire} & \textbf{2} (1 par chromosome) & \textbf{4} (2 sœurs par chromosome) \\ \hline
\textbf{Identité génétique} & Homologues : gènes mêmes loci, allèles éventuellement différents & Sœurs : \textbf{identiques} (clones) ; Homologues : semblables (père $\neq$ mère) \\ \hline
\textbf{Moment dans le cycle} & Après mitose / Phase G1 & Après phase S (réplication) / Phases G2, M \\ \hline
\end{tabularx}
\end{center}

\begin{attention}[Piège majeur au BEPC : Chromosomes vs Chromatides]
\begin{itemize}
    \item \textbf{Deux chromatides sœurs} = issues de la réplication d'un \textbf{même} chromosome $\rightarrow$ \textbf{ADN identique} (sauf erreur rare de réplication). Elles se séparent en anaphase de mitose / anaphase II de méiose.
    \item \textbf{Deux chromosomes homologues} = un vient du \textbf{père}, l'autre de la \textbf{mère} $\rightarrow$ \textbf{ADN semblable} (mêmes gènes, mêmes loci) mais \textbf{allèles souvent différents}. Ils se séparent en anaphase I de méiose.
    \item \textbf{Erreur fréquente :} Dire « 4 chromosomes » à l'état bichromatidien. \textbf{FAUX}. Tant que les centromères ne se sont pas divisés, ce sont 2 chromosomes (chacun à 2 chromatides).
    \item \textbf{Mnémotechnique :} \textbf{Centromère = Compteur de Chromosomes}.
\end{itemize}
\end{attention}

\courssub{Applications concrètes : de la maquette à la vie courante}

La structure des chromosomes et l'unicité de l'ADN (sauf jumeaux monozygotes) sont exploitées dans trois domaines majeurs au Cameroun et dans le monde.

\begin{exemplebox}[Application 1 : Test de paternité par ADN (Laboratoire du Centre Pasteur du Cameroun, Yaoundé)]
\textbf{Contexte :} Un père présumé conteste la paternité d'un enfant. Le juge ordonne un test ADN.
\textbf{Principe :} L'enfant a reçu \textbf{un chromosome homologue de son père} (donc une chromatide du chromosome paternel bichromatidien lors de la méiose) et l'autre de sa mère. On compare des \textbf{marqueurs microsatellites (STR)} — séquences répétées non codantes très polymorphes — sur une dizaine de loci autosomiques.
\textbf{Résultat :} À chaque locus, l'enfant \textbf{doit posséder un allèle retrouvé chez le père présumé}. Si 3 loci ou plus excluent le père (allèle enfant absent chez l'homme), la paternité est exclue à $> 99,99\%$. Si tous les loci sont compatibles, la probabilité de paternité dépasse $99,99\%$.
\textbf{Lien avec le TP :} L'allèle transmis correspond à l'une des chromatides du chromosome bichromatidien paternel (après méiose). La maquette montre bien que le père ne transmet \textbf{qu'un seul} de ses deux chromosomes homologues (donc une seule version de chaque gène) à chaque gamète.
\end{exemplebox}

\begin{exemplebox}[Application 2 : Identification en médecine légale (Police scientifique, DGSN)]
\textbf{Contexte :} Trace biologique (sang, salive, cheveu, sperme) sur une scène de crime (ex. braquage à Douala, agression).
\textbf{Technique :} \textbf{Profilage génétique (empreinte génétique)}. On amplifie par PCR 16 à 23 loci STR + Amelogenin (sexe). Le profil obtenu (sorte de « code-barres » numérique) est unique pour un individu (probabilité $< 10^{-18}$).
\textbf{Base de données :} Le Cameroun développe son FNAEG (Fichier National Automatisé des Empreintes Génétiques) pour croiser profils de scènes de crime et profils de personnes mises en cause / condamnées.
\textbf{Lien avec le TP :} L'unicité du profil repose sur la combinaison allélique des \textbf{paires d'homologues} : à chaque locus, on a deux allèles (un sur le chromosome paternel, un sur le maternel). La maquette bichromatidienne rappelle que chaque chromosome porte un allèle, et que la méiose n'en garde qu'un par gamète.
\end{exemplebox}

\begin{exemplebox}[Application 3 : Diagnostic prénatal par caryotype (Hôpital Central / Gynéco-Obstétrique)]
\textbf{Contexte :} Femme enceinte de 38 ans (âge maternel avancé = risque accru de non-disjonction méiotique). Amniocentèse à 16 SA (semaines d'aménorrhée).
\textbf{Technique :} Culture de lymphocytes fœtaux $\rightarrow$ blocage en métaphase (chromosomes \textbf{bichromatidiens} très condensés, bien visibles) $\rightarrow$ coloration au Giemsa (bandes G) $\rightarrow$ photographie $\rightarrow$ \textbf{caryotypage} (classement par paires d'homologues par taille et position du centromère).
\textbf{Recherche :} Nombre de chromosomes (46 ?) + Structure des paires homologues.
\begin{itemize}
    \item \textbf{Trisomie 21 :} Trois chromosomes 21 au lieu d'une paire $\rightarrow$ 47 chromosomes. Excès de matériel génétique $\rightarrow$ syndrome de Down.
    \item \textbf{Translocation, délétion, inversion :} Anomalies structurelles visibles sur la maquette bichromatidienne (bras manquants, centromère déplacé).
\end{itemize}
\textbf{Lien avec le TP :} Le caryotype \textbf{ne s'obtient que sur chromosomes bichromatidiens} (métaphase). Ta maquette de l'étape 2 est exactement ce que le biologiste observe au microscope pour compter et apparier les homologues.
\end{exemplebox}

\begin{savaistu}[Empreintes génétiques et police scientifique au Cameroun]
La première condamnation au Cameroun grâce à l'ADN remonte à 2012 (affaire de viol à Yaoundé, expertise du Centre Pasteur). Aujourd'hui, la DGSN et la Gendarmerie disposent de laboratoires mobiles de prélèvement. La loi n°2016/007 du 12 juillet 2016 encadre l'utilisation des empreintes génétiques. \textbf{Attention :} L'ADN mitochondrial (hérité de la mère uniquement) est utilisé pour les échantillons très dégradés (os, cheveux sans bulbe) ou les enquêtes sur lignée maternelle. L'ADN Y (chromosome Y, père $\rightarrow$ fils) trace la lignée paternelle. Ces particularités s'expliquent par l'hérédité des chromosomes sexuels (XY vs XX) et de l'ADNmt (cytoplasme de l'ovocyte), notions que tu verras en génétique.
\end{savaistu}

\courssub{Exercice résolu : Analyse d'un caryotype et déduction méiotique}

\begin{exoresolu}[Analyse d'une anomalie numérique]
\textbf{Énoncé :} On observe le caryotype d'un fœtus obtenu après amniocentèse. Le chromosome 21 apparaît en \textbf{trois exemplaires} (trois chromosomes 21 bichromatidiens). Les autres paires sont normales (deux chromosomes bichromatidiens par paire).\\
\textbf{Questions :}
\begin{enumerate}
    \item Combien de chromosomes au total compte ce caryotype ? Justifie.
    \item Combien de chromatides pour le chromosome 21 ? Pour l'ensemble du génome ?
    \item Quelle est l'anomalie ? Comment s'est-elle produite au niveau de la méiose ? (Utilise le vocabulaire : chromosomes homologues, chromatides sœurs, non-disjonction).
    \item Si ce fœtus est une fille, quel est son caryotype complet (formule standardisée ISCN) ?
\end{enumerate}

\tcblower
\textbf{Solution détaillée :}
\begin{enumerate}
    \item \textbf{47 chromosomes.} Le comptage se fait par centromères. 22 paires d'autosomes normales = $22 \times 2 = 44$ chromosomes. Paire 21 anormale = 3 chromosomes. Sexosomes XX = 2 chromosomes. Total $= 44 + 3 + 2 = 47$.
    \item \textbf{Chromosome 21 :} 3 chromosomes bichromatidiens $\rightarrow$ $3 \times 2 = \textbf{6 chromatides}$. \textbf{Génome total :} 47 chromosomes bichromatidiens $\rightarrow$ $47 \times 2 = \textbf{94 chromatides}$.
    \item \textbf{Trisomie 21 libre (non-disjonction).} Lors de la méiose I ou II d'un des parents (souvent mère), les \textbf{deux chromosomes homologues 21} (méiose I) ou les \textbf{deux chromatides sœurs d'un chromosome 21} (méiose II) n'ont pas migré vers des pôles opposés. Le gamète résultant porte \textbf{deux chromosomes 21} (au lieu d'un). Après fécondation par un gamète normal (1 chr 21), le zygote en a 3.
    \item \textbf{47,XX,+21.} (47 chromosomes, sexe féminin, chromosome 21 en excès).
\end{enumerate}
\end{exoresolu}

\courssub{Bilan général de la leçon : Localisation et nature de l'information génétique}

Nous arrivons au terme de cette unité. Rassemblons les preuves expérimentales et les modèles construits pour formuler une conclusion rigoureuse, comme l'exige une épreuve de BEPC.

\begin{aretenir}[Conclusion générale justifiée : Où et quoi ?]
\begin{enumerate}
    \item \textbf{Localisation : Le noyau de la cellule.}
    \begin{itemize}
        \item \textbf{Preuve expérimentale :} Les expériences de \textbf{transfert de noyau} (Gurdon, grenouille \emph{Xenopus} ; Wilmut, brebis Dolly) montrent qu'un noyau de cellule différenciée (intestinale, mammaire) introduit dans un ovocyte énucléé dirige le développement complet d'un individu \textbf{génétiquement identique au donneur du noyau}. Le cytoplasme seul ne suffit pas. L'information héréditaire est donc dans le noyau.
    \end{itemize}
    \item \textbf{Support physique : Les chromosomes.}
    \begin{itemize}
        \item \textbf{Preuve cytologique :} Corrélation stricte entre comportement des chromosomes (séparation des homologues en méiose I, des sœurs en méiose II) et transmission des caractères (lois de Mendel). Nombre constant de chromosomes par espèce (2n = 46 chez l'homme). Caryotype = « carte d'identité » chromosomique.
        \item \textbf{Structure :} Chaque chromosome = une molécule d'ADN continue (très condensée) + protéines (histones). Paire d'homologues = 2 molécules d'ADN (paternelle + maternelle) portant les mêmes gènes.
    \end{itemize}
    \item \textbf{Nature chimique : L'ADN (Acide DésoxyriboNucléique).}
    \begin{itemize}
        \item \textbf{Preuve biochimique :} Expériences d'Avery (1944) : fraction « principe transformant » = ADN pur. Hershey-Chase (1952) : phage marqué $^{32}$P (ADN) entre dans la bactérie, $^{35}$S (protéines) reste dehors. TP N°1 : extraction et coloration positive (réaction de Feulgen ou DAPI) spécifique de l'ADN nuclear.
        \item \textbf{Information :} Séquence des 4 bases (A, T, G, C) le long de la molécule $\rightarrow$ écriture des gènes $\rightarrow$ synthèse des protéines (phénotype).
    \end{itemize}
\end{enumerate}
\textbf{Synthèse :} \cle{L'information génétique déterminant les caractères héréditaires est localisée dans le noyau des cellules, supportée par les chromosomes, dont la substance chimique est l'ADN.}
\end{aretenir}

\begin{exemplebox}[Attendu type BEPC : « Rédigez une conclusion justifiée »]
\textbf{Sujet :} « À partir des résultats des expériences de transfert de noyau et de l'observation des chromosomes, conclure sur la localisation et la nature de l'information génétique. »
\textbf{Rédaction modèle :}
\begin{quote}
« Les expériences de transfert de noyau (grenouille, brebis Dolly) démontrent que le noyau d'une cellule somatique contient toute l'information nécessaire au développement d'un individu entier : \textbf{l'information génétique est localisée dans le noyau}.
L'étude des chromosomes montre qu'ils se transmettent de manière stable lors des divisions (mitose, méiose) et que leur nombre est constant pour une espèce. Ils sont les \textbf{supports physiques de l'information génétique}.
L'analyse biochimique (Avery, Hershey-Chase, TP extraction) prouve que la molécule porteuse de l'information est l'\textbf{ADN}.
\textbf{En conclusion, l'information génétique héréditaire est localisée dans le noyau, supportée par les chromosomes, et sa nature chimique est l'ADN.} »
\end{quote}
\end{exemplebox}

\begin{exercice}[Entraînement BEPC - À faire seul(e)]
\begin{enumerate}
    \item Schématise et légende une \textbf{paire de chromosomes homologues bichromatidiens} en respectant le code couleur (père/mère). Indique le centromère, les chromatides sœurs, les chromatides non-sœurs.
    \item Complète le tableau :
    \begin{center}
    \begin{tabular}{|c|c|c|}\hline
    \textbf{État} & \textbf{Nombre de chromosomes (2n=46)} & \textbf{Nombre de chromatides} \\ \hline
    G1 (Monochromatidien) & & \\ \hline
    G2 / Métaphase (Bichromatidien) & & \\ \hline
    Après Anaphase Mitose & & \\ \hline
    Après Méiose I & & \\ \hline
    Après Méiose II (Gamète) & & \\ \hline
    \end{tabular}
    \end{center}
    \item Un caryotype montre 46 chromosomes dont deux chromosomes X et un chromosome Y. Quel est le sexe ? Quelle anomalie ? Donne la formule ISCN.
    \item Explique pourquoi deux frères (même père, même mère) ont des profils génétiques différents (sauf jumeaux monozygotes), en utilisant les notions de \textbf{méiose}, \textbf{chromosomes homologues} et \textbf{fécondation aléatoire}.
\end{enumerate}
\end{exercice}

\begin{corrige}[Exercice Entraînement BEPC]
\begin{enumerate}
    \item \textbf{Schéma :} Deux chromosomes en « X » (ou V/L), même taille, même centromère. Un rouge (2 chromatides rouges superposées), un bleu (2 chromatides bleues superposées). Légendes : Centromère (commun aux 2 chromatides sœurs), Chromatides sœurs (identiques, même couleur), Chromatides non-sœurs (différentes couleurs, allèles potentiellement différents), Paire d'homologues.
    \item \textbf{Tableau :}
    \begin{center}
    \begin{tabular}{|c|c|c|}\hline
    \textbf{État} & \textbf{Chr (2n=46)} & \textbf{Chromatides} \\ \hline
    G1 & 46 & 46 \\ \hline
    G2 / Métaphase & 46 & 92 \\ \hline
    Après Anaphase Mitose & 92 (transitoire) $\rightarrow$ 46/fille & 46/fille \\ \hline
    Après Méiose I & 23 (bichromatidiens) & 46 \\ \hline
    Après Méiose II & 23 (monochromatidiens) & 23 \\ \hline
    \end{tabular}
    \end{center}
    \item \textbf{Sexe :} Masculin (présence Y). \textbf{Anomalie :} Syndrome de Klinefelter (47,XXY). \textbf{Formule :} 47,XXY. (Non-disjonction sexosomes méiose I ou II).
    \item \textbf{Explication :} Lors de la méiose I, les paires d'homologues (père/mère) se séparent \textbf{aléatoirement} (loi de l'assortiment indépendant) : chaque gamète reçoit un mélange unique de chromosomes paternels et maternels ($2^{23}$ combinaisons possibles). S'y ajoute le \textbf{brassage génétique} (crossing-over entre chromatides non-sœurs en prophase I) qui crée des chromosomes recombinés. Enfin, la \textbf{fécondation} réunit au hasard un spermatozoïde parmi $\sim 10^8$ et un ovocyte parmi $\sim 400$ ovulations. La probabilité d'avoir deux frères génétiquement identiques (hors jumeaux) est quasi nulle.
\end{enumerate}
\end{corrige}

\newpage

\coursec{Activité d'intégration}

\courssub{Objectifs de l'activité}
\begin{itemize}
    \item Mobiliser les connaissances sur la localisation (noyau) et la nature (ADN, chromosomes) de l'information génétique pour résoudre un problème concret.
    \item Exploiter des résultats d'expériences de transfert de noyau, des caryotypes et des données quantitatives sur la composition en bases de l'ADN (règles de Chargaff).
    \item Rédiger un rapport d'expertise scientifique structuré, argumenté et conclu, en respectant la démarche d'investigation.
\end{itemize}

\courssub{Situation : L'affaire du bébé échangé à la maternité de l'Hôpital Central de Yaoundé}
\begin{popfigure}
\begin{tikzpicture}[scale=0.85, every node/.style={font=\small}]
% Styles
\tikzset{
    person/.style={draw, thick, rounded corners=4pt, minimum width=2.2cm, minimum height=1.1cm, align=center, fill=popcream},
    male/.style={person, fill=popblueL},
    female/.style={person, fill=poppinkL},
    baby/.style={person, fill=popgreenL},
    arrow/.style={-Stealth, thick, popdark},
    labelbox/.style={draw, dashed, rounded corners=3pt, fill=popblue!10, inner sep=4pt, font=\footnotesize, align=center}
}
% Parents Couple 1 (Bikoi)
\node[male, align=center] (P1) at (0, 4){M. BIKOI\\Père 1};
\node[female, align=center] (M1) at (4, 4){Mme BIKOI\\Mère 1};
% Parents Couple 2 (Ngoa)
\node[male, align=center] (P2) at (0, 0){M. NGOA\\Père 2};
\node[female, align=center] (M2) at (4, 0){Mme NGOA\\Mère 2};
% Bébés
\node[baby, align=center] (B1) at (8, 4){Bébé A\\Garçon};
\node[baby, align=center] (B2) at (8, 0){Bébé B\\Fille};
% Flèches union
\draw[arrow] (P1.east) -- (M1.west) node[midway, above, font=\tiny] {Couple 1};
\draw[arrow] (P2.east) -- (M2.west) node[midway, above, font=\tiny] {Couple 2};
% Flèches filiation supposée (pointillées pour montrer le doute)
\draw[arrow, dashed, poporange] (M1.south east) .. controls (6,3) and (6,3.5) .. (B1.north west) node[midway, above right, font=\tiny] {Supposé ?};
\draw[arrow, dashed, poporange] (M2.south east) .. controls (6,-1) and (6,-0.5) .. (B2.north west) node[midway, below right, font=\tiny] {Supposé ?};
% Boîtes documents
\node[labelbox, align=center] (Doc1) at (12, 4){\textbf{Doc 1 : Transfert de noyau}\\Expérience type Gurdon :\\Noyau cellule intestinale $\to$ Œuf énucléé $\to$ Têtard\\génétiquement identique au donneur.};
\node[labelbox, align=center] (Doc2) at (12, 2){\textbf{Doc 2 : Caryotypes simplifiés}\\Paire 1 : Grande métacentrique (M. Bikoi, Bébé A).\\Paire 2 : Petite acrocentrique (Mme Ngoa, Bébé B).\\Sexe : Bébé A (XY), Bébé B (XX).};
\node[labelbox, align=center] (Doc3) at (12, 0){\textbf{Doc 3 : \% Bases ADN (Chargaff)}\\\begin{tabular}{lcccc}
& A & T & G & C \\ \hline
M. Bikoi & 30,1 & 30,1 & 19,9 & 19,9 \\
Mme Bikoi & 29,8 & 29,8 & 20,2 & 20,2 \\
M. Ngoa & 28,2 & 28,2 & 21,8 & 21,8 \\
Mme Ngoa & 31,0 & 31,0 & 19,0 & 19,0 \\
Bébé A & 30,0 & 30,0 & 20,0 & 20,0 \\
Bébé B & 29,6 & 29,6 & 20,4 & 20,4
\end{tabular}};
\end{tikzpicture}
\legende{Situation-problème : deux couples (Bikoi et Ngoa) et deux nouveau-nés (A et B) à la maternité de Yaoundé. Les documents 1, 2 et 3 sont fournis pour l'enquête.}
\end{popfigure}

Le \textbf{15 mars 2026}, à la maternité de l'Hôpital Central de Yaoundé, deux accouchements ont eu lieu quasi simultanément. Le couple \textbf{BIKOI} (M. Bikoi et Mme Bikoi) a accueilli un garçon, tandis que le couple \textbf{NGOA} (M. Ngoa et Mme Ngoa) a accueilli une fille. Quelques heures plus tard, les deux pères remarquent que les bracelets d'identification des nouveau-nés ont été retirés par erreur lors d'un changement de couche. Le personnel soignant ne peut plus garantir quel bébé appartient à quel couple. Les deux familles, très angoissées, exigent une expertise génétique rapide.

En tant qu'expert biologiste mandaté par le tribunal de grande instance du Mfoundi, vous disposez de trois documents (ci-contre) pour identifier l'appartenance de chaque enfant :
\begin{enumerate}
    \item \textbf{Document 1 :} Rappel des résultats fondamentaux d'une expérience de transfert de noyau (type John Gurdon / Ian Wilmut).
    \item \textbf{Document 2 :} Caryotypes schématisés des quatre parents et des deux bébés, mettant en évidence deux particularités chromosomiques morphologiques distinctes (taille et position du centromère) sur une paire d'autosomes, ainsi que les chromosomes sexuels.
    \item \textbf{Document 3 :} Pourcentages molaires des quatre bases azotées (Adénine, Thymine, Guanine, Cytosine) dans l'ADN extrait de globules blancs de chaque individu.
\end{enumerate}

\courssub{Consignes}
\begin{enumerate}
    \item \textbf{Localisation de l'information :} En vous appuyant \emph{uniquement} sur le Document 1, expliquez pourquoi l'analyse du noyau (et non du cytoplasme) permet de trancher l'affaire. Précisez la conclusion tirée de ce type d'expérience sur le support de l'hérédité.
    \item \textbf{Analyse caryotypique (Document 2) :}
    \begin{enumerate}
        \item Identifiez la particularité chromosomique transmissible présente chez M. Bikoi et retrouvée chez l'un des bébés. Nommez ce type de chromosome (morphologie).
        \item Identifiez la particularité chromosomique transmissible présente chez Mme Ngoa et retrouvée chez l'autre bébé.
        \item En déduire l'appartenance biologique de chaque bébé (Bébé A et Bébé B) à l'aide des chromosomes sexuels et des particularités autosomiques.
    \end{enumerate}
    \item \textbf{Analyse moléculaire (Document 3) :}
    \begin{enumerate}
        \item Vérifiez que les résultats respectent les règles de Chargaff pour chaque individu. Quelle propriété chimique de l'ADN ces règles illustrent-elles ?
        \item Comparez les profils de Bébé A et Bébé B à ceux des couples. Quel couple présente le profil le plus proche de Bébé A ? De Bébé B ? (Justifiez par une comparaison chiffrée).
    \end{enumerate}
    \item \textbf{Synthèse et rapport d'expertise :} Rédigez le rapport d'expertise officiel (maximum 10 lignes) adressé au Procureur de la République. Il doit comporter : l'énoncé du problème, la démarche suivie (documents exploités), les résultats concordants des deux approches (caryotypique et moléculaire), et la conclusion formelle d'attribution de chaque enfant à ses parents biologiques.
\end{enumerate}

\courssub{Ressources à mobiliser}
\begin{aretenir}[Rappels de cours utiles]
    \begin{itemize}
        \item \textbf{Expérience de transfert de noyau :} Le noyau d'une cellule différenciée (ex. intestin) transplanté dans un ovocyte énucléé dirige le développement d'un individu entier possédant le patrimoine génétique du donneur du noyau. \cle{Conclusion :} L'information génétique est localisée dans le \cle{noyau}.
        \item \textbf{Chromosomes :} Supports physiques de l'information génétique. \cle{Chromosomes homologues :} même taille, même forme (position centromère), mêmes gènes (allèles éventuellement différents). Un homologue vient du père, l'autre de la mère.
        \item \textbf{ADN :} Molécule constituant les chromosomes. Double hélice avec appariement spécifique des bases : \cle{A = T} et \cle{G = C} (Règles de Chargaff). L'ADN est le support chimique de l'information génétique.
        \item \cle{Caryotype :} Photographie des chromosomes classés par paires homologues (taille décroissante, centromère). Permet de voir le sexe (XX/XY) et certaines anomalies ou polymorphismes morphologiques héréditaires.
    \end{itemize}
\end{aretenir}

\courssub{Démarche et corrigé commenté}
\begin{exoresolu}[Activité d'intégration : L'affaire des bébés échangés]
\textbf{Énoncé :} À la maternité de l'Hôpital Central de Yaoundé, les bracelets de deux nouveau-nés (Bébé A, garçon ; Bébé B, fille) ont été inversés. Deux couples (Bikoi et Ngoa) réclament leur enfant. En tant qu'expert, exploitez les trois documents fournis (transfert de noyau, caryotypes, composition en bases ADN) pour identifier les parents biologiques et rédiger un rapport d'expertise.
\tcblower
\textbf{1. Localisation de l'information génétique (Doc 1)}
L'expérience de transfert de noyau (type Gurdon) montre qu'un noyau prélevé sur une cellule somatique différenciée (cellule intestinale) et introduit dans un ovocyte privé de son propre noyau (énucléé) permet le développement d'un organisme complet (têtard puis grenouille). Cet organisme présente \emph{exclusivement} les caractères génétiques du donneur du noyau. Le cytoplasme de l'ovocyte receveur (mitochondries, ribosomes, ARN maternels) n'a pas imposé ses caractères.
\par \textbf{Justification :} Puisque le phénotype de l'individu obtenu est dicté par le noyau transplanté, l'information génétique déterminant les caractères héréditaires est localisée dans le \textbf{noyau}. C'est donc l'ADN nucléaire (organisé en chromosomes) qu'il faut analyser pour identifier les parents biologiques.

\textbf{2. Analyse caryotypique (Doc 2)}
\begin{itemize}
    \item \textbf{Particularité du couple Bikoi :} M. Bikoi présente une \textbf{paire d'autosomes n°1 grande et métacentrique} (centromère médian, bras égaux). Cette particularité morphologique est retrouvée \emph{identique} chez \textbf{Bébé A} (garçon, XY). Comme un fils reçoit son chromosome Y de son père et un homologue de chaque paire autosomique de son père, cette similitude forte indique une filiation père-fils.
    \item \textbf{Particularité du couple Ngoa :} Mme Ngoa présente une \textbf{paire d'autosomes n°2 petite et acrocentrique} (centromère sub-terminal, un bras très court). Cette particularité est retrouvée chez \textbf{Bébé B} (fille, XX). Une fille reçoit un chromosome X de sa mère et un homologue de chaque paire autosomique ; la présence de cette paire acrocentrique spécifique chez la mère et l'enfant confirme la filiation mère-fille.
    \item \textbf{Attribution caryotypique :}
    \begin{itemize}
        \item \textbf{Bébé A (XY, grande paire 1 métacentrique)} $\to$ Fils biologique du couple \textbf{BIKOI}.
        \item \textbf{Bébé B (XX, petite paire 2 acrocentrique)} $\to$ Fille biologique du couple \textbf{NGOA}.
    \end{itemize}
\end{itemize}

\textbf{3. Analyse moléculaire : composition en bases de l'ADN (Doc 3)}
\begin{itemize}
    \item \textbf{Vérification des règles de Chargaff :} Pour \emph{chaque} individu du tableau, on observe : $\%A \approx \%T$ et $\%G \approx \%C$ (ex. M. Bikoi : A=30,1 \%, T=30,1 \% ; G=19,9 \%, C=19,9 \%). La somme $\%A+\%T+\%G+\%C = 100\%$.
    \item \textbf{Signification :} Ces égalités traduisent l'appariement spécifique des bases dans la double hélice : \textbf{A s'apparie avec T} (2 liaisons H) et \textbf{G s'apparie avec C} (3 liaisons H). C'est la preuve structurelle que l'ADN est une double hélice à brins complémentaires.
    \item \textbf{Comparaison des profils (Identification moléculaire) :}
    \begin{itemize}
        \item \textbf{Profil « Bikoi » :} Teneur élevée en A+T ($\approx 60\%$), faible en G+C ($\approx 40\%$). M. Bikoi (A+T=60,2\%), Mme Bikoi (A+T=59,6\%).
        \item \textbf{Profil « Ngoa » :} Teneur plus faible en A+T ($\approx 56-62\%$ variable), plus forte en G+C. M. Ngoa (A+T=56,4\%), Mme Ngoa (A+T=62,0\%).
        \item \textbf{Bébé A :} A=30,0 \%, T=30,0 \% $\to$ A+T = \textbf{60,0 \%} ; G+C = 40,0 \%. Ce profil est \textbf{quasi identique} à la moyenne du couple Bikoi (Moyenne parents Bikoi : A+T $\approx$ 59,9 \%). Il est très différent de M. Ngoa (56,4 \%) et Mme Ngoa (62,0 \%).
        \item \textbf{Bébé B :} A=29,6 \%, T=29,6 \% $\to$ A+T = \textbf{59,2 \%} ; G+C = 40,8 \%. Ce profil est \textbf{proche de la moyenne du couple Ngoa} (Moyenne parents Ngoa : A+T $\approx$ 59,2 \%). Il s'éloigne du profil Bikoi (60,0 \%).
    \end{itemize}
    \textbf{Conclusion moléculaire :} Les pourcentages de bases confirment sans ambiguïté : \textbf{Bébé A appartient au couple Bikoi} (profil AT-riche) et \textbf{Bébé B appartient au couple Ngoa}.
\end{itemize}

\textbf{4. Rapport d'expertise officiel (Synthèse)}

\begin{center}
\begin{tabularx}{\linewidth}{|X|}\hline
\rowcolor{popblueL}
\textbf{RAPPORT D'EXPERTISE GÉNÉTIQUE — AFFAIRE BÉBÉS ÉCHANGÉS — MATERNITÉ HÔPITAL CENTRAL YAOUNDÉ} \\ \hline
\textbf{Objet :} Identification des liens de filiation biologique entre les nouveau-nés A (garçon) et B (fille) et les couples Bikoi / Ngoa. \\
\textbf{Démarche :} Exploitation de trois documents scientifiques : (1) Preuve par transfert de noyau que l'information héréditaire réside dans le noyau/ADN ; (2) Analyse comparative des caryotypes (chromosomes sexuels + polymorphismes morphologiques autosomiques) ; (3) Analyse quantitative de la composition en bases de l'ADN (règles de Chargaff). \\
\textbf{Résultats :} Le caryotype de Bébé A (XY) partage une paire 1 grande métacentrique spécifique avec M. Bikoi. Le caryotype de Bébé B (XX) partage une paire 2 petite acrocentrique spécifique avec Mme Ngoa. Les profils de bases de l'ADN confirment : Bébé A (A+T=60,0\%) match le profil AT-riche du couple Bikoi (moyenne 59,9\%) ; Bébé B (A+T=59,2\%) match le profil du couple Ngoa (moyenne 59,2\%). Les deux approches sont \textbf{concordantes}. \\
\textbf{Conclusion :} \textbf{Bébé A est le fils biologique de M. et Mme BIKOI. Bébé B est la fille biologique de M. et Mme NGOA.} Il n'y a pas eu d'échange ; les bracelets ont été inversés \emph{a posteriori}. \\
Fait à Yaoundé, le 16 mars 2026. \textit{L'Expert biologiste.} \\
\hline
\end{tabularx}
\end{center}
\end{exoresolu}

\courssub{Bilan de l'activité}
\begin{aretenir}[Ce que cette activité a permis de mettre en évidence]
    \begin{itemize}
        \item \textbf{Preuve expérimentale :} L'expérience de transfert de noyau est l'argument historique majeur localisant l'information génétique dans le \textbf{noyau} (donc sur les \textbf{chromosomes} / \textbf{ADN}), excluant le cytoplasme.
        \item \textbf{Approche cytogénétique (Caryotype) :} L'observation de particularités morphologiques héréditaires sur les chromosomes (taille, position du centromère : métacentrique, acrocentrique) et des chromosomes sexuels (XX/XY) permet une identification visuelle directe de la filiation. C'est une application concrète de la notion de \cle{chromosomes homologues} (un homologue paternel, un homologue maternel).
        \item \textbf{Approche moléculaire (ADN) :} La constance des pourcentages de bases au sein d'une espèce (règles de Chargaff : A=T, G=C) et la variabilité inter-individuelle/familiale du rapport (A+T)/(G+C) constituent une \textbf{signature moléculaire} fiable. L'ADN est bien la molécule support de l'information.
        \item \textbf{Concordance des échelles :} L'activité illustre la cohérence entre l'échelle \emph{macroscopique} (chromosomes visibles au microscope) et l'échelle \emph{moléculaire} (séquence de bases), toutes deux porteuses de la même information héréditaire.
        \item \textbf{Compétence rédactionnelle :} La rédaction d'un rapport d'expertise impose une structure rigoureuse (Problème $\to$ Démarche $\to$ Résultats $\to$ Conclusion) et un langage précis, sans ambiguïté, utilisable en contexte judiciaire ou médical.
    \end{itemize}
\end{aretenir}

\begin{attention}[Pièges à éviter pour le BEPC]
    \begin{itemize}
        \item \textbf{Confondre noyau et chromosome :} Le noyau \emph{contient} les chromosomes. L'expérience de transfert de noyau montre que l'info est dans le noyau ; l'analyse caryotypique montre qu'elle est \emph{sur} les chromosomes.
        \item \textbf{Oublier les règles de Chargaff :} Ne pas vérifier A=T et G=C pour valider la fiabilité des données du Doc 3. Ne pas citer ces règles pour justifier la structure en double hélice.
        \item \textbf{Rédiger une conclusion sans chiffres :} « Bébé A ressemble à M. Bikoi » est insuffisant. Il faut écrire : « Bébé A a 60,0\% A+T comme le couple Bikoi (59,9\%) contre 59,2\% pour le couple Ngoa ».
        \item \textbf{Mélanger les échelles :} Ne pas dire « le gène est sur le noyau ». Dire : « L'ADN (molécule) est organisé en chromosomes (structure) dans le noyau (organite) ».
    \end{itemize}
\end{attention}

\newpage

\coursec{Méthodes et exercices résolus}

\courssub{Méthodes et exercices résolus}

\begin{methode}[Analyser une expérience de transfert de noyau pour localiser l'information génétique]
\textbf{Objectif :} À partir du protocole et des résultats d'une expérience de transfert de noyau (Gurdon sur Xenope ou Wilmut sur brebis Dolly), conclure sur la localisation de l'information génétique.
\begin{enumerate}
    \item \textbf{Identifier le matériel de départ :} Donneur (noyau d'une cellule différenciée : intestin, glande mammaire) et Receveur (ovocyte énucléé, noyau retiré).
    \item \textbf{Repérer la variable manipulée :} Le noyau transféré (source de l'information héréditaire testée).
    \item \textbf{Analyser le résultat :} Le développement de l'embryon et les caractères de l'individu obtenu (clone) sont-ils ceux du \textbf{donneur du noyau} ou ceux du \textbf{receveur (ovocyte)} ?
    \item \textbf{Conclure :} Si l'individu obtenu présente les caractères du donneur du noyau, l'information génétique se trouve dans le \textbf{noyau} (plus précisément dans les chromosomes). Le cytoplasme de l'ovocyte assure seulement la machinerie cellulaire (mitochondries, ribosomes, facteurs de transcription).
    \item \textbf{Rédiger la conclusion type :} « L'expérience montre que le noyau d'une cellule différenciée contient toute l'information génétique nécessaire au développement complet d'un individu. L'information génétique est donc localisée dans le noyau, supportée par les chromosomes. »
\end{enumerate}
\end{methode}

\begin{exoresolu}[Expérience de Gurdon (1962) sur Xenope laevis]
\textbf{Énoncé :} John Gurdon prélève le noyau d'une cellule intestinale d'un têtard de Xenope (Donneur A) et l'injecte dans un ovocyte énucléé (Receveur B). L'ovocyte se développe en un têtard normal. On répète l'expérience avec un noyau de cellule cutanée d'une grenouille adulte (Donneur C). 
\begin{enumerate}
    \item Préciser la nature des cellules donneuses et receveuses.
    \item Quel est le génotype des têtards obtenus par rapport aux donneurs ?
    \item Conclure sur la localisation et la potentialité de l'information génétique dans les noyaux transférés.
\end{enumerate}
\tcblower
\textbf{Solution détaillée :}
\begin{enumerate}
    \item \textbf{Cellules donneuses :} Noyau de cellule intestinale de têtard (cellule différenciée, stade larvaire) pour la 1ère expérience ; noyau de cellule cutanée de grenouille adulte (cellule hautement différenciée) pour la 2nde. \textbf{Cellule receveuse :} Ovocyte énucléé (ovocyte dont le noyau maternel a été détruit ou retiré par UV/aspiration).
    \item Les têtards obtenus sont des \textbf{clones} des donneurs. Ils possèdent le \textbf{même génotype} (même ADN nucléaire) que la cellule donneuse (A ou C). Ils expriment les caractères héréditaires du donneur (ex: marqueurs génétiques, sexe).
    \item \textbf{Localisation :} Puisque le cytoplasme de l'ovocyte (Receveur B) est identique dans les deux expériences mais que les descendants sont différents (clones de A ou de C), l'information génétique déterminant les caractères se trouve dans le \textbf{noyau transféré}.
    \item \textbf{Potentialité (Totipotence) :} Le noyau d'une cellule intestinale de têtard (déjà différenciée) et même celui d'une cellule cutanée d'adulte (très différenciée) conservent \textbf{toute l'information génétique} nécessaire pour piloter le développement complet d'un nouvel individu. Ils sont \textbf{totipotents}. L'information génétique n'est pas perdue ni modifiée irréversiblement lors de la différenciation cellulaire ; seuls certains gènes sont exprimés ou réprimés.
\end{enumerate}
\textbf{Conclusion :} L'information génétique est localisée dans le \textbf{noyau} (sur les chromosomes) et conserve sa \textbf{totipotence} même dans les cellules différenciées.
\end{exoresolu}

\begin{methode}[Identifier et caractériser une paire de chromosomes homologues sur un caryotype]
\textbf{Objectif :} Sur une photographie de caryotype (ex: caryotype humain 2n=46), repérer une paire homologue, décrire sa structure (mono/bichromatidien) et justifier l'homologie.
\begin{enumerate}
    \item \textbf{Classer les chromosomes :} Les ranger par ordre de taille décroissante et position du centromère (métacentrique, sous-métacentrique, acrocentrique).
    \item \textbf{Repérer les paires :} Chercher \textbf{deux chromosomes} de \textbf{même taille}, \textbf{même position du centromère} (donc même rapport bras court/bras long) et \textbf{même profil de bandage} (bande G ou bande R).
    \item \textbf{Vérifier l'origine :} Rappeler qu'un chromosome de la paire vient du \textbf{père} (chromosome paternel) et l'autre de la \textbf{mère} (chromosome maternel).
    \item \textbf{Préciser l'état (Schéma) :}
        \begin{itemize}
            \item \textbf{Monochromatidien (1 chromatide) :} En $G_1$ (avant réplication) ou en fin de mitose/méiose II. Forme en \textbf{bâtonnet} ou \textbf{V} simple.
            \item \textbf{Bichromatidien (2 chromatides sœurs) :} En $G_2$ (après réplication S) ou en début de mitose/méiose I (prophase, métaphase). Forme en \textbf{X} (centromère médian) ou \textbf{Y} (centromère sub-terminal). Les deux chromatides sont \textbf{identiques} (réplication semi-conservatrice).
        \end{itemize}
    \item \textbf{Distinguer homologues et chromatides sœurs :} Homologues = même gènes (mêmes loci), allèles éventuellement différents, origine parentale différente. Chromatides sœurs = copies identiques (mêmes allèles), issues de la réplication d'un même chromosome.
\end{enumerate}
\end{methode}

\begin{exoresolu}[Analyse d'un caryotype humain et état des chromosomes]
\textbf{Énoncé :} On observe le caryotype d'une cellule humaine en métaphase de mitose (2n=46). La paire n°1 est métacentrique, la paire n°13 est acrocentrique.
\begin{enumerate}
    \item Combien de chromosomes et de chromatides compte cette cellule ? Justifier.
    \item Décrire les critères morphologiques permettant d'identifier la paire n°1 comme homologue.
    \item Schématiser un chromosome de la paire n°13 (acrocentrique) bichromatidien en légendant : centromère, bras court (p), bras long (q), chromatides sœurs.
    \item Quelle est la différence fondamentale entre les deux chromosomes de la paire n°1 et les deux chromatides d'un chromosome de la paire n°1 ?
\end{enumerate}
\tcblower
\textbf{Solution détaillée :}
\begin{enumerate}
    \item \textbf{Chromosomes :} 46 (2n=46, c'est le nombre de centromères). \textbf{Chromatides :} 92. \textbf{Justification :} En métaphase de mitose, la réplication de l'ADN (phase S) a déjà eu lieu. Chaque chromosome est \textbf{bichromatidien} (constitué de 2 chromatides sœurs). $46 \times 2 = 92$ chromatides.
    \item \textbf{Critères d'homologie (Paire n°1) :}
        \begin{itemize}
            \item \textbf{Taille identique} (plus grands chromosomes).
            \item \textbf{Position du centromère identique} : métacentrique (centromère au milieu, bras p et q de longueur égale).
            \item \textbf{Bandage identique} (même alternance bandes claires/sombres).
        \end{itemize}
        L'un est d'origine paternelle, l'autre maternelle.
    \item \textbf{Schéma (description textuelle pour le code TikZ) :}
        \begin{center}
        \begin{tikzpicture}[scale=1.2, thick]
            % Chromosome acrocentrique bichromatidien (X shape but centromere high)
            \draw[popblue, fill=popblueL] (0,0) -- (0.3, 1.5) -- (0.6, 1.5) -- (0.3, 0) -- cycle; % lower left
            \draw[popblue, fill=popblueL] (0.6,0) -- (0.9, 1.5) -- (0.6, 1.5) -- (0.3, 0) -- cycle; % lower right? No, acrocentric: short arm very small.
            % Better representation: Two long arms, tiny short arms.
            \draw[popdark, fill=popblue!20] (0.4, 0) -- (0.3, 0.3) -- (0.5, 0.3) -- cycle; % p arm left
            \draw[popdark, fill=popblue!20] (0.6, 0) -- (0.5, 0.3) -- (0.7, 0.3) -- cycle; % p arm right
            \draw[popdark, fill=popblue!20] (0.3, 0.3) -- (0.2, 2.5) -- (0.5, 2.5) -- (0.5, 0.3) -- cycle; % q arm left chromatid
            \draw[popdark, fill=popblue!20] (0.5, 0.3) -- (0.5, 2.5) -- (0.8, 2.5) -- (0.7, 0.3) -- cycle; % q arm right chromatid
            % Centromere region
            \draw[poporange, fill=poporange] (0.3, 0.3) rectangle (0.7, 0.35);
            \node[poporange, font=\tiny] at (0.5, 0.15) {Centromère};
            % Labels
            \node[left] at (0.2, 1.5) {Bras long (q)};
            \node[left] at (0.25, 0.2) {Bras court (p)};
            \node[right] at (0.8, 1.5) {Chromatides sœurs};
            \draw[<-, popdark] (0.8, 2.0) -- (1.2, 2.2) node[right, font=\small] {Chromatide 1};
            \draw[<-, popdark] (0.8, 1.0) -- (1.2, 0.8) node[right, font=\small] {Chromatide 2 (identique)};
        \end{tikzpicture}
        \end{center}
        \legende{Schéma d'un chromosome acrocentrique bichromatidien (paire 13-15, 21-22). Le centromère est sub-terminal. Les bras courts (p) portent les gènes des ARN ribosomiques (NOR).}
    \item \textbf{Différence fondamentale :}
        \begin{itemize}
            \item \textbf{Deux chromosomes homologues (Paire n°1) :} Portent les \textbf{mêmes gènes aux mêmes loci}, mais les \textbf{allèles peuvent être différents} (ex: allèle groupe sanguin A sur l'un, allèle B sur l'autre). Origine : un paternel, un maternel. \textbf{Hétérozygotie possible}.
            \item \textbf{Deux chromatides sœurs (d'un même chromosome) :} Sont des \textbf{copies identiques} issues de la réplication de l'ADN (semi-conservatrice). Elles portent \textbf{les mêmes allèles} à chaque locus (sauf erreur de réplication/mutation). \textbf{Identiques génétiquement}.
        \end{itemize}
\end{enumerate}
\end{exoresolu}

\begin{methode}[Construire une maquette de paire de chromosomes homologues (Monochromatidien vs Bichromatidien) — TP N°2]
\textbf{Objectif :} Réaliser une maquette physique (fil de fer, pâte à modeler, perles, papier) respectant l'échelle et la structure pour visualiser la réplication et l'homologie.
\begin{enumerate}
    \item \textbf{Matériel :} Fil de fer rigide (squelette), perles de couleurs (gènes/loci), pâte à modeler (chromatine), étiquettes (p/q, centromère, allèles).
    \item \textbf{Étape 1 : Chromosome monochromatidien (Modèle $G_1$).}
        \begin{itemize}
            \item Former un seul filament (double brin d'ADN + histones).
            \item Placer le \textbf{centromère} (constriction primaire) à la position définie (métacentrique/acrocentrique).
            \item Insérer des \textbf{perles/repères} le long des bras pour symboliser des \textbf{gènes (loci)}. Noter les allèles (ex: $A$ sur le paternel, $a$ sur le maternel).
        \end{itemize}
    \item \textbf{Étape 2 : Réplication (Phase S) $\rightarrow$ Chromosome bichromatidien (Modèle Mitose Métaphase).}
        \begin{itemize}
            \item Dupliquer \textbf{fidèlement} le filament initial : on obtient \textbf{deux chromatides sœurs} côte à côte.
            \item Les deux chromatides restent jointes au niveau du \textbf{centromère} (kinétochore).
            \item \textbf{Vérification cruciale :} Les \textbf{allèles} sur les loci correspondants des deux chromatides sœurs sont \textbf{identiques} ($A$ et $A$ sur les deux chromatides du chromosome paternel ; $a$ et $a$ sur les deux chromatides du chromosome maternel).
        \end{itemize}
    \item \textbf{Étape 3 : Assembler la paire homologue.}
        \begin{itemize}
            \item Placer côte à côte le chromosome paternel (bichromatidien) et le chromosome maternel (bichromatidien).
            \item Vérifier : Même taille, même centromère, mêmes loci (gènes aux mêmes positions), allèles potentiellement différents entre homologues ($A$ vs $a$), allèles identiques entre chromatides sœurs.
        \end{itemize}
    \item \textbf{Légender :} Indiquer $2n=2$ (pour une paire), $2c$ (contenu ADN double par rapport à $G_1$), centromère, chromatides sœurs, chromosomes homologues, allèles.
\end{enumerate}
\end{methode}

\begin{exoresolu}[Conception et analyse de maquettes de chromosomes homologues]
\textbf{Énoncé :} Un élève réalise une maquette d'une paire de chromosomes homologues métacentriques bichromatidiens pour un gène unique à deux allèles ($A$ et $a$). Il place l'allèle $A$ sur les deux chromatides du premier chromosome et l'allèle $a$ sur une seule chromatide du second chromosome (l'autre chromatide n'ayant pas d'allèle marqué).
\begin{enumerate}
    \item Identifier les erreurs dans cette maquette par rapport à la réalité biologique de la réplication de l'ADN.
    \item Corriger la maquette en décrivant la disposition correcte des allèles sur les 4 chromatides.
    \item Préciser le nombre de molécules d'ADN double brin présentes dans cette paire de chromosomes bichromatidiens.
\end{enumerate}
\tcblower
\textbf{Solution détaillée :}
\begin{enumerate}
    \item \textbf{Erreurs :}
        \begin{enumerate}
            \item \textbf{Non-respect de la réplication semi-conservatrice :} Les deux chromatides sœurs sont des copies identiques. Si l'une porte l'allèle $A$, l'autre \textbf{doit obligatoirement} porter l'allèle $A$. L'élève a mis $A$ sur une seule chromatide (ou a oublié de le copier).
            \item \textbf{Allèle manquant sur le chromosome maternel :} Le chromosome maternel (porteur de l'allèle $a$) est bichromatidien. Ses deux chromatides sœurs doivent \textbf{toutes deux} porter l'allèle $a$. L'absence d'allèle sur une chromatide implique une perte d'information génétique, ce qui est impossible en mitose normale.
        \end{enumerate}
    \item \textbf{Maquette corrigée (Disposition des allèles) :}
        \begin{itemize}
            \item \textbf{Chromosome Paternel (Homologue 1) :} 2 chromatides sœurs. \textbf{Chromatide 1 : Allèle $A$} — \textbf{Chromatide 2 : Allèle $A$}. (Génotype du chromosome : $[A//A]$ au niveau des chromatides).
            \item \textbf{Chromosome Maternel (Homologue 2) :} 2 chromatides sœurs. \textbf{Chromatide 1 : Allèle $a$} — \textbf{Chromatide 2 : Allèle $a$}. (Génotype du chromosome : $[a//a]$ au niveau des chromatides).
            \item \textbf{Paire homologue (Cellule en $G_2$/Métaphase) :} Génotype cellulaire $[A//a]$ (hétérozygote). 4 chromatides au total : $A, A, a, a$.
        \end{itemize}
    \item \textbf{Nombre de molécules d'ADN double brin :} 
        \begin{itemize}
            \item 1 chromatide = 1 molécule d'ADN double brin (continue, très condensée).
            \item 1 chromosome bichromatidien = 2 molécules d'ADN double brin.
            \item 1 paire de chromosomes homologues bichromatidiens = 2 chromosomes $\times$ 2 chromatides = \textbf{4 molécules d'ADN double brin}.
        \end{itemize}
        \textbf{Rappel :} En $G_1$ (monochromatidien), la même paire ne contient que 2 molécules d'ADN double brin. La phase S a doublé la quantité d'ADN ($2c \rightarrow 4c$).
\end{enumerate}
\end{exoresolu}

\begin{methode}[Mettre en évidence le matériel génétique (ADN) par la réaction de Feulgen — TP N°1]
\textbf{Objectif :} Réaliser et interpréter la coloration de Feulgen (spécifique de l'ADN) sur une lame de cellules (racine d'oignon, rate de rat, hépatiques) pour localiser l'ADN au microscope.
\begin{enumerate}
    \item \textbf{Principe biochimique :} La réaction de Feulgen repose sur l'hydrolyse de l'ADN par l'acide chlorhydrique (HCl) chaud $\rightarrow$ libération des bases désoxyribose $\rightarrow$ aldéhydes (groupements \ce{-CHO} sur le désoxyribose) $\rightarrow$ coloration par le réactif de Schiff (fuchsine décolorée par $\ce{SO2}$) $\rightarrow$ couleur \textbf{rose/violet magenta} (leuchofuchsine $\rightarrow$ fuchsine). \textbf{L'ARN n'est pas coloré} (ribose pas d'aldéhyde libre stable dans ces conditions).
    \item \textbf{Protocole clé (Étapes critiques) :}
        \begin{enumerate}
            \item Fixation (Éthanol/Acide acétique ou Carnoy) : figer les structures.
            \item \textbf{Hydrolyse acide (HCl 1N, 60°C, 8-10 min) :} \textbf{Étape décisive}. Trop court = pas de détection. Trop long = dégradation de l'ADN et perte du signal. \textit{Contrôle : Témoin sans hydrolyse (reste incolore)}.
            \item Rinçage eau distillée (arrêter l'hydrolyse).
            \item Coloration Réactif de Schiff (15-30 min, obscurité) : fixation des aldéhydes.
            \item Rinçage eau sulfureuse ($\ce{SO2}$) / eau distillée : éliminer le fond.
            \item Contre-coloration (Vert lumière / Fast Green) : visualiser le cytoplasme et les parois (vert) vs noyau (rose).
            \item Montage entre lame et lamelle (Canada balsam / Euparal).
        \end{enumerate}
    \item \textbf{Observation au microscope (Objectif $\times 40$ ou $\times 100$) :}
        \begin{itemize}
            \item \textbf{Noyau :} Coloré en \textbf{rose/violet foncé} (chromatine, chromosomes en mitose). Preuve de présence d'ADN.
            \item \textbf{Cytoplasme / Parois :} Coloré en \textbf{vert} (contre-coloration) ou incolore. \textbf{Pas de rose} $\rightarrow$ Pas d'ADN (sauf mitochondries/chloroplastes non visibles ou peu colorés par Feulgen standard).
            \item \textbf{Mitose :} Chromosomes bichromatidiens roses alignés à l'équateur (métaphase).
        \end{itemize}
    \item \textbf{Conclusion type :} « La coloration de Feulgen, spécifique de l'ADN, ne colore que le noyau (et les chromosomes en division). L'ADN, support de l'information génétique, est donc localisé dans le noyau. »
\end{enumerate}
\end{methode}

\begin{exoresolu}[Interprétation d'un TP Feulgen et contrôles]
\textbf{Énoncé :} Lors d'un TP de mise en évidence de l'ADN par la réaction de Feulgen sur des pointes de racines d'oignon (\textit{Allium cepa}), un élève obtient les résultats suivants sur deux lames :
\begin{itemize}
    \item \textbf{Lame 1 (Protocole complet) :} Noyaux des cellules en interphase colorés en rose vif ; chromosomes de cellules en métaphase colorés en rose vif ; cytoplasme et parois cellulaires colorés en vert clair.
    \item \textbf{Lame 2 (Oubli de l'hydrolyse HCl) :} Noyaux incolores (transparents) ; chromosomes invisibles ; cytoplasme et parois colorés en vert.
\end{itemize}
\begin{enumerate}
    \item Expliquer le rôle de l'hydrolyse acide (HCl chaud) dans la réaction de Feulgen.
    \item Pourquoi la Lame 2 ne montre-t-elle pas de coloration rose dans les noyaux ?
    \item À quoi sert la contre-coloration au Vert Lumière (Fast Green) ?
    \item En déduire la localisation de l'ADN dans la cellule interphasique et en division.
\end{enumerate}
\tcblower
\textbf{Solution détaillée :}
\begin{enumerate}
    \item \textbf{Rôle de l'hydrolyse acide (HCl 1N, 60°C) :} Elle provoque la \textbf{dépurination} de l'ADN : rupture des liaisons glycosidiques entre les bases azotées (purines surtout) et le désoxyribose. Cela libère des groupements \textbf{aldéhydes (-\ce{CHO})} sur les carbones C1' des désoxyriboses de la chaîne sucre-phosphate. Ces aldéhydes sont les sites de fixation du réactif de Schiff.
    \item \textbf{Lame 2 (Sans hydrolyse) :} Sans hydrolyse, les groupements aldéhydes ne sont pas libérés (les bases restent fixées sur les sucres). Le réactif de Schiff (fuchsine décolorée) ne peut pas réagir (pas de groupement \ce{-CHO} libre). Il n'y a donc \textbf{pas de restauration de la couleur rose (fuchsine)}. Le noyau reste incolore. Cela prouve que la coloration rose observée sur la Lame 1 est bien due à la réaction de Feulgen spécifique de l'ADN, et non à une coloration aspecificque.
    \item \textbf{Rôle de la contre-coloration (Vert Lumière) :} Elle colore les structures \textbf{non colorées par Feulgen} : cytoplasme, parois cellulaires, membranes, protéines. Elle permet de \textbf{visualiser la morphologie cellulaire} (limites de la cellule, forme du noyau) et de contraster le noyau rose sur fond vert. C'est un contrôle de la qualité de l'étalement et de la présence de cellules.
    \item \textbf{Localisation de l'ADN :}
        \begin{itemize}
            \item \textbf{Cellule en interphase :} ADN localisé dans le \textbf{noyau} (chromatine diffuse, nucléole non coloré car riche en ARN/protéines).
            \item \textbf{Cellule en division (mitose) :} ADN localisé dans les \textbf{chromosomes} (chromatine condensée), visibles comme des structures discrètes roses à l'équateur (métaphase) ou en migration (anaphase).
            \item \textbf{Absence :} Pas d'ADN détecté dans le cytoplasme (ni dans les mitochondries/chloroplastes par cette technique standard de Feulgen sur matériel végétal fixé standard).
        \end{itemize}
        \textbf{Conclusion :} L'ADN, molécule support de l'information génétique, est strictement localisé dans le \textbf{noyau} (interphase) et les \textbf{chromosomes} (division).
\end{enumerate}
\end{exoresolu}

\begin{methode}[Établir la relation hiérarchique : ADN $\rightarrow$ Chromosome $\rightarrow$ Gène $\rightarrow$ Caractère]
\textbf{Objectif :} Structurer ses connaissances pour expliquer comment l'information génétique est organisée physiquement et fonctionnellement, du macromolécule au phénotype.
\begin{enumerate}
    \item \textbf{Niveau 1 : L'ADN (Molécule).} Double hélice (Watson \& Crick). Suite de nucléotides (A, T, G, C). \textbf{Information = Séquence des bases}. Unité de l'hérédité chimique.
    \item \textbf{Niveau 2 : Le Gène (Unité fonctionnelle).} Segment d'ADN (locus précis sur un chromosome). Portion codant pour une protéine (ou ARN fonctionnel). \textbf{1 gène $\rightarrow$ 1 polypeptide (protéine)}. Allèles = versions différentes d'un même gène (mutation).
    \item \textbf{Niveau 3 : Le Chromosome (Support physique).} Molécule d'ADN très condensée (enroulée sur histones $\rightarrow$ nucléosomes $\rightarrow$ fibre de chromatine $\rightarrow$ boucles $\rightarrow$ chromatide). \textbf{Porteur linéaire de gènes} (des centaines à des milliers). Existe en paires homologues (diploïdie 2n).
    \item \textbf{Niveau 4 : Le Génome (Ensemble).} Ensemble du matériel génétique (ADN nucléaire + ADN mitochondrial/chloroplastique). Caryotype = photographie des chromosomes.
    \item \textbf{Niveau 5 : Le Caractère (Phénotype).} Trait observable (couleur yeux, groupe sanguin, résistance maladie). Résultat de l'expression des gènes (protéines) $\times$ Environnement.
    \item \textbf{Flux de l'information :} \textbf{ADN (Gène)} $\xrightarrow{\text{Transcription}}$ \textbf{ARNm} $\xrightarrow{\text{Traduction}}$ \textbf{Protéine} $\xrightarrow{\text{Fonction}}$ \textbf{Caractère}.
\end{enumerate}
\end{methode}

\begin{exoresolu}[De la molécule au caractère : Le cas de l'hémoglobine et de la drépanocytose]
\textbf{Énoncé :} La drépanocytose est une maladie génétique héréditaire due à une mutation du gène $\beta$-globine ($HBB$) sur le chromosome 11. L'allèle normal est noté $A$ (codant pour $\beta^A$), l'allèle muté $S$ (codant pour $\beta^S$). La mutation est une substitution de base (A $\rightarrow$ T) au 6ème codon (GAG $\rightarrow$ GTG) $\rightarrow$ Remplacement Glutamate (acide) par Valine (hydrophobe) en position 6 de la chaîne $\beta$.
\begin{enumerate}
    \item Classer les éléments suivants du plus petit au plus grand : Chromosome 11, Gène $HBB$, Nucléotide, Génome, Chaîne $\beta$-globine, Allèle $S$.
    \item Expliquer comment une modification au niveau du \textbf{nucléotide} (Niveau 1) entraîne une modification du \textbf{caractère} (Niveau 5 : forme des globules rouges, anémie).
    \item Préciser la localisation exacte de l'allèle $S$ dans le noyau d'une cellule d'un individu hétérozygote $AS$.
\end{enumerate}
\tcblower
\textbf{Solution détaillée :}
\begin{enumerate}
    \item \textbf{Ordre croissant :}
        \textbf{Nucléotide} (base + sucre + phosphate) $<$ \textbf{Allèle $S$} (version mutée du gène, suite de nucléotides) $=$ \textbf{Gène $HBB$} (locus, suite de nucléotides) $<$ \textbf{Chaîne $\beta$-globine} (protéine, produit du gène) $<$ \textbf{Chromosome 11} (support physique portant le gène $HBB$ et des milliers d'autres) $<$ \textbf{Génome} (ensemble des 23 paires de chromosomes + ADN mitochondrial).
        \textit{Note :} L'allèle est une forme du gène, même échelle physique (locus), mais le gène est l'entité générique.
    \item \textbf{Chaîne causale (Mécanisme) :}
        \begin{enumerate}
            \item \textbf{ADN (Nucléotide) :} Mutation ponctuelle (Substitution A$\rightarrow$T) dans l'exon 1 du gène $HBB$ sur le chromosome 11. $\rightarrow$ Changement du codon GAG (Glu) en GTG (Val).
            \item \textbf{ARNm (Transcription) :} L'ARNm mature porte le codon GUG (Val) au lieu de GAG (Glu).
            \item \textbf{Protéine (Traduction) :} Synthèse d'une chaîne $\beta^S$ anormale (Valine en position 6 au lieu de Glutamate). Perte d'une charge négative à la surface de la protéine.
            \item \textbf{Biochimie/Cellulaire :} En condition de faible $\ce{O2}$ (hypoxie), la Valine hydrophobe s'insère dans une poche hydrophobe d'une autre molécule d'hémoglobine $\rightarrow$ \textbf{Polymérisation de l'HbS} en fibres rigides $\rightarrow$ \textbf{Déformation du globule rouge} en forme de faucille (drépanocyte) $\rightarrow$ Rigidité, hémolyse, vaso-occlusion.
            \item \textbf{Caractère (Phénotype) :} Anémie hémolytique chronique, crises vaso-occlusives douloureuses, sensibilité aux infections. L'hétérozygote $AS$ (porteur sain) a un avantage sélectif contre le paludisme (sélection équilibrante, pertinent au Cameroun).
        \end{enumerate}
    \item \textbf{Localisation de l'allèle $S$ chez $AS$ :}
        \begin{itemize}
            \item Dans le \textbf{noyau} de chaque cellule somatique.
            \item Sur l'un des deux \textbf{chromosomes 11 homologues} (le chromosome maternel OU paternel, selon l'héritage).
            \item Au \textbf{locus précis du gène $HBB$} (bande 11p15.5 sur le bras court p du chromosome 11).
            \item Sur les \textbf{deux chromatides sœurs} de ce chromosome 11 (si la cellule est en $G_2$ ou M), car la réplication a copié la mutation.
            \item L'autre chromosome 11 homologue porte l'allèle $A$ (normal) au même locus.
        \end{itemize}
\end{enumerate}
\end{exoresolu}

\begin{methode}[Rédiger une conclusion scientifique argumentée (Savoir-faire transversal)]
\textbf{Objectif :} Structurer une réponse écrite (type Bac) qui démontre le raisonnement : Observation/Donnée $\rightarrow$ Interprétation $\rightarrow$ Conclusion.
\begin{enumerate}
    \item \textbf{Reformuler la question/problème} (1 phrase).
    \item \textbf{Citer les faits/observations/résultats expérimentaux} utilisés comme preuves (données chiffrées, observations au microscope, résultats de TP, documents fournis). Utiliser \textbf{« Car »} ou \textbf{« En effet »}.
    \item \textbf{Expliquer le mécanisme / L'interprétation scientifique} (mobiliser les savoirs du cours : structure ADN, réplication, mitose, Feulgen, transfert noyau). Utiliser le vocabulaire précis (chromosome, chromatide, allèle, locus, hydrolyse, totipotence).
    \item \textbf{Conclure explicitement} en répondant à la question initiale. Commencer par \textbf{« En conclusion, »} ou \textbf{« Ainsi, »}.
    \item \textbf{Relire :} Vérifier l'orthographe des termes scientifiques (chromosome, chromatide, centromère, allèle, locus, hydrolyse, Feulgen, totipotence), les accords, et l'absence de contradiction.
\end{enumerate}
\end{methode}

\begin{exoresolu}[Rédaction argumentée : Preuve que l'ADN est le support de l'information génétique]
\textbf{Énoncé :} À partir de l'expérience de transformation de Griffith (1928) et de l'identification du principe transformant par Avery, MacLeod et McCarty (1944), rédigez une conclusion argumentée montrant que l'ADN est le support de l'information génétique. (On rappelle : Griffith : Souche R vivante + Souche S tuée $\rightarrow$ Mort souris + Souche S vivante récupérée. Avery : Extrait de S purifié $\rightarrow$ Seule la fraction ADN transforme R en S. Destruction ADN $\rightarrow$ Perte transformation. Destruction Protéines/ARN $\rightarrow$ Transformation conservée).
\tcblower
\textbf{Solution détaillée :}
\textbf{Problème :} Quelle est la nature chimique du « principe transformant » responsable du transfert du caractère de virulence (capsule lisse) de la souche S à la souche R chez \textit{Streptococcus pneumoniae} ?

\textbf{Argumentation :}
\begin{enumerate}
    \item \textbf{Fait expérimental 1 (Griffith) :} Un facteur thermostable (résistant à la chaleur qui tue les bactéries S) présent dans l'extrait de souche S morte permet de transformer durablement la souche R non virulente en souche S virulente (capsulée). Ce facteur est l'information génétique de la virulence.
    \item \textbf{Fait expérimental 2 (Avery et al. - Fractionnement) :} L'extrait de S est fractionné (protéines, ARN, ADN, polysaccharides). \textbf{Seule la fraction ADN} conserve le pouvoir transformant. Les fractions protéines, ARN, polysaccharides sont inactives.
    \item \textbf{Fait expérimental 3 (Avery et al. - Destruction enzymatique) :} Le traitement de l'extrait actif par la \textbf{DNase} (enzyme détruisant l'ADN) \textbf{abolit totalement} le pouvoir transformant. En revanche, le traitement par la Protéinase (protéines) ou la RNase (ARN) \textbf{ne l'affecte pas}.
\end{enumerate}
\textbf{Interprétation :} Puisque la destruction spécifique de l'ADN supprime l'activité transformante (transfert de l'information de virulence), et que la destruction des autres macromolécules ne l'empêche pas, l'ADN est la \textbf{seule molécule nécessaire et suffisante} pour porter cette information héréditaire.

\textbf{Conclusion :} \textbf{L'ADN (Acide DésoxyriboNucléique) est le support chimique de l'information génétique.} C'est la molécule du gène. (Confirmé plus tard par Hershey \& Chase 1952 sur bactériophage).
\end{exoresolu}

\begin{attention}[Les 5 erreurs qui coûtent des points au BEPC]
\begin{enumerate}
    \item \textbf{Confondre Chromosome et Chromatide (et Homologue).}
    \textit{Erreur :} « La paire de chromosomes est constituée de 2 chromatides. » ou « Les chromatides sœurs sont des chromosomes homologues. »
    \textbf{Correct :} 1 Chromosome = 1 centromère. En $G_1$ : 1 chromatide (monochromatidien). En $G_2/M$ : 2 chromatides sœurs (bichromatidien). 2 Chromosomes homologues = 1 paternel + 1 maternel (même taille, même centromère, mêmes gènes, allèles éventuellement différents). \textbf{Compter les centromères pour compter les chromosomes.}
    
    \item \textbf{Oublier le rôle de l'hydrolyse HCl dans Feulgen (ou la confondre avec la coloration).}
    \textit{Erreur :} « Le réactif de Schiff colore l'ADN en rose. » (Sans mentionner l'hydrolyse).
    \textbf{Correct :} L'hydrolyse acide (HCl chaud) \textbf{libère les aldéhydes} sur le désoxyribose. C'est \textbf{ensuite} que le réactif de Schiff (fuchsine décolorée) se fixe sur ces aldéhydes et redevient rose. Sans hydrolyse $\rightarrow$ pas d'aldéhydes $\rightarrow$ pas de couleur. C'est la spécificité de la réaction.
    
    \item \textbf{Dire que le noyau de la cellule receveuse (ovocyte) participe au patrimoine génétique du clone.}
    \textit{Erreur :} « L'ovocyte apporte la moitié des chromosomes. » ou « Le clone a le patrimoine du noyau donneur ET de l'ovocyte. »
    \textbf{Correct :} L'ovocyte est \textbf{énucléé} (son noyau maternel est retiré/détruit). Le clone ne reçoit que le noyau du donneur. Le cytoplasme de l'ovocyte apporte les organites (mitochondries $\rightarrow$ ADN mitochondrial maternel) et le stock protéique/ARNm maternel pour les premières divisions, \textbf{pas les chromosomes nucléaires}.
    
    \item \textbf{Confondre Allèles différents (Hétérozygotie) et Chromatides sœurs différentes.}
    \textit{Erreur :} Sur un schéma de chromosome bichromatidien en métaphase, écrire allèle $A$ sur une chromatide et allèle $a$ sur l'autre chromatide du \textbf{même} chromosome.
    \textbf{Correct :} Les chromatides sœurs sont des \textbf{copies identiques} (réplication semi-conservatrice). Elles portent les \textbf{mêmes allèles}. Les allèles différents ($A$ et $a$) se trouvent sur les \textbf{deux chromosomes homologues} (l'un paternel, l'autre maternel), chacun étant bichromatidien.
    
    \item \textbf{Vocabulaire imprécis : « Gène » au lieu de « Allèle », « Caractère » au lieu de « Gène », « Noyau » au lieu de « Chromosome/ADN ».}
    \textit{Erreur :} « Le gène de la drépanocytose est sur le chromosome 11. » (Le gène $HBB$ est sur le chr11, l'allèle $S$ est la version mutée). « L'information génétique est dans le noyau. » (Trop vague : dans l'ADN des chromosomes \textbf{dans} le noyau).
    \textbf{Correct :} Utiliser les termes exacts du programme : \textbf{Locus, Allèle, Gène, Chromosome, Chromatide, Centromère, ADN, Nucléotide, Réplication, Mitose, Méiose, Feulgen, Hydrolyse, Transfert de noyau, Totipotence, Clonage.}
\end{enumerate}
\end{attention}

\newpage

\coursec{Exercices}

\courssub{Je vérifie mes connaissances}

\begin{exercice}[QCM : localisation de l'information génétique]
Pour chaque question, une seule réponse est exacte. Entoure la lettre correspondante.

\textbf{1.} L'information génétique d'une cellule est contenue principalement dans :
\begin{itemize}
\item[a)] le cytoplasme ;
\item[b)] le noyau ;
\item[c)] la membrane cellulaire.
\end{itemize}

\textbf{2.} Le support de l'information génétique est :
\begin{itemize}
\item[a)] le chromosome ;
\item[b)] la vacuole ;
\item[c)] le chloroplaste.
\end{itemize}

\textbf{3.} La molécule qui porte l'information génétique est :
\begin{itemize}
\item[a)] une protéine ;
\item[b)] un lipide ;
\item[c)] l'ADN.
\end{itemize}

\textbf{4.} Chez l'Homme, une cellule possède normalement :
\begin{itemize}
\item[a)] 23 chromosomes ;
\item[b)] 46 chromosomes ;
\item[c)] 48 chromosomes.
\end{itemize}
\end{exercice}

\begin{exercice}[Vrai ou faux ? Justifie ta réponse]
Réponds par \textbf{vrai} ou \textbf{faux}, puis justifie en une ou deux phrases.

\textbf{1.} Dans l'expérience de Hammerling sur l'algue \emph{Acetabularia}, c'est le pied (qui contient le noyau) qui détermine la forme du chapeau régénéré.

\textbf{2.} Deux chromosomes homologues portent toujours exactement les mêmes allèles.

\textbf{3.} Dans une molécule d'ADN, la quantité d'adénine (A) est toujours égale à la quantité de thymine (T).

\textbf{4.} Un chromosome bichromatidien est formé de deux chromatides sœurs réunies par un centromère.
\end{exercice}

\begin{exercice}[Définitions à compléter]
Recopie et complète chaque définition avec les mots qui conviennent.

\textbf{1.} Un \textbf{chromosome} est une structure du \dots\ldots\ldots qui porte l'information \dots\ldots\ldots ; il est constitué d'une molécule d'\dots\ldots\ldots associée à des protéines.

\textbf{2.} L'\textbf{ADN} est une molécule formée de deux brins enroulés en \dots\ldots\ldots, constitués de nucléotides dont les bases sont : \dots\ldots\ldots, \dots\ldots\ldots, \dots\ldots\ldots et \dots\ldots\ldots.

\textbf{3.} Deux chromosomes sont dits \textbf{homologues} lorsqu'ils ont la même \dots\ldots\ldots et portent les mêmes \dots\ldots\ldots aux mêmes emplacements, l'un étant d'origine \dots\ldots\ldots et l'autre d'origine \dots\ldots\ldots.
\end{exercice}

\begin{exercice}[Calcul direct : complémentarité des bases]
L'analyse d'un fragment d'ADN extrait de cellules de la moelle osseuse d'un patient du Centre Hospitalier d'Essos (Yaoundé) donne les pourcentages de bases suivants : A = \SI{28}{\%} et G = \SI{22}{\%}.

\textbf{1.} Rappelle la règle de complémentarité des bases de l'ADN.

\textbf{2.} Calcule le pourcentage de thymine (T) et le pourcentage de cytosine (C) de ce fragment, en justifiant chaque calcul.
\end{exercice}

\courssub{Je m'entraîne}

\begin{exercice}[Le brin complémentaire]
On donne la séquence d'un brin d'ADN : A -- T -- G -- C -- C -- A -- G -- T.

\textbf{1.} Écris la séquence du brin complémentaire, en alignant les bases deux à deux.

\textbf{2.} Combien de liaisons A--T et de liaisons G--C comporte ce fragment d'ADN double brin ?

\textbf{3.} Explique en deux phrases en quoi la succession des bases le long d'un brin d'ADN constitue une \cle{information génétique}.
\end{exercice}

\begin{exercice}[Schéma d'une paire de chromosomes homologues]
\textbf{1.} Schématise et légende une paire de chromosomes homologues \textbf{bichromatidiens} (en forme de X), en utilisant deux couleurs différentes pour distinguer le chromosome d'origine paternelle de celui d'origine maternelle.

\textbf{2.} Sur ton schéma, place : le centromère, les chromatides sœurs, et deux gènes (couleur des yeux, groupe sanguin) aux mêmes emplacements sur les deux homologues.

\textbf{3.} Donne une différence entre un chromosome monochromatidien et un chromosome bichromatidien.
\end{exercice}

\begin{exercice}[Protocole du TP d'extraction de l'ADN]
Au cours du TP de mise en évidence du matériel génétique, on réalise l'extraction de l'ADN de cellules d'oignon (ou de banane écrasée) selon le protocole suivant : broyage du matériel dans une solution contenant du détergent (liquide vaisselle) et du sel ; filtration ; ajout d'alcool glacé sur le filtrat.

\textbf{1.} Quel est le rôle du détergent dans ce protocole ?

\textbf{2.} Pourquoi utilise-t-on de l'alcool \emph{glacé} ? Qu'observe-t-on à l'interface entre le filtrat et l'alcool ?

\textbf{3.} Pourquoi faut-il broyer le matériel végétal avant l'extraction ?
\end{exercice}

\begin{exercice}[Maquette de chromosomes]
Pour le TP de conception de maquettes, un groupe d'élèves de 3ème du collège de Bafoussam fabrique des chromosomes avec des boudins de pâte à modeler de deux couleurs (rouge et bleu).

\textbf{1.} Propose une utilisation de ces boudins pour représenter une paire de chromosomes homologues monochromatidiens, puis bichromatidiens.

\textbf{2.} Combien de boudins faut-il pour représenter les chromosomes d'une cellule humaine à 46 chromosomes, tous monochromatidiens ?

\textbf{3.} Que représente la différence de couleur entre les boudins ?
\end{exercice}

\courssub{Je me prépare au BEPC}

\begin{exercice}[Type BEPC : l'expérience de Hammerling]
L'algue verte \emph{Acetabularia} est constituée de trois parties : un \textbf{pied} (qui contient le noyau), un \textbf{stipe} (tige) et un \textbf{chapeau}. Il existe deux espèces : \emph{A. acetabulum} (chapeau arrondi) et \emph{A. crenulata} (chapeau découpé). Le biologiste Hammerling a réalisé des greffes dont les résultats sont résumés dans le tableau ci-dessous.

\begin{center}
\begin{tabularx}{\linewidth}{|c|X|X|X|}
\hline
\rowcolor{popblueL}
\textbf{Expérience} & \textbf{Pied greffé} & \textbf{Stipe + chapeau} & \textbf{Chapeau régénéré} \\
\hline
1 & \emph{A. acetabulum} & \emph{A. crenulata} & arrondi (type \emph{acetabulum}) \\
\hline
2 & \emph{A. crenulata} & \emph{A. acetabulum} & découpé (type \emph{crenulata}) \\
\hline
3 & pied enlevé (sans noyau) & \emph{A. acetabulum} & pas de régénération \\
\hline
\end{tabularx}
\end{center}

\textbf{1.} Dégager le problème posé par cette série d'expériences. (1 pt)

\textbf{2.} Formuler une hypothèse permettant d'expliquer les résultats des expériences 1 et 2. (1 pt)

\textbf{3.} Comparer les résultats des expériences 1 et 2 : quel élément du greffon détermine la forme du chapeau régénéré ? (1 pt)

\textbf{4.} Quel est le rôle de l'expérience 3 ? Que montre-t-elle ? (1 pt)

\textbf{5.} Conclure sur la localisation de l'information génétique dans la cellule. (1 pt)
\end{exercice}

\begin{exercice}[Type BEPC : analyse d'un caryotype humain]
Au laboratoire d'analyses d'un hôpital de Douala, on réalise le caryotype (photographie ordonnée des chromosomes d'une cellule) d'un nouveau-né présentant des signes particuliers. Le document obtenu montre des chromosomes bichromatidiens classés par paires d'homologues, numérotées de 1 à 22, plus les chromosomes sexuels. Le technicien dénombre : 22 paires de chromosomes ordinaires, deux chromosomes sexuels identiques (XX), et \textbf{trois} exemplaires du chromosome 21.

\textbf{1.} Définir le terme \cle{caryotype}. (1 pt)

\textbf{2.} Combien de chromosomes au total compte une cellule de ce nouveau-né ? Comparer avec le nombre normal. (1 pt)

\textbf{3.} Déterminer le sexe chromosomique de cet enfant. Justifier. (1 pt)

\textbf{4.} Nommer l'anomalie observée et la maladie qui en résulte. (1 pt)

\textbf{5.} Rappeler ce que sont deux chromosomes homologues et préciser leur origine chez un individu. (1 pt)
\end{exercice}

\begin{exercice}[Type BEPC : nature de l'information génétique]
Dans le cadre de la semaine scientifique de son collège à Ngaoundéré, un élève de 3ème étudie la molécule d'ADN. On lui fournit les données suivantes, obtenues par analyse chimique de l'ADN de deux organismes différents.

\begin{center}
\begin{tabularx}{\linewidth}{|l|X|X|X|X|}
\hline
\rowcolor{popblueL}
\textbf{Organisme} & \textbf{A (\%)} & \textbf{T (\%)} & \textbf{G (\%)} & \textbf{C (\%)} \\
\hline
Homme & 30,4 & 30,1 & 19,6 & 19,9 \\
\hline
Maïs & 26,8 & 27,2 & 22,8 & 23,2 \\
\hline
\end{tabularx}
\end{center}

\textbf{1.} Décrire brièvement la structure de la molécule d'ADN (nombre de brins, disposition, éléments constitutifs). (1,5 pt)

\textbf{2.} À partir du tableau, dégager deux relations numériques entre les pourcentages des bases, valables pour les deux organismes. Quelle règle ces relations illustrent-elles ? (1,5 pt)

\textbf{3.} Un brin d'ADN humain porte la séquence : T -- A -- C -- G -- G -- T. Écrire la séquence du brin complémentaire. (1 pt)

\textbf{4.} Expliquer pourquoi les pourcentages de bases sont différents entre l'Homme et le maïs, alors que la règle de la question 2 est respectée dans les deux cas. (1 pt)

\textbf{5.} En t'appuyant sur tes connaissances, expliquer en quelques phrases le lien entre l'ADN, les chromosomes et les caractères héréditaires. (1 pt)
\end{exercice}



\newpage

\coursec{Corrigés des exercices}

\coursec{Corrigés des exercices — Leçon 02 : Localisation et nature de l'information génétique}

\begin{corrige}[Exercice 1 — QCM : localisation de l'information génétique]
\textbf{1. b)} Les expériences de transfert de noyau (Hammerling sur \emph{Acetabularia}) montrent que l'information génétique se trouve dans le \cle{noyau}.

\textbf{2. a)} Les \cle{chromosomes}, visibles dans le noyau au moment de la division cellulaire, sont les supports de l'information génétique.

\textbf{3. c)} L'\cle{ADN} (acide désoxyribonucléique) est la molécule qui constitue les chromosomes et qui porte l'information génétique.

\textbf{4. b)} Une cellule humaine possède 46 chromosomes, soit 23 paires de chromosomes homologues (22 paires ordinaires + 1 paire de chromosomes sexuels).

\textbf{Point de vigilance :} ne pas confondre le nombre de chromosomes (46) avec le nombre de paires (23). Au BEPC, cette confusion est la plus fréquente.
\end{corrige}

\begin{corrige}[Exercice 2 — Vrai ou faux ? Justifie ta réponse]
\textbf{1. Vrai.} Le chapeau régénéré a toujours la forme caractéristique de l'espèce qui a fourni le pied ; or le pied est la seule partie qui contient le noyau. C'est donc le noyau qui détient l'information déterminant la forme du chapeau.

\textbf{2. Faux.} Deux chromosomes homologues portent les mêmes \emph{gènes} aux mêmes emplacements, mais les \emph{allèles} (versions d'un gène) peuvent être différents : par exemple, un allèle « yeux bruns » sur l'un et un allèle « yeux bleus » sur l'autre.

\textbf{3. Vrai.} Dans l'ADN double brin, l'adénine est toujours appariée à la thymine (A avec T) et la guanine à la cytosine (G avec C) : c'est la règle de \cle{complémentarité des bases}. On a donc toujours A = T et G = C.

\textbf{4. Vrai.} Après la duplication de l'ADN, chaque chromosome est formé de deux chromatides sœurs identiques, réunies au niveau du centromère : on dit qu'il est \cle{bichromatidien} (aspect en X).

\textbf{Point de vigilance :} une réponse « vrai » ou « faux » sans justification ne rapporte aucun point : la justification est toujours exigée.
\end{corrige}

\begin{corrige}[Exercice 3 — Définitions à compléter]
\textbf{1.} Un chromosome est une structure du \textbf{noyau} qui porte l'information \textbf{génétique} ; il est constitué d'une molécule d'\textbf{ADN} associée à des protéines.

\textbf{2.} L'ADN est une molécule formée de deux brins enroulés en \textbf{double hélice}, constitués de nucléotides dont les bases sont : \textbf{adénine (A)}, \textbf{thymine (T)}, \textbf{guanine (G)} et \textbf{cytosine (C)}.

\textbf{3.} Deux chromosomes sont dits homologues lorsqu'ils ont la même \textbf{forme (taille et aspect)} et portent les mêmes \textbf{gènes} aux mêmes emplacements, l'un étant d'origine \textbf{paternelle} et l'autre d'origine \textbf{maternelle}.

\textbf{Point de vigilance :} à la question 3, « mêmes gènes » et non « mêmes allèles » : c'est précisément ce qui distingue l'homologie de l'identité totale.
\end{corrige}

\begin{corrige}[Exercice 4 — Calcul direct : complémentarité des bases]
\textbf{1.} Dans une molécule d'ADN double brin, les bases sont appariées deux à deux : l'adénine avec la thymine (A--T) et la guanine avec la cytosine (G--C). Il en résulte que les quantités sont égales deux à deux : A = T et G = C.

\textbf{2.} D'après la règle de complémentarité :
\[ T = A = \SI{28}{\%} \qquad \text{et} \qquad C = G = \SI{22}{\%}. \]
Vérification : A + T + G + C $= 28 + 28 + 22 + 22 = \SI{100}{\%}$. Le résultat est cohérent.

\textbf{Point de vigilance :} toujours vérifier que la somme des quatre pourcentages vaut \SI{100}{\%} ; c'est un contrôle rapide qui évite les erreurs de recopie.
\end{corrige}

\begin{corrige}[Exercice 5 — Le brin complémentaire]
\textbf{1.} En appliquant la règle A avec T et G avec C :
\begin{center}
\begin{tabular}{cccccccc}
A & T & G & C & C & A & G & T \\
$\vert$ & $\vert$ & $\vert$ & $\vert$ & $\vert$ & $\vert$ & $\vert$ & $\vert$ \\
T & A & C & G & G & T & C & A \\
\end{tabular}
\end{center}

\textbf{2.} On compte les paires de bases : paires A--T : A, T, A, T soit \textbf{4 paires A--T} ; paires G--C : G, C, C, G soit \textbf{4 paires G--C}. Le fragment comporte donc 8 paires de bases au total.

\textbf{3.} L'ordre des bases le long du brin forme une sorte de « texte » écrit avec quatre lettres (A, T, G, C). Cette succession, différente d'un gène à l'autre et d'un individu à l'autre, constitue le message qui détermine les caractères héréditaires : c'est l'information génétique.

\textbf{Point de vigilance :} ne jamais apparier A avec G ni T avec C ; seuls les couples A--T et G--C sont corrects.
\end{corrige}

\begin{corrige}[Exercice 6 — Schéma d'une paire de chromosomes homologues]
\textbf{1. et 2.} Schéma attendu :

\begin{popfigure}
\begin{center}
\begin{tikzpicture}[scale=1.0]
% Chromosome paternel (bleu) en X
\draw[line width=4pt, popblue, rounded corners=2pt] (0,1.6) -- (0.55,0) -- (0, -1.6);
\draw[line width=4pt, popblue, rounded corners=2pt] (1.1,1.6) -- (0.55,0) -- (1.1,-1.6);
\fill[popdark] (0.55,0) circle (0.09);
% Chromosome maternel (rouge/orange) en X
\draw[line width=4pt, poporange, rounded corners=2pt] (2.6,1.6) -- (3.15,0) -- (2.6,-1.6);
\draw[line width=4pt, poporange, rounded corners=2pt] (3.7,1.6) -- (3.15,0) -- (3.7,-1.6);
\fill[popdark] (3.15,0) circle (0.09);
% Légendes fléchées
\draw[->, popdark] (0.55,0.15) -- (-1.3,0.9) node[left, font=\small] {1};
\draw[->, popdark] (0,1.2) -- (-1.3,1.9) node[left, font=\small] {2};
\draw[->, popdark] (1.1,0.9) -- (4.9,1.9) node[right, font=\small] {3};
\draw[->, popdark] (3.7,-0.9) -- (4.9,-1.4) node[right, font=\small] {4};
\node[font=\small, text=popblue] at (0.55,-2.1) {origine paternelle};
\node[font=\small, text=poporange] at (3.15,-2.1) {origine maternelle};
\end{tikzpicture}
\end{center}
\legende{Paire de chromosomes homologues bichromatidiens. 1 : centromère (point de réunion des chromatides) ; 2 : chromatide sœur ; 3 : emplacement du gène « couleur des yeux » (même emplacement sur les deux homologues) ; 4 : emplacement du gène « groupe sanguin ».}
\end{popfigure}

\textbf{3.} Un chromosome \cle{monochromatidien} est formé d'une \emph{seule} chromatide (aspect en bâtonnet simple), alors qu'un chromosome \cle{bichromatidien} est formé de \emph{deux} chromatides sœurs réunies par le centromère (aspect en X), après duplication de l'ADN.

\textbf{Point de vigilance :} les deux gènes doivent être placés \emph{aux mêmes emplacements} sur les deux homologues ; un schéma sans légende ni titre perd des points au BEPC.
\end{corrige}

\begin{corrige}[Exercice 7 — Protocole du TP d'extraction de l'ADN]
\textbf{1.} Le détergent détruit (dissout) les membranes des cellules et des noyaux, ce qui libère l'ADN dans la solution. Le sel favorise la séparation de l'ADN d'avec les protéines qui lui sont associées.

\textbf{2.} L'ADN est insoluble dans l'alcool, et d'autant plus que l'alcool est froid : l'alcool glacé provoque donc la précipitation de l'ADN. À l'interface entre le filtrat et l'alcool, on observe un dépôt blanchâtre et filamenteux : c'est l'ADN mis en évidence.

\textbf{3.} Le broyage brise les parois et les tissus végétaux et libère les cellules, ce qui augmente la surface de contact avec la solution et permet d'extraire une plus grande quantité d'ADN.

\textbf{Point de vigilance :} bien distinguer les rôles : broyage = libérer les cellules ; détergent = détruire les membranes ; alcool glacé = faire précipiter l'ADN.
\end{corrige}

\begin{corrige}[Exercice 8 — Maquette de chromosomes]
\textbf{1.} Paire monochromatidienne : un boudin rouge et un boudin bleu de même longueur, placés côte à côte. Paire bichromatidienne : deux boudins rouges réunis en X (par une boule de pâte figurant le centromère) et deux boudins bleus réunis de la même façon, les deux X étant placés côte à côte.

\textbf{2.} Chaque chromosome monochromatidien est représenté par un seul boudin : il faut donc \textbf{46 boudins} (23 rouges et 23 bleus, soit 23 paires d'homologues).

\textbf{3.} La différence de couleur représente l'\textbf{origine parentale} des chromosomes homologues : par exemple, le rouge pour les chromosomes d'origine maternelle et le bleu pour ceux d'origine paternelle.

\textbf{Point de vigilance :} dans une paire bichromatidienne, il faut 4 boudins au total (2 par chromosome) ; beaucoup d'élèves n'en mettent que 2.
\end{corrige}

\begin{corrige}[Exercice 9 — Type BEPC : l'expérience de Hammerling]
\textbf{1.} Problème : dans quelle partie de la cellule se trouve l'information qui détermine les caractères héréditaires (ici, la forme du chapeau) ? \hfill (1 pt)

\textbf{2.} Hypothèse : l'information déterminant la forme du chapeau se trouve dans le pied, car c'est la partie qui contient le noyau. \hfill (1 pt)

\textbf{3.} Dans les deux expériences, le chapeau régénéré a toujours la forme de l'espèce qui a fourni le \textbf{pied}, et non celle de l'espèce qui a fourni le stipe et le chapeau. C'est donc le pied, c'est-à-dire la partie contenant le noyau, qui détermine la forme du chapeau régénéré. \hfill (1 pt)

\textbf{4.} L'expérience 3 est une \textbf{expérience témoin}. Sans pied (donc sans noyau), il n'y a aucune régénération du chapeau : cela montre que le noyau est indispensable à la formation du chapeau et confirme qu'il détient l'information nécessaire. \hfill (1 pt)

\textbf{5.} Conclusion : l'information génétique qui détermine les caractères héréditaires est localisée dans le \cle{noyau} de la cellule. \hfill (1 pt)

\textbf{Barème indicatif :} 5 points au total, 1 point par question.

\textbf{Points de vigilance :}
\begin{itemize}
\item Le problème doit être formulé sous forme de question portant sur la \emph{localisation} de l'information, pas sur la forme du chapeau elle-même.
\item À la question 3, la comparaison doit citer explicitement les deux expériences et leurs résultats opposés : c'est le croisement des greffes qui constitue la preuve.
\item Le mot « témoin » est attendu à la question 4 ; dire simplement « expérience de contrôle » est accepté.
\end{itemize}
\end{corrige}

\begin{corrige}[Exercice 10 — Type BEPC : analyse d'un caryotype humain]
\textbf{1.} Le \cle{caryotype} est la photographie ordonnée de l'ensemble des chromosomes d'une cellule, classés par paires d'homologues selon leur taille et leur forme. \hfill (1 pt)

\textbf{2.} La cellule compte : 22 paires ordinaires (44 chromosomes) + 1 chromosome 21 surnuméraire + 2 chromosomes sexuels, soit $44 + 1 + 2 = \textbf{47 chromosomes}$. Normalement, une cellule humaine en compte 46 : il y a donc un chromosome en trop. \hfill (1 pt)

\textbf{3.} L'enfant possède deux chromosomes sexuels identiques XX : c'est une \textbf{fille} (un garçon aurait XY). \hfill (1 pt)

\textbf{4.} L'anomalie est la présence de trois exemplaires du chromosome 21 : c'est la \textbf{trisomie 21}, responsable du \textbf{syndrome de Down}. \hfill (1 pt)

\textbf{5.} Deux chromosomes homologues sont deux chromosomes de même forme et de même taille, portant les mêmes gènes aux mêmes emplacements. Chez un individu, l'un des deux homologues est d'origine \textbf{paternelle} (apporté par le spermatozoïde) et l'autre d'origine \textbf{maternelle} (apportée par l'ovule). \hfill (1 pt)

\textbf{Barème indicatif :} 5 points au total, 1 point par question.

\textbf{Points de vigilance :}
\begin{itemize}
\item Le décompte de la question 2 doit être détaillé ($44 + 1 + 2$) : un résultat « 47 » sans calcul visible peut être pénalisé.
\item Ne pas écrire « 23 paires + 1 » : la paire 21 compte ici trois chromosomes, ce n'est plus une paire.
\item Pour le sexe, la justification (XX = fille, XY = garçon) est obligatoire.
\end{itemize}
\end{corrige}

\begin{corrige}[Exercice 11 — Type BEPC : nature de l'information génétique]
\textbf{1.} La molécule d'ADN est formée de \textbf{deux brins} enroulés l'un autour de l'autre en \textbf{double hélice}. Chaque brin est un enchaînement de \textbf{nucléotides} ; chaque nucléotide comporte une base azotée parmi quatre : adénine (A), thymine (T), guanine (G) et cytosine (C). Les deux brins sont unis par l'appariement des bases : A avec T et G avec C. \hfill (1,5 pt)

\textbf{2.} Pour les deux organismes, on constate :
\begin{itemize}
\item A $\approx$ T (Homme : 30,4 $\approx$ 30,1 ; maïs : 26,8 $\approx$ 27,2) ;
\item G $\approx$ C (Homme : 19,6 $\approx$ 19,9 ; maïs : 22,8 $\approx$ 23,2).
\end{itemize}
Ces relations illustrent la \textbf{règle de complémentarité des bases} : dans l'ADN double brin, chaque A est associé à un T et chaque G à un C, donc A = T et G = C (les petits écarts du tableau sont dus aux imprécisions des mesures expérimentales). \hfill (1,5 pt)

\textbf{3.} Brin complémentaire : \hfill (1 pt)
\begin{center}
\begin{tabular}{cccccc}
T & A & C & G & G & T \\
$\vert$ & $\vert$ & $\vert$ & $\vert$ & $\vert$ & $\vert$ \\
A & T & G & C & C & A \\
\end{tabular}
\end{center}

\textbf{4.} La règle A = T et G = C concerne les \emph{relations entre les bases} à l'intérieur d'une même molécule d'ADN ; elle est donc valable pour tous les organismes. En revanche, la \emph{proportion globale} des bases (par exemple le rapport (A+T)/(G+C)) dépend de la succession des bases, qui est propre à chaque espèce : c'est pourquoi les pourcentages diffèrent entre l'Homme et le maïs. \hfill (1 pt)

\textbf{5.} L'ADN est la molécule qui porte l'information génétique, inscrite dans la succession de ses bases. Les molécules d'ADN, associées à des protéines, constituent les \textbf{chromosomes} situés dans le noyau. Chaque chromosome porte de nombreux \textbf{gènes}, et ce sont ces gènes qui déterminent les \textbf{caractères héréditaires} transmis des parents aux enfants. \hfill (1 pt)

\textbf{Barème indicatif :} 6 points au total (1,5 + 1,5 + 1 + 1 + 1).

\textbf{Points de vigilance :}
\begin{itemize}
\item À la question 2, il faut citer les valeurs du tableau pour appuyer chaque relation : une règle énoncée sans appui sur les données n'est pas une exploitation de document.
\item À la question 4, l'attendu est la distinction entre une règle \emph{universelle} (complémentarité) et des proportions \emph{spécifiques} à chaque espèce.
\item À la question 5, la chaîne ADN $\rightarrow$ chromosome $\rightarrow$ gène $\rightarrow$ caractère doit apparaître clairement et dans cet ordre logique.
\end{itemize}
\end{corrige}

\newpage

\coursec{Fiche bilan}

\begin{popfigure}
\begin{tikzpicture}[mindmap, grow cyclic, every node/.style={concept, font=\small},
  root concept/.append style={concept color=popblue, font=\bfseries\small, minimum size=3.2cm},
  level 1/.append style={level distance=3.4cm, sibling angle=51.4, font=\footnotesize},
  level 2/.append style={level distance=2.2cm, font=\scriptsize}]
\node[root concept]{Information génétique : localisation et nature}
  child[concept color=poporange]{ node {Problème : où est l'information ?}
    child { node {Hypothèses : noyau ? cytoplasme ?} } }
  child[concept color=popgreen]{ node {Transfert de noyau}
    child { node {Le caractère suit le noyau} }
    child { node {Info dans le noyau} } }
  child[concept color=popblue]{ node {Chromosomes}
    child { node {Visibles en division} }
    child { node {Paires homologues} }
    child { node {Mono / bichromatidiens} } }
  child[concept color=popgold]{ node {ADN}
    child { node {Molécule support} }
    child { node {Dans les chromosomes} } }
  child[concept color=popgreen!70!black]{ node {TP 1 : matériel génétique}
    child { node {Coloration + observation} } }
  child[concept color=poporange!80!black]{ node {TP 2 : maquettes}
    child { node {Paires homologues} } }
  child[concept color=popmuted]{ node {Applications}
    child { node {Hérédité, diagnostic} } };
\end{tikzpicture}
\legende{Carte mentale de la leçon : l'information génétique est localisée dans le noyau, portée par les chromosomes, constitués d'ADN.}
\end{popfigure}

\begin{bilan}
\textbf{Définitions clés}
\begin{itemize}
\item \cle{Information génétique} : ensemble des instructions qui déterminent les caractères héréditaires d'un individu et se transmettent à sa descendance.
\item \cle{Chromosome} : structure du noyau cellulaire, visible au microscope pendant la division cellulaire, support de l'information génétique.
\item \cle{Chromosomes homologues} : chromosomes d'une même paire, de même forme et de même taille, l'un d'origine paternelle, l'autre d'origine maternelle.
\item \cle{Chromosome monochromatidien} : chromosome formé d'une seule chromatide ; \cle{chromosome bichromatidien} : chromosome formé de deux chromatides identiques réunies par le centromère.
\item \cle{ADN} (acide désoxyribonucléique) : molécule constituant les chromosomes ; c'est elle qui porte l'information génétique.
\end{itemize}

\textbf{Résultats expérimentaux à connaître par cœur}
\begin{center}
\begin{tabularx}{\linewidth}{|l|X|X|}
\hline
\rowcolor{popblueL} \textbf{Expérience} & \textbf{Protocole} & \textbf{Conclusion} \\
\hline
Transfert de noyau & On transfère le noyau d'une cellule d'un individu A dans une cellule énucléée d'un individu B & L'individu obtenu ressemble à A (donneur du noyau) \\
\hline
Interprétation & Seul le noyau a été changé & L'information génétique se trouve dans le \cle{noyau} \\
\hline
Observation de cellules en division & Coloration + microscope & L'information est portée par les \cle{chromosomes}, faits d'\cle{ADN} \\
\hline
\end{tabularx}
\end{center}

\textbf{Méthodes express}
\begin{itemize}
\item Pour interpréter une expérience de transfert de noyau : identifier le donneur du noyau $\Rightarrow$ comparer les caractères de l'individu obtenu $\Rightarrow$ conclure que le caractère suit le noyau.
\item Pour reconnaître deux chromosomes homologues : même taille, même forme, même position du centromère.
\item Pour distinguer mono/bichromatidien : compter les chromatides (1 ou 2) réunies au centromère.
\item Démarche scientifique : observation $\rightarrow$ hypothèse $\rightarrow$ expérience $\rightarrow$ interprétation $\rightarrow$ conclusion.
\end{itemize}
\end{bilan}

\begin{aretenir}[Checklist avant le BEPC]
\begin{itemize}
\item[$\square$] Je sais définir l'information génétique et citer un caractère héréditaire.
\item[$\square$] Je sais décrire le protocole d'une expérience de transfert de noyau.
\item[$\square$] Je sais conclure de cette expérience que l'information génétique est dans le noyau.
\item[$\square$] Je sais définir un chromosome et dire quand il est observable.
\item[$\square$] Je sais reconnaître deux chromosomes homologues sur un schéma.
\item[$\square$] Je sais distinguer un chromosome monochromatidien d'un chromosome bichromatidien.
\item[$\square$] Je sais que les chromosomes sont constitués d'ADN.
\item[$\square$] Je sais que l'ADN est la molécule qui porte l'information génétique.
\item[$\square$] Je sais décrire le TP de mise en évidence du matériel génétique (coloration, observation).
\item[$\square$] Je sais construire ou légender une maquette de paire de chromosomes homologues.
\item[$\square$] Je sais rédiger une démarche complète : observation, hypothèse, expérience, interprétation, conclusion.
\item[$\square$] Je vérifie que ma conclusion répond exactement au problème posé, sans hors-programme.
\end{itemize}
\end{aretenir}

\findecours

\end{document}
